% ****************************************************************************************************
% Analyse_bipartit_symmetrisch.tex
%
% Teil A: Kritische Analyse des \Pi_2^P-Haertebeweises fuer bipartite AFs
%         (Proposition \ref{prop:bip-st-ind-hard} in Kapitel \ref{ch:graphclasses}).
% Teil B: Symmetrische irreflexive AFs: Ist \sigma-Unabhaengigkeit dort einfacher?
%
% Einbinden ueber % ****************************************************************************************************
% Analyse_bipartit_symmetrisch.tex
%
% Teil A: Kritische Analyse des \Pi_2^P-Haertebeweises fuer bipartite AFs
%         (Proposition \ref{prop:bip-st-ind-hard} in Kapitel \ref{ch:graphclasses}).
% Teil B: Symmetrische irreflexive AFs: Ist \sigma-Unabhaengigkeit dort einfacher?
%
% Einbinden ueber % ****************************************************************************************************
% Analyse_bipartit_symmetrisch.tex
%
% Teil A: Kritische Analyse des \Pi_2^P-Haertebeweises fuer bipartite AFs
%         (Proposition \ref{prop:bip-st-ind-hard} in Kapitel \ref{ch:graphclasses}).
% Teil B: Symmetrische irreflexive AFs: Ist \sigma-Unabhaengigkeit dort einfacher?
%
% Einbinden ueber % ****************************************************************************************************
% Analyse_bipartit_symmetrisch.tex
%
% Teil A: Kritische Analyse des \Pi_2^P-Haertebeweises fuer bipartite AFs
%         (Proposition \ref{prop:bip-st-ind-hard} in Kapitel \ref{ch:graphclasses}).
% Teil B: Symmetrische irreflexive AFs: Ist \sigma-Unabhaengigkeit dort einfacher?
%
% Einbinden ueber \include{chapters/examples/Analyse_bipartit_symmetrisch} in thesis.tex,
% nach \include{chapters/examples/Graphklassen}.
% ****************************************************************************************************

\chapter{Analyse des Härtebeweises für bipartite AFs}
\label{ch:bip-analyse}

\section{Gegenstand und Vorgehen}
\label{sec:bip-analyse-gegenstand}

Dieses Kapitel untersucht Proposition~\ref{prop:bip-st-ind-hard} aus
Kapitel~\ref{ch:graphclasses}, also die Behauptung, dass
\textit{stable}-Unabhängigkeit auch auf bipartiten AFs
\(\Pi_2^\text{P}\)-vollständig ist. Die Untersuchung erfolgt in drei Schritten.
Zuerst wird die Konstruktion so rekonstruiert, dass ihre Semantik durch ein
explizites Strukturlemma erfasst wird (Abschnitt~\ref{sec:bip-analyse-struktur}).
Anschließend wird die Korrektheit vollständig nachgerechnet
(Abschnitt~\ref{sec:bip-analyse-korrektheit}) und die
Graphklassenzugehörigkeit verifiziert
(Abschnitt~\ref{sec:bip-analyse-bipartit}). Erst danach werden die gefundenen
Mängel gesammelt (Abschnitt~\ref{sec:bip-analyse-kritik}) und behoben
(Abschnitt~\ref{sec:bip-analyse-reparatur}).

Das Ergebnis vorweg: \emph{Die Reduktion ist inhaltlich korrekt.} Die
Beweisführung weist jedoch eine Reihe formaler Mängel auf, von denen einer
(die unbelegte Ausgangsprobleminstanz) den Beweis in seiner jetzigen Form
angreifbar macht, während die übrigen die Nachvollziehbarkeit betreffen.

\section{Rekonstruktion der Konstruktion}
\label{sec:bip-analyse-struktur}

Wir wiederholen die Konstruktion in fixierter Notation. Sei
\(\Psi=\forall X\exists Y\,\varphi\) mit
\(\varphi=\bigwedge_{C\in\mathcal{C}}C\) in konjunktiver Normalform, wobei jede
Klausel \emph{monoton} ist, also entweder ausschließlich positive Literale
(\(C\in\mathcal{C}^{+}\)) oder ausschließlich negative Literale
(\(C\in\mathcal{C}^{-}\)) enthält. Sei \(V:=X\cup Y\cup\{p,n\}\) mit zwei
frischen Variablen \(p,n\). Das AF \(F_\Psi=(A_\Psi,R_\Psi)\) ist gegeben durch
\[
\begin{aligned}
    A_{\Psi}
    ={}&
    \{v,\overline{v}\mid v\in V\}
    \cup
    \{c_C\mid C\in\mathcal{C}\}
    \cup
    \{\varphi^{+},\varphi^{-},t^{+},t^{-}\}
\end{aligned}
\]
und der in Proposition~\ref{prop:bip-st-ind-hard} angegebenen Angriffsmenge
\(R_\Psi\). Zentral ist die Beobachtung, dass die Argumente \(p\) und
\(\overline{n}\) genau so wirken, als käme das Literal \(p\) in jeder positiven
und das Literal \(\neg n\) in jeder negativen Klausel vor. Wir machen das
explizit und setzen
\[
    \widehat{\varphi}
    :=
    \bigwedge_{C\in\mathcal{C}^{+}}(p\lor C)
    \;\land\;
    \bigwedge_{C\in\mathcal{C}^{-}}(\neg n\lor C).
\]

\begin{Lemma}
    \label{lem:bip-qbf-aequiv}
    Es gilt
    \(
        \Psi\text{ ist gültig}
        \iff
        \forall (X\cup\{p,n\})\,\exists Y\,\widehat{\varphi}
        \text{ ist gültig.}
    \)
\end{Lemma}

\begin{proof}
    \enquote{\(\Rightarrow\)}: Seien \(\alpha\) eine Belegung von \(X\) und
    \(\pi,\nu\) beliebige Wahrheitswerte für \(p\) und \(n\). Da \(\Psi\)
    gültig ist, existiert \(\beta\) mit \(\varphi(\alpha,\beta)=\top\). Dann ist
    jede Klausel \(C\in\mathcal{C}\) unter \((\alpha,\beta)\) erfüllt, also erst
    recht \(p\lor C\) beziehungsweise \(\neg n\lor C\). Somit gilt
    \(\widehat{\varphi}=\top\).

    \enquote{\(\Leftarrow\)}: Sei \(\alpha\) eine Belegung von \(X\). Wir wählen
    die Belegung \(p=\bot\) und \(n=\top\). Dann gilt
    \(\widehat{\varphi}\equiv\varphi\), und die Voraussetzung liefert ein
    \(\beta\) mit \(\varphi(\alpha,\beta)=\top\).
\end{proof}

Der Beweis in Kapitel~\ref{ch:graphclasses} argumentiert an mehreren Stellen mit
Aussagen der Form \enquote{die Argumente \(t^{+}\) und \(t^{-}\) erlauben es,
die Labels von \(\varphi^{+}\) und \(\varphi^{-}\) unabhängig zu setzen}, ohne
diese zu beweisen. Das folgende Lemma stellt die vollständige Struktur von
\(\Lambda_{\mathrm{st}}(F_\Psi)\) bereit und macht alle späteren Schritte
elementar.

\begin{Lemma}[Strukturlemma]
    \label{lem:bip-struktur}
    Sei \(\lambda\) ein Labelling von \(F_\Psi\). Dann ist \(\lambda\) genau
    dann \textit{stable}, wenn es eine totale Belegung \(\tau:V\to\{\top,\bot\}\)
    gibt, sodass die folgenden vier Bedingungen gelten.
    \begin{enumerate}
        \item Für alle \(v\in V\) ist
        \(\lambda(v)=\mathit{in}\iff\tau(v)=\top\)
        und
        \(\lambda(\overline{v})=\mathit{in}\iff\tau(v)=\bot\).
        \item Für alle \(C\in\mathcal{C}^{+}\) ist
        \(\lambda(c_C)=\mathit{out}\) genau dann, wenn \(\tau\models p\lor C\);
        andernfalls \(\lambda(c_C)=\mathit{in}\). Analog ist für
        \(C\in\mathcal{C}^{-}\) genau dann \(\lambda(c_C)=\mathit{out}\), wenn
        \(\tau\models\neg n\lor C\).
        \item \(\lambda(t^{+})=\mathit{in}\iff\lambda(\varphi^{+})=\mathit{out}\)
        und \(\lambda(t^{-})=\mathit{in}\iff\lambda(\varphi^{-})=\mathit{out}\).
        \item \(\lambda(\varphi^{+})=\mathit{in}\) impliziert
        \(\lambda(c_C)=\mathit{out}\) für alle \(C\in\mathcal{C}^{+}\);
        \(\lambda(\varphi^{-})=\mathit{in}\) impliziert
        \(\lambda(c_C)=\mathit{out}\) für alle \(C\in\mathcal{C}^{-}\).
    \end{enumerate}
    Umgekehrt liefert jede Wahl von \(\tau\) und jede mit (3) und (4)
    verträgliche Wahl von \(\lambda(\varphi^{+}),\lambda(\varphi^{-})\)
    ein \textit{stable}-Labelling. Insbesondere ist
    \(\lambda(\varphi^{+})=\mathit{out}\) stets realisierbar, und
    \(\lambda(\varphi^{+})=\mathit{in}\) ist genau dann realisierbar, wenn
    \(\tau\) alle Klauseln \(p\lor C\) mit \(C\in\mathcal{C}^{+}\) erfüllt.
\end{Lemma}

\begin{proof}
    Sei \(\lambda\in\Lambda_{\mathrm{st}}(F_\Psi)\); insbesondere ist
    \(\mathit{und}(\lambda)=\emptyset\).

    Zu (1): Die einzigen Angreifer von \(v\) sind \(\overline{v}\) und
    umgekehrt. Wären beide \(\mathit{in}\), so griffe ein
    \(\mathit{in}\)-Argument ein \(\mathit{in}\)-Argument an. Wären beide
    \(\mathit{out}\), so besäße \(v\) keinen \(\mathit{in}\)-Angreifer. Also ist
    genau eines der beiden \(\mathit{in}\), und wir setzen
    \(\tau(v):=\top\) genau dann, wenn \(\lambda(v)=\mathit{in}\).

    Zu (2): Die Angreifer von \(c_C\) mit \(C\in\mathcal{C}^{+}\) sind genau
    die Argumente \(v\) mit \(v\in C\) sowie \(p\). Nach (1) ist ein solcher
    Angreifer genau dann \(\mathit{in}\), wenn das zugehörige Literal unter
    \(\tau\) wahr ist. Da \(\lambda\) \textit{stable} ist, gilt
    \(\lambda(c_C)=\mathit{out}\) genau dann, wenn mindestens ein Angreifer
    \(\mathit{in}\) ist, also genau dann, wenn \(\tau\models p\lor C\). Der Fall
    \(C\in\mathcal{C}^{-}\) verläuft analog mit den Angreifern \(\overline{v}\)
    für \(\neg v\in C\) und \(\overline{n}\).

    Zu (3): Der einzige Angreifer von \(t^{+}\) ist \(\varphi^{+}\). Also ist
    \(t^{+}\) genau dann \(\mathit{out}\), wenn \(\varphi^{+}\) \(\mathit{in}\)
    ist, und wegen \(\mathit{und}(\lambda)=\emptyset\) genau dann
    \(\mathit{in}\), wenn \(\varphi^{+}\) \(\mathit{out}\) ist.

    Zu (4): Die Angreifer von \(\varphi^{+}\) sind die \(c_C\) mit
    \(C\in\mathcal{C}^{+}\) sowie \(t^{+}\). Ist \(\varphi^{+}\)
    \(\mathit{in}\), so sind alle Angreifer \(\mathit{out}\).

    Für die Umkehrung sei \(\tau\) beliebig und seien
    \(\lambda(\varphi^{\pm})\) gemäß (3) und (4) gewählt; die übrigen Labels
    seien durch (1)--(3) festgelegt. Man verifiziert unmittelbar für jedes
    Argument die \textit{stable}-Bedingung: Argumente \(v,\overline{v}\) sind
    genau dann \(\mathit{out}\), wenn ihr Partner \(\mathit{in}\) ist;
    Klauselargumente genau dann, wenn ein Literalargument \(\mathit{in}\) ist;
    \(t^{\pm}\) genau dann, wenn \(\varphi^{\pm}\) \(\mathit{in}\) ist; und
    \(\varphi^{+}\) ist entweder \(\mathit{in}\) (dann sind nach (4) alle
    \(c_C\) und nach (3) auch \(t^{+}\) \(\mathit{out}\)) oder \(\mathit{out}\)
    (dann ist nach (3) \(t^{+}\) \(\mathit{in}\) und greift \(\varphi^{+}\) an).
    Es tritt kein \(\mathit{und}\) auf.

    Der Zusatz folgt: \(\lambda(\varphi^{+})=\mathit{out}\) ist wegen
    \(\lambda(t^{+})=\mathit{in}\) immer zulässig, und
    \(\lambda(\varphi^{+})=\mathit{in}\) ist nach (4) und (2) genau dann
    zulässig, wenn \(\tau\) jede Klausel \(p\lor C\), \(C\in\mathcal{C}^{+}\),
    erfüllt.
\end{proof}

Lemma~\ref{lem:bip-struktur} zeigt insbesondere: Die Labels von
\(\varphi^{+}\) und \(\varphi^{-}\) sind \emph{unabhängig voneinander} frei
wählbar, sobald die jeweiligen Klauselargumente \(\mathit{out}\) sind, da
\(\varphi^{+},t^{+}\) und \(\varphi^{-},t^{-}\) disjunkte Angreifermengen
besitzen. Genau diese Aussage benutzt der Originalbeweis, ohne sie zu begründen.

\section{Korrektheit der Reduktion}
\label{sec:bip-analyse-korrektheit}

\begin{Proposition}
    \label{prop:bip-korrekt}
    Es gilt
    \[
        \Psi\text{ ist gültig}
        \quad\Longleftrightarrow\quad
        \{\varphi^{+},\varphi^{-}\}\bot_{\mathrm{st}}(X\cup\{p,n\})\mid\emptyset
        \text{ in }F_\Psi.
    \]
\end{Proposition}

\begin{proof}
    Wir schreiben \(\mathcal{X}:=\{\varphi^{+},\varphi^{-}\}\) und
    \(\mathcal{Y}:=X\cup\{p,n\}\); beide Mengen sind disjunkt, und
    \(Z=\emptyset\), sodass die Bedingung \(\lambda_1|_Z=\lambda_2|_Z\) aus
    Definition~\ref{def:BU:AFs} für jedes Paar erfüllt ist. Gesucht ist also
    stets ein \(\lambda_3\in\Lambda_{\mathrm{st}}(F_\Psi)\) mit
    \(\lambda_3|_{\mathcal{X}}=\lambda_1|_{\mathcal{X}}\) und
    \(\lambda_3|_{\mathcal{Y}}=\lambda_2|_{\mathcal{Y}}\).

    \enquote{\(\Rightarrow\)}: Sei \(\Psi\) gültig und seien
    \(\lambda_1,\lambda_2\in\Lambda_{\mathrm{st}}(F_\Psi)\). Nach
    Lemma~\ref{lem:bip-struktur} induziert \(\lambda_2\) eine totale Belegung
    \(\tau_2\) von \(V\); insbesondere legt \(\lambda_2|_{\mathcal{Y}}\) die
    Werte \(\tau_2|_{X\cup\{p,n\}}\) fest. Nach Lemma~\ref{lem:bip-qbf-aequiv}
    existiert eine Belegung \(\beta\) von \(Y\), sodass
    \(\tau_3:=\tau_2|_{X\cup\{p,n\}}\cup\beta\) die Formel
    \(\widehat{\varphi}\) erfüllt. Wir wählen \(\lambda_3\) als das durch
    \(\tau_3\) und
    \(\lambda_3(\varphi^{+}):=\lambda_1(\varphi^{+})\),
    \(\lambda_3(\varphi^{-}):=\lambda_1(\varphi^{-})\)
    bestimmte Labelling. Da \(\tau_3\models\widehat{\varphi}\), sind nach
    Lemma~\ref{lem:bip-struktur}(2) sämtliche Klauselargumente
    \(\mathit{out}\); damit sind beide Wahlen für \(\varphi^{+}\) und
    \(\varphi^{-}\) mit (4) verträglich, und \(\lambda_3\) ist
    \textit{stable}. Nach Konstruktion gilt
    \(\lambda_3|_{\mathcal{X}}=\lambda_1|_{\mathcal{X}}\) und
    \(\lambda_3|_{\mathcal{Y}}=\lambda_2|_{\mathcal{Y}}\).

    \enquote{\(\Leftarrow\)}: Sei \(\Psi\) nicht gültig. Dann existiert eine
    Belegung \(\alpha\) von \(X\), sodass \(\varphi(\alpha,\beta)=\bot\) für
    jede Belegung \(\beta\) von \(Y\) gilt. Wir wählen:
    \begin{itemize}
        \item \(\lambda_1\) als ein beliebiges \textit{stable}-Labelling mit
        \(\lambda_1(\varphi^{+})=\lambda_1(\varphi^{-})=\mathit{in}\). Ein
        solches existiert: Man setze \(\tau_1(p)=\top\) und \(\tau_1(n)=\bot\)
        sowie \(\tau_1\) auf \(X\cup Y\) beliebig. Dann sind alle Klauseln
        \(p\lor C\) und \(\neg n\lor C\) erfüllt, also nach
        Lemma~\ref{lem:bip-struktur} alle Klauselargumente \(\mathit{out}\)
        und beide \(\varphi\)-Argumente \(\mathit{in}\) realisierbar.
        \item \(\lambda_2\) als ein \textit{stable}-Labelling zur Belegung
        \(\tau_2:=\alpha\cup\{p\mapsto\bot,\;n\mapsto\top\}\cup\beta_0\) mit
        beliebigem \(\beta_0\).
    \end{itemize}
    Angenommen, es gäbe ein \(\lambda_3\in\Lambda_{\mathrm{st}}(F_\Psi)\) mit
    \(\lambda_3|_{\mathcal{X}}=\lambda_1|_{\mathcal{X}}\) und
    \(\lambda_3|_{\mathcal{Y}}=\lambda_2|_{\mathcal{Y}}\). Sei \(\tau_3\) die
    zugehörige Belegung. Aus
    \(\lambda_3|_{\mathcal{Y}}=\lambda_2|_{\mathcal{Y}}\) folgt
    \(\tau_3|_X=\alpha\), \(\tau_3(p)=\bot\) und \(\tau_3(n)=\top\). Aus
    \(\lambda_3(\varphi^{+})=\lambda_3(\varphi^{-})=\mathit{in}\) folgt mit
    Lemma~\ref{lem:bip-struktur}(4) und (2), dass \(\tau_3\) sämtliche Klauseln
    \(p\lor C\) und \(\neg n\lor C\) erfüllt, also
    \(\tau_3\models\widehat{\varphi}\). Wegen \(\tau_3(p)=\bot\) und
    \(\tau_3(n)=\top\) gilt aber \(\widehat{\varphi}\equiv\varphi\) unter
    \(\tau_3\), also \(\varphi(\alpha,\tau_3|_Y)=\top\) im Widerspruch zur Wahl
    von \(\alpha\). Also existiert kein solches \(\lambda_3\), und es gilt
    \(\mathcal{X}\not\bot_{\mathrm{st}}\mathcal{Y}\mid\emptyset\).
\end{proof}

Die Konstruktion ist offensichtlich in Polynomialzeit berechenbar: Sie erzeugt
\(2|V|+|\mathcal{C}|+4\) Argumente und
\(\mathcal{O}(|V|+\sum_{C\in\mathcal{C}}|C|)\) Angriffe.

\section{Graphklassenzugehörigkeit}
\label{sec:bip-analyse-bipartit}

\begin{Proposition}
    \label{prop:bip-korrekt-bipartit}
    \(F_\Psi\) ist bipartit.
\end{Proposition}

\begin{proof}
    Wir setzen
    \[
    \begin{aligned}
        L:={}& \{v\mid v\in V\}\cup\{c_C\mid C\in\mathcal{C}^{-}\}
              \cup\{\varphi^{+},t^{-}\},\\
        R:={}& \{\overline{v}\mid v\in V\}\cup\{c_C\mid C\in\mathcal{C}^{+}\}
              \cup\{\varphi^{-},t^{+}\}.
    \end{aligned}
    \]
    Offensichtlich gilt \(L\cap R=\emptyset\) und \(L\cup R=A_\Psi\). Für jede
    Art von Angriff prüft man die Seitenwechsel-Eigenschaft:
    \((v,\overline{v})\) und \((\overline{v},v)\) verlaufen zwischen \(L\) und
    \(R\); \((v,c_C)\) mit \(C\in\mathcal{C}^{+}\) sowie \((p,c_C)\) verlaufen
    von \(L\) nach \(R\); \((\overline{v},c_C)\) mit \(C\in\mathcal{C}^{-}\)
    sowie \((\overline{n},c_C)\) verlaufen von \(R\) nach \(L\);
    \((c_C,\varphi^{+})\) mit \(C\in\mathcal{C}^{+}\) verläuft von \(R\) nach
    \(L\); \((c_C,\varphi^{-})\) mit \(C\in\mathcal{C}^{-}\) von \(L\) nach
    \(R\); \((\varphi^{+},t^{+})\) von \(L\) nach \(R\) und
    \((\varphi^{-},t^{-})\) von \(R\) nach \(L\), jeweils mit den
    Gegenangriffen. Also gilt
    \(R_\Psi\subseteq(L\times R)\cup(R\times L)\).
\end{proof}

Hier zeigt sich der eigentliche Grund für die Monotonieforderung: Nur weil
positive Klauseln ausschließlich von unquerstrichenen und negative Klauseln
ausschließlich von querstrichenen Literalargumenten angegriffen werden, lassen
sich die Klauselargumente überhaupt widerspruchsfrei auf die beiden Seiten
verteilen. Bei gemischten Klauseln entstünde ein Klauselargument mit Angreifern
aus \(L\) \emph{und} \(R\), und die Bipartition bräche zusammen.

\section{Kritikpunkte}
\label{sec:bip-analyse-kritik}

Die Reduktion ist nach den Abschnitten~\ref{sec:bip-analyse-korrektheit}
und~\ref{sec:bip-analyse-bipartit} korrekt. Der Beweis in
Kapitel~\ref{ch:graphclasses} weist jedoch die folgenden Mängel auf. Wir
ordnen sie nach abnehmender Schwere.

\begin{enumerate}
    \item \textbf{Unbelegtes Ausgangsproblem.} Der Beweis reduziert von
    \textsc{Balanced Monotone \(\forall\exists\)-3-SAT-\((1,1,2,2)\)} und
    behauptet ohne Quellenangabe und ohne Definition, dieses Problem sei
    \(\Pi_2^\text{P}\)-vollständig. Weder wird gesagt, worauf sich
    \enquote{balanced} bezieht, noch was das Tupel \((1,1,2,2)\) zählt. Das ist
    der einzige Mangel, der die \emph{Gültigkeit} der Aussage betrifft: Eine
    Reduktion von einem Problem unbekannter Komplexität beweist nichts.
    Vorkommensbeschränkungen sind zudem heikel. Nach dem Satz von Tovey ist
    jede \(k\)-CNF-Formel, in der jede Klausel \(k\) verschiedene Literale
    besitzt und jede Variable in höchstens \(k\) Klauseln vorkommt, erfüllbar;
    zu enge Vorkommensschranken können ein Problem also trivialisieren.
    \item \textbf{Die Zusatzrestriktionen werden nicht benötigt.} Die
    Konstruktion verwendet ausschließlich die Monotonie der Klauseln (siehe
    Abschnitt~\ref{sec:bip-analyse-bipartit}); weder die Ausgeglichenheit von
    \(|\mathcal{C}^{+}|\) und \(|\mathcal{C}^{-}|\) noch irgendeine
    Vorkommensschranke noch die Klauselbreite \(3\) geht in Korrektheit oder
    Bipartitheit ein. Der Beweis macht sich also ohne Not von einer stärkeren
    und schwerer zu belegenden Voraussetzung abhängig.
    Abschnitt~\ref{sec:bip-analyse-reparatur} zeigt, dass die Monotonie
    elementar und selbsttragend erreicht werden kann.
    \item \textbf{Vertauschte Rollen von \(\lambda_1\) und \(\lambda_2\).} In
    der Richtung \enquote{\(\Psi\) gültig} übernimmt \(\lambda_1\) die Labels
    von \(\varphi^{+},\varphi^{-}\) und \(\lambda_2\) die Belegung von
    \(X\cup\{p,n\}\), passend zur Reihenfolge der Mengen in
    \(\{\varphi^{+},\varphi^{-}\}\bot_{\mathrm{st}}X\cup\{p,n\}\mid\emptyset\)
    und zu Definition~\ref{def:BU:AFs}. In der Gegenrichtung sind die Rollen
    stillschweigend vertauscht: Dort trägt \(\lambda_1\) die Belegung und
    \(\lambda_2\) die \(\varphi\)-Labels. Inhaltlich ist das folgenlos, weil
    die bedingte Unabhängigkeit symmetrisch ist, formal ist es jedoch ein
    Bruch mit der eigenen Definition. Entweder man benennt konsistent um (wie
    in Proposition~\ref{prop:bip-korrekt}) oder man beruft sich explizit auf
    das Symmetrieaxiom.
    \item \textbf{Undefinierte Mengen \(A\) und \(B\).} Der Satz
    \enquote{[\ldots] das die Einschränkung von \(\lambda_1\) auf \(A\) und die
    Einschränkung von \(\lambda_2\) auf \(B\) kombiniert} verwendet zwei nie
    eingeführte Symbole. Erschwerend kommt hinzu, dass \(A\) in der gesamten
    Arbeit die Argumentmenge eines AFs bezeichnet; gemeint sind offenbar
    \(\{\varphi^{+},\varphi^{-}\}\) und \(X\cup\{p,n\}\).
    \item \textbf{Namenskollision \(L,R\).} Die Bipartitionsklassen heißen
    \(L\) und \(R\), während \(R\) beziehungsweise \(R_\Psi\) die
    Angriffsrelation bezeichnet. Die Schlussformel
    \(R_\Psi\subseteq(L\times R)\cup(R\times L)\) mischt beide Bedeutungen in
    einer Zeile. Vorzuziehen wären etwa \(A_1,A_2\) oder \(L_1,L_2\).
    \item \textbf{Fehlendes Strukturlemma.} Der Beweis benutzt mehrfach
    Eigenschaften von \(\Lambda_{\mathrm{st}}(F_\Psi)\), die er nicht beweist:
    dass in jedem \textit{stable}-Labelling genau eines von \(v,\overline{v}\)
    das Label \(\mathit{in}\) trägt; dass \(c_C\) genau dann \(\mathit{out}\)
    ist, wenn ein Literalargument \(\mathit{in}\) ist; und dass \(t^{+},t^{-}\)
    die freie und \emph{voneinander unabhängige} Wahl der \(\varphi\)-Labels
    erlauben. Lemma~\ref{lem:bip-struktur} schließt diese Lücke.
    \item \textbf{Unvollständige Fallunterscheidung über \(p\) und \(n\).} Der
    Abschnitt \enquote{Konstruktion} diskutiert nur die Kombinationen
    \((p,n)=(\mathit{out},\mathit{in})\) und
    \((p,n)=(\mathit{in},\mathit{out})\). Die ebenfalls in
    \textit{stable}-Labellings auftretenden Kombinationen
    \((\mathit{in},\mathit{in})\) und \((\mathit{out},\mathit{out})\) bleiben
    unerwähnt. Der Korrektheitsbeweis benötigt diese Fallunterscheidung
    tatsächlich nicht --- er läuft über die Formeläquivalenz aus
    Lemma~\ref{lem:bip-qbf-aequiv} ---, der Text suggeriert jedoch eine
    erschöpfende Analyse, die er nicht liefert.
    \item \textbf{Zugehörigkeit fehlt.} Die Proposition behauptet
    \(\Pi_2^\text{P}\)-\emph{Vollständigkeit}, der Beweis zeigt nur die Härte.
    Ein Verweis auf Theorem~\ref{th:ind-st-is-pi2p-c} (die obere Schranke gilt
    für alle AFs, also insbesondere für bipartite) fehlt, anders als beim
    ungeradezykelfreien Fall.
    \item \textbf{Formulierungs- und Notationsmängel.} Die Instanz wird als
    Behauptung formuliert (\enquote{wählen wir
    \(\{\varphi^{+},\varphi^{-}\}\perp_{st}\ldots\)}) statt als Anfrage; die
    Polynomialität der Konstruktion wird nicht erwähnt; \(\perp\) wird sowohl
    für die Unabhängigkeitsrelation als auch (in \enquote{\(p=\perp\)}) für den
    Wahrheitswert \emph{falsch} verwendet, während an anderer Stelle
    \(\bot_{\mathrm{st}}\) für die Unabhängigkeit steht; die Argumentmenge
    heißt hier \(\mathcal{A}_\Psi\), im ungeradezykelfreien Fall dagegen
    \(A_\Psi\); und der Satz \enquote{dann werden die Klauselargumente
    ausschließlich durch die ursprünglichen Literale kontrolliert werden}
    enthält ein doppeltes Verb.
\end{enumerate}

\begin{Corollary}
    \label{cor:bip-subsumiert-ocf}
    Proposition~\ref{prop:bip-st-ind-hard} impliziert
    Proposition~\ref{prop:ocf-af-ci-st-pi^2}.
\end{Corollary}

\begin{proof}
    Ist \(F\) bipartit, so wechselt jeder gerichtete Zyklus in jedem Schritt die
    Bipartitionsklasse und hat daher gerade Länge. Bipartite AFs sind also
    ungeradezykelfrei, und die Härte auf der Teilklasse überträgt sich auf die
    Oberklasse.
\end{proof}

Das ist ein struktureller Hinweis, kein Fehler: Der eigenständige Beweis für
ungeradezykelfreie AFs für \(\sigma=\mathrm{st}\) und \(\sigma=\mathrm{pr}\) ist
entbehrlich, sobald das bipartite Resultat gesichert ist. Umgekehrt bleibt der
gesonderte Beweis für \(\sigma=\mathrm{co}\)
(Proposition~\ref{prop:ocf-af-ci-co-pi^2-c}) nötig, denn das dort verwendete AF
ist wegen der Argumente \(d_i\), die sowohl von \(x_i\) als auch von
\(\overline{x_i}\) angegriffen werden, nicht bipartit.

Schließlich sei angemerkt, dass die Herleitung von
Proposition~\ref{prop:bip-prf-ind-hard} aus dem \textit{stable}-Fall die
Kohärenz ungeradezykelfreier AFs voraussetzt, also
\(\Lambda_{\mathrm{st}}(F)=\Lambda_{\mathrm{pr}}(F)\). Diese Tatsache stammt von
\textcite{dung:1995:af} und sollte dort zitiert werden.

\section{Reparatur: Monotonie ohne fremde Voraussetzung}
\label{sec:bip-analyse-reparatur}

Der schwerwiegendste Kritikpunkt lässt sich vollständig beheben. Statt sich auf
eine exotische Restriktion zu berufen, genügt eine elementare Normalform, die
sich in wenigen Zeilen beweisen lässt.

\begin{Lemma}[Monotone Normalform]
    \label{lem:monotone-normalform}
    Zu jeder Formel \(\Psi=\forall X\exists Y\,\varphi\) mit \(\varphi\) in KNF
    lässt sich in Polynomialzeit eine Formel
    \(\Psi'=\forall X\exists Y'\,\varphi'\) berechnen, sodass jede Klausel von
    \(\varphi'\) entweder nur positive oder nur negative Literale enthält und
    \(\Psi\) genau dann gültig ist, wenn \(\Psi'\) gültig ist.
\end{Lemma}

\begin{proof}
    Für jede Variable \(v\in X\cup Y\) sei \(v'\) eine frische Variable und
    \(Y':=Y\cup\{v'\mid v\in X\cup Y\}\). Die Formel \(\varphi'\) entstehe aus
    \(\varphi\), indem jedes negative Literal \(\neg v\) durch \(v'\) ersetzt
    wird; zusätzlich nehmen wir für jedes \(v\in X\cup Y\) die beiden Klauseln
    \(
        (v\lor v')
    \) und \(
        (\neg v\lor\neg v')
    \)
    auf. Alle aus \(\varphi\) entstandenen Klauseln sind rein positiv, die
    ersten Zusatzklauseln ebenfalls, die zweiten rein negativ. Also ist
    \(\varphi'\) monoton. Die beiden Zusatzklauseln sind zusammen äquivalent zu
    \(v'\leftrightarrow\neg v\).

    Sei \(\Psi\) gültig und \(\alpha\) eine Belegung von \(X\). Wähle \(\beta\)
    mit \(\varphi(\alpha,\beta)=\top\) und setze \(v'\) auf den zu \(v\)
    komplementären Wert. Die Zusatzklauseln sind dann erfüllt. Ist eine
    ursprüngliche Klausel \(C\) durch ein positives Literal \(v\) erfüllt, so
    auch ihr Bild; ist sie durch \(\neg v\) erfüllt, so ist \(v\) falsch, also
    \(v'\) wahr, und das Bild ist ebenfalls erfüllt. Somit ist \(\Psi'\) gültig.

    Sei umgekehrt \(\Psi'\) gültig und \(\alpha\) eine Belegung von \(X\). Es
    existiert eine Belegung von \(Y'\), die \(\varphi'\) erfüllt; wegen der
    Zusatzklauseln gilt dort \(v'=\neg v\) für alle \(v\in X\cup Y\). Sei
    \(\beta\) die Einschränkung auf \(Y\). Jede ursprüngliche Klausel \(C\) hat
    ein erfülltes Bildliteral; ist dieses ein \(v\), so erfüllt \(v\) auch
    \(C\); ist es ein \(v'\), so ist \(v\) falsch und \(\neg v\in C\) erfüllt.
    Also gilt \(\varphi(\alpha,\beta)=\top\) und \(\Psi\) ist gültig.

    Die Konstruktion verdoppelt die Variablenzahl und fügt \(2|X\cup Y|\)
    Klauseln hinzu, ist also polynomiell.
\end{proof}

Damit ist \textsc{Monotone \(\forall\exists\)-SAT} \(\Pi_2^\text{P}\)-hart,
sobald \(\forall\exists\)-SAT es ist, und
Proposition~\ref{prop:bip-st-ind-hard} steht auf einer belegten Grundlage. Wird
zusätzlich Klauselbreite \(3\) gewünscht, füllt man die zweielementigen
Zusatzklauseln durch Wiederholung eines Literals auf; die Monotonie bleibt
erhalten. Alle übrigen Kritikpunkte sind redaktioneller Natur und durch die
Fassung in den Abschnitten~\ref{sec:bip-analyse-struktur}
bis~\ref{sec:bip-analyse-bipartit} bereits ausgeräumt.

\section{Fazit}
\label{sec:bip-analyse-fazit}

Die Aussage von Proposition~\ref{prop:bip-st-ind-hard} ist richtig, und der
Kern der Konstruktion --- die Aufspaltung des Formelarguments in \(\varphi^{+}\)
und \(\varphi^{-}\) sowie die Steuerung durch die frischen Variablen \(p\) und
\(n\) --- ist elegant und trägt die Bipartitheit. Der Beweis sollte jedoch um
Lemma~\ref{lem:bip-qbf-aequiv}, Lemma~\ref{lem:bip-struktur} und
Lemma~\ref{lem:monotone-normalform} ergänzt, in den Rollen von
\(\lambda_1,\lambda_2\) vereinheitlicht und um den Zugehörigkeitsverweis
ergänzt werden. Danach ist er lückenlos.

\chapter{Symmetrische irreflexive AFs}
\label{ch:sym-analyse}

\section{Fragestellung}
\label{sec:sym-frage}

Abschnitt~\ref{sec:sym-afs} hat symmetrische AFs bereits als Extremfall hoher
Angriffsdichte betrachtet und für vollständig gerichtete AFs gezeigt, dass
\textit{stable}-Unabhängigkeit dort stets verletzt ist
(Proposition~\ref{prop:cdg-afs-not-st-ci}). Offen blieb die
komplexitätstheoretische Einordnung der gesamten Klasse. Dieses Kapitel
schließt diese Lücke.

Der Ausgangsverdacht ist, dass die Klasse einfach sein müsste. Symmetrische
irreflexive AFs sind kohärent \cite{Coste:2005:symmetric-af-irflx-coherent},
besitzen stets \textit{stable}-Labellings und ihre klassischen
Entscheidungsprobleme --- Existenz, Verifikation, kredulöse und skeptische
Akzeptanz --- sind sämtlich in Polynomialzeit lösbar
(Proposition~\ref{prop:sym-klassisch-p}). Für alle bisher betrachteten
Graphklassen (Abschnitt~\ref{sec:acyc-afs} bis
Abschnitt~\ref{sec:odd-cycle-free-afs}) war eine solche strukturelle
Vereinfachung stets mit einer Vereinfachung der Unabhängigkeit einhergegangen
oder zumindest verträglich.

Das Hauptresultat dieses Kapitels widerlegt diesen Verdacht für die
\textit{stable}-Semantik: \(\textit{Ind}_{\mathrm{st}}\) bleibt auf
symmetrischen irreflexiven AFs \(\Pi_2^\text{P}\)-vollständig
(Theorem~\ref{th:sym-st-pi2p}). Für die \textit{preferred}-Semantik dagegen
tritt eine echte Vereinfachung ein: Das Problem fällt von
\(\Pi_3^\text{P}\)-Vollständigkeit auf \(\Pi_2^\text{P}\)-Vollständigkeit
(Korollar~\ref{cor:sym-pr-pi2p}). Ergänzend wird ein Fragment identifiziert, in
dem die Unabhängigkeit tatsächlich in Polynomialzeit entscheidbar wird
(Proposition~\ref{prop:sym-singleton-p}).

Wir betrachten für \(\sigma\in\{co,st,pr\}\) das Entscheidungsproblem
\[
\begin{array}{ll}
\textit{Ind}^{\mathrm{sym}}_{\sigma}\\
\text{Input:}  & \text{symmetrisches, irreflexives AF } F=(A,R),\
                 \text{disjunkte } X,Y,Z\subseteq A \\
\text{Output:} & \text{TRUE} \Longleftrightarrow X \bot_{\sigma} Y \mid Z.
\end{array}
\]

\section{Struktur der Labellingmengen}
\label{sec:sym-struktur}

\begin{Definition}
    \label{def:sym-af}
    Ein AF \(F=(A,R)\) heißt \emph{symmetrisch}, falls
    \((a,b)\in R\Rightarrow(b,a)\in R\) für alle \(a,b\in A\) gilt, und
    \emph{irreflexiv}, falls \((a,a)\notin R\) für alle \(a\in A\) gilt. Für
    \(a\in A\) schreiben wir \(N(a):=\{b\in A\mid (a,b)\in R\}\) für die
    \emph{Nachbarschaft} von \(a\) und \(N[a]:=N(a)\cup\{a\}\). Für
    \(S\subseteq A\) sei \(N(S):=\bigcup_{a\in S}N(a)\).
\end{Definition}

In einem symmetrischen irreflexiven AF ist \(R\) die Kantenrelation eines
schlichten ungerichteten Graphen. Eine Menge \(S\subseteq A\) ist genau dann
konfliktfrei, wenn sie in diesem Graphen unabhängig ist, und wir nennen \(S\)
\emph{naiv}, falls \(S\) \(\subseteq\)-maximal konfliktfrei ist. Für
\(S\subseteq A\) bezeichne \(\lambda_S\) das Labelling mit
\(\mathit{in}(\lambda_S)=S\) und \(\mathit{out}(\lambda_S)=A\setminus S\).

\begin{Lemma}
    \label{lem:sym-cf-adm}
    Sei \(F=(A,R)\) symmetrisch und irreflexiv. Dann ist jede konfliktfreie
    Menge \(S\subseteq A\) admissibel.
\end{Lemma}

\begin{proof}
    Sei \(a\in S\) und sei \(b\) ein Angreifer von \(a\), also \((b,a)\in R\).
    Wegen der Konfliktfreiheit von \(S\) gilt \(b\notin S\). Aus der Symmetrie
    folgt \((a,b)\in R\), also greift \(a\in S\) den Angreifer \(b\) an. Damit
    verteidigt \(S\) jedes seiner Elemente.
\end{proof}

\begin{Proposition}
    \label{prop:sym-naiv-st-pr}
    Sei \(F=(A,R)\) symmetrisch und irreflexiv. Dann gilt
    \[
        \Lambda_{\mathrm{st}}(F)
        =\Lambda_{\mathrm{pr}}(F)
        =\{\lambda_S\mid S\subseteq A\text{ naiv}\}.
    \]
    Insbesondere ist \(\Lambda_{\mathrm{st}}(F)\neq\emptyset\).
\end{Proposition}

\begin{proof}
    Sei \(S\) naiv und \(a\in A\setminus S\). Dann ist \(S\cup\{a\}\) nicht
    konfliktfrei. Da \(S\) konfliktfrei und \(F\) irreflexiv ist, muss der
    Konflikt zwischen \(a\) und einem \(b\in S\) bestehen; wegen der Symmetrie
    gilt \((b,a)\in R\). Also greift \(S\) jedes Argument außerhalb von \(S\)
    an, das heißt \(\lambda_S\) ist ein \textit{stable}-Labelling.

    Umgekehrt ist jede \textit{stable}-Extension konfliktfrei und maximal
    konfliktfrei, denn ein hinzunehmbares Argument wäre unangegriffen. Also
    stimmen die \textit{stable}-Extensionen genau mit den naiven Mengen
    überein.

    Für die \textit{preferred}-Semantik gilt allgemein
    \(\Lambda_{\mathrm{st}}(F)\subseteq\Lambda_{\mathrm{pr}}(F)\). Sei
    umgekehrt \(E\) eine \textit{preferred}-Extension, also
    \(\subseteq\)-maximal admissibel. Wäre \(E\) nicht maximal konfliktfrei, so
    gäbe es eine konfliktfreie Menge \(E\subsetneq T\); nach
    Lemma~\ref{lem:sym-cf-adm} wäre \(T\) admissibel, im Widerspruch zur
    Maximalität von \(E\). Also ist \(E\) naiv und damit \textit{stable}. Da
    eine \textit{stable}-Extension \(E\) die Labellingdarstellung
    \(\mathit{in}=E\), \(\mathit{out}=A\setminus E\), \(\mathit{und}=\emptyset\)
    besitzt und dies zugleich das zu \(E\) gehörige
    \textit{preferred}-Labelling ist, stimmen auch die Labellingmengen überein.

    Da \(A\) endlich ist, besitzt jede konfliktfreie Menge eine maximale
    Obermenge; insbesondere existiert mindestens eine naive Menge.
\end{proof}

\begin{Corollary}
    \label{cor:sym-extend}
    Sei \(F=(A,R)\) symmetrisch und irreflexiv und sei \(T\subseteq A\)
    konfliktfrei. Dann existiert ein \(\lambda\in\Lambda_{\mathrm{st}}(F)\) mit
    \(T\subseteq\mathit{in}(\lambda)\). Ist zusätzlich \(a\in A\setminus T\) und
    besitzt \(a\) einen Nachbarn in \(T\), so gilt
    \(\lambda(a)=\mathit{out}\) für jedes solche \(\lambda\).
\end{Corollary}

\begin{proof}
    Man erweitere \(T\) zu einer naiven Menge \(S\) und wende
    Proposition~\ref{prop:sym-naiv-st-pr} an. Besitzt \(a\) einen Nachbarn in
    \(T\subseteq S\), so verhindert die Konfliktfreiheit \(a\in S\).
\end{proof}

Korollar~\ref{cor:sym-extend} ist das Arbeitspferd aller folgenden Beweise: Es
erlaubt, ein Labelling durch Angabe eines konfliktfreien \enquote{Kerns} zu
spezifizieren, sofern jedes Argument, das ausgeschlossen werden soll, einen
Nachbarn im Kern besitzt.

\begin{Proposition}
    \label{prop:sym-klassisch-p}
    Sei \(F=(A,R)\) symmetrisch und irreflexiv und sei \(\sigma\in\{st,pr\}\).
    Dann gilt:
    \begin{enumerate}
        \item \(\Lambda_\sigma(F)\neq\emptyset\).
        \item Die Verifikation, ob ein gegebenes Labelling in
        \(\Lambda_\sigma(F)\) liegt, ist in Polynomialzeit möglich.
        \item Jedes \(a\in A\) ist kredulös akzeptiert.
        \item Ein Argument \(a\in A\) ist genau dann skeptisch akzeptiert, wenn
        \(N(a)=\emptyset\).
    \end{enumerate}
\end{Proposition}

\begin{proof}
    (1) und (2) folgen aus Proposition~\ref{prop:sym-naiv-st-pr}: Zu prüfen ist
    lediglich, ob \(\mathit{in}(\lambda)\) konfliktfrei ist und ob jedes
    Argument außerhalb einen Nachbarn darin besitzt; beides ist in
    \(\mathcal{O}(|A|+|R|)\) möglich.

    (3) Wegen der Irreflexivität ist \(\{a\}\) konfliktfrei; nach
    Korollar~\ref{cor:sym-extend} existiert ein \(\lambda\in\Lambda_\sigma(F)\)
    mit \(\lambda(a)=\mathit{in}\).

    (4) Ist \(N(a)=\emptyset\), so ist \(a\) in jeder maximal konfliktfreien
    Menge enthalten, denn \(a\) steht mit keinem Argument in Konflikt. Ist
    dagegen \(b\in N(a)\), so liefert Korollar~\ref{cor:sym-extend} ein
    \(\lambda\) mit \(\lambda(b)=\mathit{in}\) und folglich
    \(\lambda(a)=\mathit{out}\).
\end{proof}

Proposition~\ref{prop:sym-klassisch-p} formalisiert den Befund von
\textcite{Coste:2005:symmetric-af-irflx-coherent}, dass in symmetrischen
irreflexiven AFs nur zwei sehr einfache und effizient prüfbare Formen der
Akzeptanz auftreten. Nach diesem Bild müsste auch die bedingte Unabhängigkeit
leicht sein: Sie ist nach Definition~\ref{def:BU:AFs} eine Aussage über
\(\Lambda_\sigma(F)\), und alle punktweisen Fragen über \(\Lambda_\sigma(F)\)
sind trivial.

\section{Die eigentliche Schwierigkeit: exakte Realisierbarkeit}
\label{sec:sym-realisierbarkeit}

Die bedingte Unabhängigkeit fragt nicht nach dem Label eines einzelnen
Arguments, sondern danach, ob ein \emph{vollständig vorgeschriebenes
Teillabelling} realisiert werden kann. Genau hier bricht die Einfachheit
zusammen.

\begin{Definition}
    \label{def:sym-real}
    Das Realisierbarkeitsproblem \(\textit{Real}^{\mathrm{sym}}_{\mathrm{st}}\)
    ist wie folgt definiert:
    \[
    \begin{array}{ll}
    \text{Input:}  & \text{symmetrisches, irreflexives AF } F=(A,R),\
                     W\subseteq A,\ \eta:W\to\{\mathit{in},\mathit{out}\}\\
    \text{Output:} & \text{TRUE} \Longleftrightarrow
                     \exists\lambda\in\Lambda_{\mathrm{st}}(F):\lambda|_W=\eta.
    \end{array}
    \]
\end{Definition}

Der Unterschied zur kredulösen Akzeptanz ist der Zwang, Argumente auch
\emph{ausschließen} zu müssen: Ein Argument \(a\) mit \(\eta(a)=\mathit{out}\)
verlangt einen Nachbarn im \(\mathit{in}\)-Teil, und mehrere solche Forderungen
müssen gleichzeitig durch eine einzige konfliktfreie Menge erfüllt werden.

\begin{Proposition}
    \label{prop:sym-real-npc}
    \(\textit{Real}^{\mathrm{sym}}_{\mathrm{st}}\) ist NP-vollständig.
\end{Proposition}

\begin{proof}
    Zugehörigkeit: Man rät \(\lambda\) und verifiziert in Polynomialzeit
    (Proposition~\ref{prop:sym-klassisch-p}(2)), dass
    \(\lambda\in\Lambda_{\mathrm{st}}(F)\) und \(\lambda|_W=\eta\).

    Härte: Wir reduzieren von 3-SAT. Sei \(\varphi=\bigwedge_{k=1}^{m}C_k\) eine
    KNF-Formel über den Variablen \(y_1,\dots,y_l\). Wir definieren ein
    symmetrisches irreflexives AF \(F_\varphi=(A,R)\) durch
    \[
        A:=\{w_j,\overline{w_j}\mid 1\leq j\leq l\}\cup\{c_k\mid 1\leq k\leq m\}
    \]
    und die (symmetrisch abgeschlossene) Angriffsrelation, die genau die
    folgenden Konflikte enthält:
    \begin{itemize}
        \item \(\{w_j,\overline{w_j}\}\) für jedes \(j\),
        \item \(\{w_j,c_k\}\), falls \(y_j\) positiv in \(C_k\) vorkommt,
        \item \(\{\overline{w_j},c_k\}\), falls \(y_j\) negativ in \(C_k\)
        vorkommt.
    \end{itemize}
    Wir setzen \(W:=\{c_1,\dots,c_m\}\) und \(\eta\equiv\mathit{out}\).

    Sei \(\lambda\in\Lambda_{\mathrm{st}}(F_\varphi)\) mit \(\lambda|_W=\eta\)
    und \(S:=\mathit{in}(\lambda)\). Da \(S\) konfliktfrei ist, gilt nicht
    gleichzeitig \(w_j\in S\) und \(\overline{w_j}\in S\); wir definieren
    \(\beta(y_j):=\top\) genau dann, wenn \(w_j\in S\). Nach
    Proposition~\ref{prop:sym-naiv-st-pr} ist \(S\) naiv, also besitzt jedes
    \(c_k\notin S\) einen Nachbarn in \(S\). Die Nachbarn von \(c_k\) sind
    ausschließlich Literalargumente, also enthält \(S\) ein Literalargument von
    \(C_k\). Ist dies \(w_j\), so gilt \(\beta(y_j)=\top\) und \(y_j\in C_k\);
    ist es \(\overline{w_j}\), so ist \(w_j\notin S\), also
    \(\beta(y_j)=\bot\) und \(\neg y_j\in C_k\). In beiden Fällen ist \(C_k\)
    unter \(\beta\) erfüllt. Somit ist \(\varphi\) erfüllbar.

    Sei umgekehrt \(\beta\) eine erfüllende Belegung und
    \(T:=\{w_j\mid\beta(y_j)=\top\}\cup\{\overline{w_j}\mid\beta(y_j)=\bot\}\).
    \(T\) ist konfliktfrei, denn \(T\) enthält aus jedem Paar genau ein
    Argument und die Paare stehen untereinander nicht in Konflikt. Jedes
    \(c_k\) besitzt ein unter \(\beta\) wahres Literal, dessen Literalargument
    in \(T\) liegt. Nach Korollar~\ref{cor:sym-extend} existiert ein
    \(\lambda\in\Lambda_{\mathrm{st}}(F_\varphi)\) mit
    \(T\subseteq\mathit{in}(\lambda)\), und alle \(c_k\) tragen dort
    \(\mathit{out}\). Also gilt \(\lambda|_W=\eta\).

    Die Konstruktion ist offensichtlich polynomiell.
\end{proof}

Proposition~\ref{prop:sym-real-npc} erklärt, warum die Tractability der
klassischen Aufgaben nicht auf die Unabhängigkeit durchschlägt: Der innere
Existenzquantor aus Definition~\ref{def:BU:AFs} ist bereits auf dieser Klasse
NP-vollständig. Was noch fehlt, ist ein Mechanismus, der den äußeren
Allquantor über \(\lambda_1,\lambda_2\) in einen echten \(\forall\)-Quantor
über Belegungen verwandelt. Dieser Mechanismus ist der Gegenstand des nächsten
Abschnitts.

\section{Hauptresultat}
\label{sec:sym-hauptresultat}

Die zentrale konstruktive Schwierigkeit in symmetrischen AFs ist, dass jeder
Angriff in beide Richtungen wirkt. Zwei Konsequenzen sind dabei entscheidend.

Erstens lässt sich die für Reduktionen übliche Invariante \enquote{aus jedem
Paar \(v,\overline{v}\) ist genau eines \(\mathit{in}\)} nicht mehr durch einen
bloßen Zweierzyklus erzwingen: Sobald \(v\) oder \(\overline{v}\) weitere
Nachbarn besitzen, können in einer naiven Menge beide fehlen, weil sie von
anderer Seite dominiert werden. Wir lösen das, indem wir Wert- und Signalknoten
trennen und zu einem Viererkreis verbinden.

Zweitens lässt sich die in gerichteten Konstruktionen übliche Aggregation
\enquote{\(\varphi\) ist \(\mathit{in}\), also sind alle Klauselargumente
\(\mathit{out}\), also ist jede Klausel erfüllt} nicht nachbauen: Ein
Aggregationsargument \(\varphi\), das alle Klauselargumente angreift,
dominiert diese selbst und macht damit jede Forderung an die Literale
gegenstandslos. Wir verzichten deshalb auf ein Aggregationsargument und
schreiben die Klauselargumente stattdessen \emph{direkt} als
\(\mathit{out}\)-Teil der Unabhängigkeitsanfrage vor. Damit die Menge der so
vorschreibbaren Muster klein und beherrschbar bleibt, verbinden wir die
Klauselargumente zu einer Clique.

\begin{Definition}
    \label{def:sym-konstruktion}
    Sei \(\Psi=\forall X\exists Y\,\varphi\) mit
    \(X=\{x_1,\dots,x_n\}\), \(Y=\{y_1,\dots,y_l\}\) und
    \(\varphi=\bigwedge_{k=1}^{m}C_k\) in KNF, \(m\geq 1\). Wir setzen
    \(X^{+}:=X\cup\{x_0\}\) mit einer frischen universellen Variablen \(x_0\)
    und
    \[
        \widehat{\varphi}:=\bigwedge_{k=1}^{m}\bigl(x_0\lor C_k\bigr).
    \]
    Das AF \(G_\Psi=(A_\Psi,R_\Psi)\) besitzt die Argumentmenge
    \[
    \begin{aligned}
        A_\Psi:={}&
        \{u_i,\overline{u_i},s_i,\overline{s_i}\mid x_i\in X^{+}\}
        \;\cup\;
        \{w_j,\overline{w_j}\mid y_j\in Y\}
        \;\cup\;
        \{c_1,\dots,c_m\},
    \end{aligned}
    \]
    und \(R_\Psi\) ist die symmetrische Hülle genau der folgenden Konflikte:
    \begin{enumerate}
        \item[(K1)] \(\{u_i,\overline{u_i}\}\),
        \(\{u_i,\overline{s_i}\}\),
        \(\{\overline{u_i},s_i\}\),
        \(\{s_i,\overline{s_i}\}\)
        \quad für alle \(x_i\in X^{+}\),
        \item[(K2)] \(\{w_j,\overline{w_j}\}\) \quad für alle \(y_j\in Y\),
        \item[(K3)] \(\{s_i,c_k\}\), falls \(x_i\) positiv in \(C_k\) vorkommt,
        und \(\{\overline{s_i},c_k\}\), falls \(x_i\) negativ in \(C_k\)
        vorkommt \quad(\(1\leq i\leq n\)); zusätzlich \(\{s_0,c_k\}\) für alle
        \(k\),
        \item[(K4)] \(\{w_j,c_k\}\), falls \(y_j\) positiv in \(C_k\) vorkommt,
        und \(\{\overline{w_j},c_k\}\), falls \(y_j\) negativ in \(C_k\)
        vorkommt,
        \item[(K5)] \(\{c_k,c_{k'}\}\) für alle \(k\neq k'\).
    \end{enumerate}
    Als Unabhängigkeitsanfrage setzen wir
    \[
        \mathcal{X}:=\{c_1,\dots,c_m\},\qquad
        \mathcal{Y}:=\{u_i,\overline{u_i}\mid x_i\in X^{+}\},\qquad
        Z:=\emptyset.
    \]
\end{Definition}

Das AF \(G_\Psi\) ist nach Konstruktion symmetrisch und irreflexiv, besitzt
\(4(n+1)+2l+m\) Argumente und ist in Polynomialzeit berechenbar. Die Blöcke
(K1) bilden für jede universelle Variable einen Viererkreis
\(u_i-\overline{u_i}-s_i-\overline{s_i}-u_i\); die Wertknoten
\(u_i,\overline{u_i}\) besitzen \emph{keine} weiteren Nachbarn, während die
Signalknoten \(s_i,\overline{s_i}\) die Verbindung zu den Klauselargumenten
herstellen.

\begin{Lemma}[Wertknoten]
    \label{lem:sym-wertknoten}
    Sei \(\lambda\in\Lambda_{\mathrm{st}}(G_\Psi)\) und
    \(S:=\mathit{in}(\lambda)\). Dann gilt für jedes \(x_i\in X^{+}\):
    \(|S\cap\{u_i,\overline{u_i}\}|=1\).
\end{Lemma}

\begin{proof}
    Wegen \(\{u_i,\overline{u_i}\}\in R_\Psi\) und der Konfliktfreiheit von
    \(S\) liegen nicht beide in \(S\). Seien nun beide nicht in \(S\). Nach
    Proposition~\ref{prop:sym-naiv-st-pr} ist \(S\) naiv, also besitzt
    \(u_i\) einen Nachbarn in \(S\). Die Nachbarn von \(u_i\) sind nach (K1)
    genau \(\overline{u_i}\) und \(\overline{s_i}\); folglich gilt
    \(\overline{s_i}\in S\). Analog besitzt \(\overline{u_i}\) die Nachbarn
    \(u_i\) und \(s_i\), also \(s_i\in S\). Wegen
    \(\{s_i,\overline{s_i}\}\in R_\Psi\) widerspricht das der
    Konfliktfreiheit von \(S\).
\end{proof}

Für \(\lambda\in\Lambda_{\mathrm{st}}(G_\Psi)\) mit
\(S=\mathit{in}(\lambda)\) definieren wir daher die \emph{induzierte Belegung}
\[
    \alpha_\lambda:X^{+}\to\{\top,\bot\},\quad
    \alpha_\lambda(x_i)=\top:\iff u_i\in S,
\]
sowie
\[
    \beta_\lambda:Y\to\{\top,\bot\},\quad
    \beta_\lambda(y_j)=\top:\iff w_j\in S.
\]

\begin{Lemma}[Signaltreue]
    \label{lem:sym-signaltreue}
    Sei \(\lambda\in\Lambda_{\mathrm{st}}(G_\Psi)\), \(S=\mathit{in}(\lambda)\).
    Dann gilt:
    \begin{enumerate}
        \item \(s_i\in S\Rightarrow\alpha_\lambda(x_i)=\top\) und
        \(\overline{s_i}\in S\Rightarrow\alpha_\lambda(x_i)=\bot\).
        \item \(w_j\in S\Rightarrow\beta_\lambda(y_j)=\top\) und
        \(\overline{w_j}\in S\Rightarrow\beta_\lambda(y_j)=\bot\).
    \end{enumerate}
\end{Lemma}

\begin{proof}
    (1) Aus \(s_i\in S\) und \(\{\overline{u_i},s_i\}\in R_\Psi\) folgt
    \(\overline{u_i}\notin S\), mit Lemma~\ref{lem:sym-wertknoten} also
    \(u_i\in S\) und damit \(\alpha_\lambda(x_i)=\top\). Der zweite Fall
    verläuft symmetrisch über \(\{u_i,\overline{s_i}\}\in R_\Psi\).

    (2) \(w_j\in S\) liefert \(\beta_\lambda(y_j)=\top\) definitionsgemäß; aus
    \(\overline{w_j}\in S\) folgt \(w_j\notin S\) wegen
    \(\{w_j,\overline{w_j}\}\in R_\Psi\), also \(\beta_\lambda(y_j)=\bot\).
\end{proof}

\begin{Lemma}[Realisierbare Teillabellings]
    \label{lem:sym-teile}
    Sei \(\lambda\in\Lambda_{\mathrm{st}}(G_\Psi)\) und
    \(S=\mathit{in}(\lambda)\). Dann gilt:
    \begin{enumerate}
        \item \(|S\cap\mathcal{X}|\leq 1\).
        \item Zu jedem \(k\in\{1,\dots,m\}\) und jeder Belegung
        \(\alpha:X^{+}\to\{\top,\bot\}\) existiert
        \(\lambda'\in\Lambda_{\mathrm{st}}(G_\Psi)\) mit
        \(\mathit{in}(\lambda')\cap\mathcal{X}=\{c_k\}\) und
        \(\alpha_{\lambda'}=\alpha\).
        \item Zu jeder Belegung \(\alpha:X^{+}\to\{\top,\bot\}\) existiert
        genau dann ein \(\lambda'\in\Lambda_{\mathrm{st}}(G_\Psi)\) mit
        \(\mathit{in}(\lambda')\cap\mathcal{X}=\emptyset\) und
        \(\alpha_{\lambda'}=\alpha\), wenn eine Belegung \(\beta\) von \(Y\)
        mit \(\widehat{\varphi}(\alpha,\beta)=\top\) existiert.
        \item Zu jeder Belegung \(\alpha:X^{+}\to\{\top,\bot\}\) existiert
        \(\lambda'\in\Lambda_{\mathrm{st}}(G_\Psi)\) mit
        \(\alpha_{\lambda'}=\alpha\).
    \end{enumerate}
\end{Lemma}

\begin{proof}
    Wir schreiben \(U_\alpha:=\{u_i\mid\alpha(x_i)=\top\}
    \cup\{\overline{u_i}\mid\alpha(x_i)=\bot\}\) und
    \(S_\alpha:=\{s_i\mid\alpha(x_i)=\top\}
    \cup\{\overline{s_i}\mid\alpha(x_i)=\bot\}\).
    Nach (K1) ist \(U_\alpha\cup S_\alpha\) konfliktfrei: Innerhalb eines
    Viererkreises werden nur gegenüberliegende Knoten gewählt
    (\(u_i\) mit \(s_i\) beziehungsweise \(\overline{u_i}\) mit
    \(\overline{s_i}\)), und zwischen verschiedenen Blöcken bestehen keine
    Konflikte. Ferner besitzt jedes Argument aus
    \(\{u_i,\overline{u_i}\}\setminus U_\alpha\) seinen Partner in
    \(U_\alpha\) als Nachbarn. Nach Korollar~\ref{cor:sym-extend} gilt
    daher für jede Erweiterung eines \(U_\alpha\) enthaltenden konfliktfreien
    Kerns zu einem \textit{stable}-Labelling \(\lambda'\) bereits
    \(\alpha_{\lambda'}=\alpha\).

    (1) Die Klauselargumente bilden nach (K5) eine Clique; eine konfliktfreie
    Menge enthält daher höchstens eines von ihnen.

    (2) Der Kern \(T:=U_\alpha\cup\{c_k\}\) ist konfliktfrei, da Wertknoten
    nach (K1) nur Wert- und Signalknoten desselben Blocks als Nachbarn haben
    und insbesondere zu keinem Klauselargument benachbart sind. Sei
    \(\lambda'\) eine Erweiterung gemäß Korollar~\ref{cor:sym-extend}. Dann gilt
    \(\alpha_{\lambda'}=\alpha\) nach der Vorbemerkung, und jedes \(c_{k'}\)
    mit \(k'\neq k\) besitzt den Nachbarn \(c_k\in T\), ist also
    \(\mathit{out}\). Somit ist
    \(\mathit{in}(\lambda')\cap\mathcal{X}=\{c_k\}\).

    (3) \enquote{\(\Leftarrow\)}: Sei \(\beta\) mit
    \(\widehat{\varphi}(\alpha,\beta)=\top\) und sei
    \[
        T:=U_\alpha\cup S_\alpha\cup
        \{w_j\mid\beta(y_j)=\top\}\cup\{\overline{w_j}\mid\beta(y_j)=\bot\}.
    \]
    \(T\) ist konfliktfrei: Konflikte bestehen nach (K3), (K4) nur zwischen
    Signal- beziehungsweise Literalargumenten und Klauselargumenten, und
    \(T\) enthält keine Klauselargumente; innerhalb der Blöcke wurde die
    Konfliktfreiheit bereits festgestellt. Da \(\widehat{\varphi}\) unter
    \((\alpha,\beta)\) erfüllt ist, besitzt jede Klausel \(x_0\lor C_k\) ein
    wahres Literal, dessen zugehöriges Signal- oder Literalargument nach (K3),
    (K4) in \(T\) liegt und zu \(c_k\) benachbart ist. Nach
    Korollar~\ref{cor:sym-extend} existiert ein \textit{stable}-Labelling
    \(\lambda'\) mit \(T\subseteq\mathit{in}(\lambda')\); dort sind alle
    \(c_k\) \(\mathit{out}\) und es gilt \(\alpha_{\lambda'}=\alpha\).

    \enquote{\(\Rightarrow\)}: Sei \(\lambda'\) mit
    \(\mathit{in}(\lambda')\cap\mathcal{X}=\emptyset\) und
    \(\alpha_{\lambda'}=\alpha\), und sei \(S'=\mathit{in}(\lambda')\). Da
    \(S'\) naiv ist, besitzt jedes \(c_k\) einen Nachbarn in \(S'\). Die
    Nachbarn von \(c_k\) sind nach (K3)--(K5) Signalargumente,
    Literalargumente und andere Klauselargumente; letztere liegen nicht in
    \(S'\). Also enthält \(S'\) ein Signal- oder Literalargument zu einem
    Literal von \(x_0\lor C_k\). Nach Lemma~\ref{lem:sym-signaltreue} ist
    dieses Literal unter \((\alpha,\beta_{\lambda'})\) wahr. Da \(k\) beliebig
    war, gilt \(\widehat{\varphi}(\alpha,\beta_{\lambda'})=\top\).

    (4) Man erweitere \(U_\alpha\) nach Korollar~\ref{cor:sym-extend}.
\end{proof}

\begin{Theorem}
    \label{th:sym-st-pi2p}
    \(\textit{Ind}^{\mathrm{sym}}_{\mathrm{st}}\) ist
    \(\Pi_2^\text{P}\)-vollständig. Die Härte gilt bereits für
    \(Z=\emptyset\).
\end{Theorem}

\begin{proof}
    \emph{Zugehörigkeit.} Symmetrische irreflexive AFs bilden eine Teilklasse
    aller AFs; die obere Schranke folgt daher aus
    Theorem~\ref{th:ind-st-is-pi2p-c}.

    \emph{Härte.} Sei \(\Psi=\forall X\exists Y\,\varphi\) eine Instanz des
    \(\Pi_2^\text{P}\)-vollständigen Problems \(\mathrm{QBF}^2_\forall\) und sei
    \(G_\Psi\) mit \(\mathcal{X},\mathcal{Y},Z=\emptyset\) wie in
    Definition~\ref{def:sym-konstruktion}. Die Mengen \(\mathcal{X}\),
    \(\mathcal{Y}\) und \(Z\) sind paarweise disjunkt, und \(G_\Psi\) ist
    symmetrisch und irreflexiv. Wegen \(Z=\emptyset\) ist die Voraussetzung
    \(\lambda_1|_Z=\lambda_2|_Z\) aus Definition~\ref{def:BU:AFs} für jedes
    Paar erfüllt. Wir zeigen
    \[
        \Psi\text{ ist gültig}
        \quad\Longleftrightarrow\quad
        \mathcal{X}\bot_{\mathrm{st}}\mathcal{Y}\mid\emptyset
        \text{ in }G_\Psi.
    \]

    Vorab halten wir fest, dass wegen
    \(\widehat{\varphi}=\bigwedge_k(x_0\lor C_k)\) gilt:
    \(\Psi\) ist genau dann gültig, wenn
    \(\forall X^{+}\exists Y\,\widehat{\varphi}\) gültig ist. Für
    \(\alpha(x_0)=\top\) ist \(\widehat{\varphi}\) unter jeder Belegung von
    \(Y\) erfüllt, für \(\alpha(x_0)=\bot\) gilt
    \(\widehat{\varphi}\equiv\varphi\).

    \enquote{\(\Rightarrow\)}: Sei \(\Psi\) gültig und seien
    \(\lambda_1,\lambda_2\in\Lambda_{\mathrm{st}}(G_\Psi)\). Sei
    \(\alpha:=\alpha_{\lambda_2}\); nach Lemma~\ref{lem:sym-wertknoten} ist
    \(\lambda_2|_{\mathcal{Y}}\) durch \(\alpha\) eindeutig bestimmt und
    umgekehrt. Nach Lemma~\ref{lem:sym-teile}(1) ist
    \(\mathit{in}(\lambda_1)\cap\mathcal{X}\) entweder leer oder von der Form
    \(\{c_k\}\).

    Im Fall \(\mathit{in}(\lambda_1)\cap\mathcal{X}=\{c_k\}\) liefert
    Lemma~\ref{lem:sym-teile}(2) ein \(\lambda_3\) mit
    \(\mathit{in}(\lambda_3)\cap\mathcal{X}=\{c_k\}\) und
    \(\alpha_{\lambda_3}=\alpha\), also
    \(\lambda_3|_{\mathcal{X}}=\lambda_1|_{\mathcal{X}}\) und
    \(\lambda_3|_{\mathcal{Y}}=\lambda_2|_{\mathcal{Y}}\).

    Im Fall \(\mathit{in}(\lambda_1)\cap\mathcal{X}=\emptyset\) existiert wegen
    der Gültigkeit von \(\forall X^{+}\exists Y\,\widehat{\varphi}\) eine
    Belegung \(\beta\) mit \(\widehat{\varphi}(\alpha,\beta)=\top\);
    Lemma~\ref{lem:sym-teile}(3) liefert das gesuchte \(\lambda_3\).

    Damit gilt \(\mathcal{X}\bot_{\mathrm{st}}\mathcal{Y}\mid\emptyset\).

    \enquote{\(\Leftarrow\)}: Sei \(\Psi\) nicht gültig. Dann existiert eine
    Belegung \(\alpha\) von \(X\), sodass \(\varphi(\alpha,\beta)=\bot\) für
    alle \(\beta\) gilt. Wir setzen
    \(\alpha^{+}:=\alpha\cup\{x_0\mapsto\bot\}\); dann gilt
    \(\widehat{\varphi}(\alpha^{+},\beta)=\bot\) für alle \(\beta\).

    Wähle \(\lambda_2\) nach Lemma~\ref{lem:sym-teile}(4) mit
    \(\alpha_{\lambda_2}=\alpha^{+}\). Wähle \(\lambda_1\) nach
    Lemma~\ref{lem:sym-teile}(3), angewandt auf eine beliebige Belegung
    \(\alpha'\) mit \(\alpha'(x_0)=\top\); wegen
    \(\widehat{\varphi}(\alpha',\beta)=\top\) für jedes \(\beta\) existiert ein
    solches \(\lambda_1\) mit
    \(\mathit{in}(\lambda_1)\cap\mathcal{X}=\emptyset\).

    Angenommen, es existierte \(\lambda_3\in\Lambda_{\mathrm{st}}(G_\Psi)\) mit
    \(\lambda_3|_{\mathcal{X}}=\lambda_1|_{\mathcal{X}}\) und
    \(\lambda_3|_{\mathcal{Y}}=\lambda_2|_{\mathcal{Y}}\). Dann gilt
    \(\mathit{in}(\lambda_3)\cap\mathcal{X}=\emptyset\) und
    \(\alpha_{\lambda_3}=\alpha^{+}\). Nach Lemma~\ref{lem:sym-teile}(3)
    existierte dann ein \(\beta\) mit
    \(\widehat{\varphi}(\alpha^{+},\beta)=\top\), im Widerspruch zur Wahl von
    \(\alpha\). Also gilt
    \(\mathcal{X}\not\bot_{\mathrm{st}}\mathcal{Y}\mid\emptyset\).

    Die Konstruktion ist in Polynomialzeit berechenbar, womit die Härte gezeigt
    ist.
\end{proof}

\begin{Corollary}
    \label{cor:sym-nicht-ocf}
    Theorem~\ref{th:sym-st-pi2p} folgt nicht aus
    Proposition~\ref{prop:ocf-af-ci-st-pi^2}.
\end{Corollary}

\begin{proof}
    Für \(m\geq 3\) enthält \(G_\Psi\) wegen (K5) den Dreierkreis
    \(c_1\to c_2\to c_3\to c_1\) und ist damit nicht ungeradezykelfrei. Ebenso
    ist \(G_\Psi\) wegen dieses Dreiecks nicht bipartit. Die Klasse der
    symmetrischen irreflexiven AFs ist somit weder in der Klasse der
    ungeradezykelfreien noch in der der bipartiten AFs enthalten.
\end{proof}

Bemerkenswert ist, dass symmetrische irreflexive AFs dennoch kohärent sind: Die
Kohärenz wird hier nicht über die Abwesenheit ungerader Zyklen erreicht, sondern
über die Symmetrie selbst (Lemma~\ref{lem:sym-cf-adm}). Die Zyklenfreiheit ist
also hinreichend, aber nicht notwendig für Kohärenz, und für die Komplexität
der Unabhängigkeit ist offenkundig nicht die Kohärenz entscheidend.

\section{Konsequenz für die \textit{preferred}-Semantik}
\label{sec:sym-preferred}

\begin{Lemma}[Transferlemma]
    \label{lem:sym-transfer}
    Sei \(F=(A,R)\) ein AF und seien \(\sigma,\tau\) Semantiken mit
    \(\Lambda_\sigma(F)=\Lambda_\tau(F)\). Dann gilt für alle paarweise
    disjunkten \(X,Y,Z\subseteq A\):
    \(X\bot_\sigma Y\mid Z\) genau dann, wenn \(X\bot_\tau Y\mid Z\).
\end{Lemma}

\begin{proof}
    Die definierende Bedingung aus Definition~\ref{def:BU:AFs} quantifiziert
    ausschließlich über Elemente von \(\Lambda_\sigma(F)\). Ersetzt man diese
    Menge durch die nach Voraussetzung identische Menge \(\Lambda_\tau(F)\), so
    geht die Bedingung Wort für Wort in die entsprechende Bedingung für
    \(\tau\) über.
\end{proof}

\begin{Corollary}
    \label{cor:sym-pr-pi2p}
    \(\textit{Ind}^{\mathrm{sym}}_{\mathrm{pr}}\) ist
    \(\Pi_2^\text{P}\)-vollständig.
\end{Corollary}

\begin{proof}
    Nach Proposition~\ref{prop:sym-naiv-st-pr} gilt
    \(\Lambda_{\mathrm{pr}}(F)=\Lambda_{\mathrm{st}}(F)\) für jedes
    symmetrische irreflexive AF \(F\). Mit Lemma~\ref{lem:sym-transfer} sind
    \(\textit{Ind}^{\mathrm{sym}}_{\mathrm{pr}}\) und
    \(\textit{Ind}^{\mathrm{sym}}_{\mathrm{st}}\) dieselbe Sprache; die
    Behauptung folgt aus Theorem~\ref{th:sym-st-pi2p}.
\end{proof}

Für \(\sigma=\mathrm{pr}\) ist die Klasse also \emph{echt} einfacher: Das
allgemeine Problem ist nach Theorem~\ref{th:ind-pr-is-pi3p-c}
\(\Pi_3^\text{P}\)-vollständig, auf symmetrischen irreflexiven AFs sinkt es um
eine Stufe der polynomiellen Hierarchie. Der Grund ist strukturell: Die
Maximalitätsbedingung der \textit{preferred}-Semantik, die im allgemeinen Fall
den zusätzlichen Quantorenwechsel erzwingt, kollabiert hier zu einer in
Polynomialzeit prüfbaren Bedingung (Proposition~\ref{prop:sym-klassisch-p}(2)).
Unter der üblichen Annahme \(\Pi_2^\text{P}\neq\Pi_3^\text{P}\) ist dieser
Gewinn nicht nur formal.

Für \(\sigma=\mathrm{st}\) tritt dagegen kein Gewinn ein. Die Erklärung liefert
Proposition~\ref{prop:sym-real-npc}: Die Symmetrie vereinfacht alle
\emph{punktweisen} Fragen über \(\Lambda_{\mathrm{st}}(F)\), lässt aber die
\emph{simultane} Realisierbarkeit vorgeschriebener Teillabellings NP-hart. Da
die bedingte Unabhängigkeit genau von dieser simultanen Realisierbarkeit
handelt, überträgt sich die Vereinfachung nicht.

\section{Ein in Polynomialzeit entscheidbares Fragment}
\label{sec:sym-fragment}

Die Härte in Theorem~\ref{th:sym-st-pi2p} benötigt eine unbeschränkt große
Menge \(\mathcal{Y}\). Beschränkt man beide Mengen auf einzelne Argumente und
setzt \(Z=\emptyset\), so wird das Problem tatsächlich einfach. Damit ist die
Klasse in einem präzisen Sinn doch einfacher als der allgemeine Fall, in dem
schon die zugrundeliegenden Realisierbarkeitsfragen NP-vollständig sind.

\begin{Proposition}
    \label{prop:sym-singleton-p}
    Sei \(F=(A,R)\) symmetrisch und irreflexiv und seien \(x,y\in A\) mit
    \(x\neq y\). Dann ist \(\{x\}\bot_{\mathrm{st}}\{y\}\mid\emptyset\) in Zeit
    \(\mathcal{O}(|A|^2+|A|\cdot|R|)\) entscheidbar.
\end{Proposition}

\begin{proof}
    Wegen \(Z=\emptyset\) ist die Bedingung \(\lambda_1|_Z=\lambda_2|_Z\)
    stets erfüllt. Setze
    \[
        \mathcal{R}:=\bigl\{(\lambda(x),\lambda(y))\;\big|\;
        \lambda\in\Lambda_{\mathrm{st}}(F)\bigr\}
        \subseteq\{\mathit{in},\mathit{out}\}^2 .
    \]
    Nach Definition~\ref{def:BU:AFs} gilt
    \(\{x\}\bot_{\mathrm{st}}\{y\}\mid\emptyset\) genau dann, wenn aus
    \((\alpha_1,\beta_1),(\alpha_2,\beta_2)\in\mathcal{R}\) stets
    \((\alpha_1,\beta_2)\in\mathcal{R}\) folgt. Da \(\mathcal{R}\) höchstens
    vier Elemente besitzt, genügt es, die vier Zugehörigkeiten zu bestimmen.
    Wir benutzen durchgehend Proposition~\ref{prop:sym-naiv-st-pr} und
    Korollar~\ref{cor:sym-extend}.

    \emph{\((\mathit{in},\mathit{in})\).} Genau dann realisierbar, wenn
    \((x,y)\notin R\): Ist \(\{x,y\}\) konfliktfrei, erweitere man es; andernfalls
    verbietet die Konfliktfreiheit beide gemeinsam.

    \emph{\((\mathit{in},\mathit{out})\).} Ist \((x,y)\in R\), so erweitere man
    \(\{x\}\); dann ist \(y\) \(\mathit{out}\). Ist \((x,y)\notin R\), so ist
    das Paar genau dann realisierbar, wenn \(N(y)\setminus N[x]\neq\emptyset\).
    Existiert nämlich ein solches \(w\), so ist \(\{x,w\}\) konfliktfrei und
    jede Erweiterung enthält \(x\), aber wegen \(w\in N(y)\) nicht \(y\).
    Existiert umgekehrt \(\lambda\) mit \(\lambda(x)=\mathit{in}\) und
    \(\lambda(y)=\mathit{out}\), so besitzt \(y\) einen Nachbarn
    \(w\in\mathit{in}(\lambda)\); wegen \(x\notin N(y)\) ist \(w\neq x\), und
    wegen der Konfliktfreiheit gilt \(w\notin N(x)\).

    \emph{\((\mathit{out},\mathit{in})\).} Symmetrisch mit vertauschten Rollen.

    \emph{\((\mathit{out},\mathit{out})\).} Genau dann realisierbar, wenn
    \(w_1\in N(x)\setminus\{y\}\) und \(w_2\in N(y)\setminus\{x\}\) mit
    \(w_1=w_2\) oder \((w_1,w_2)\notin R\) existieren. Sind solche \(w_1,w_2\)
    gegeben, so ist \(\{w_1,w_2\}\) konfliktfrei, und jede Erweiterung schließt
    \(x\) und \(y\) aus. Existiert umgekehrt \(\lambda\) mit
    \(\lambda(x)=\lambda(y)=\mathit{out}\), so besitzen \(x\) und \(y\) je einen
    Nachbarn \(w_1,w_2\in\mathit{in}(\lambda)\); diese liegen wegen
    \(\lambda(x)=\lambda(y)=\mathit{out}\) nicht in \(\{x,y\}\) und stehen
    untereinander nicht in Konflikt.

    Alle vier Tests sind durch Absuchen von \(N(x)\), \(N(y)\) beziehungsweise
    aller Paare \((w_1,w_2)\in N(x)\times N(y)\) in der angegebenen Zeit
    durchführbar. Anschließend prüft man die Rechteckbedingung auf der höchstens
    vierelementigen Menge \(\mathcal{R}\) in konstanter Zeit.
\end{proof}

Im allgemeinen Fall sind bereits die vier Realisierbarkeitstests
NP-vollständig, da \((\mathit{in},\cdot)\) die kredulöse Akzeptanz unter
\textit{stable}-Semantik enthält. Proposition~\ref{prop:sym-singleton-p} ist
also eine echte Vereinfachung; sie zeigt zugleich, dass die Härte aus
Theorem~\ref{th:sym-st-pi2p} nicht aus einzelnen Argumentpaaren stammt, sondern
aus der \emph{gemeinsamen} Vorschrift vieler Labels.

Als Konsistenzprüfung sei angemerkt, dass
Proposition~\ref{prop:cdg-afs-not-st-ci} ein Spezialfall dieser Analyse ist: Ein
vollständig gerichtetes AF ist symmetrisch und irreflexiv, seine naiven Mengen
sind genau die einelementigen Teilmengen von \(A\), und für \(x\neq y\) ist
\((\mathit{in},\mathit{in})\) wegen \((x,y)\in R\) nicht realisierbar, während
\((\mathit{in},\mathit{out})\) und \((\mathit{out},\mathit{in})\) realisierbar
sind. Die Rechteckbedingung ist damit verletzt.

\section{Die \textit{complete}-Semantik bleibt offen}
\label{sec:sym-complete}

Für \(\sigma=\mathrm{co}\) lässt sich die Konstruktion aus
Definition~\ref{def:sym-konstruktion} nicht übernehmen. Wir halten zunächst die
Struktur der \textit{complete}-Labellings fest.

\begin{Proposition}
    \label{prop:sym-co-struktur}
    Sei \(F=(A,R)\) symmetrisch und irreflexiv und sei \(\lambda\) ein
    Labelling mit \(S:=\mathit{in}(\lambda)\). Dann ist \(\lambda\) genau dann
    \textit{complete}, wenn \(S\) konfliktfrei ist,
    \(\mathit{out}(\lambda)=N(S)\) gilt und kein \(a\in A\setminus S\) mit
    \(N(a)\subseteq N(S)\) existiert.
\end{Proposition}

\begin{proof}
    Nach Definition~\ref{def:complete-labelling} ist \(a\) genau dann
    \(\mathit{out}\), wenn \(a\) einen \(\mathit{in}\)-Angreifer besitzt, also
    genau dann, wenn \(a\in N(S)\); das ist gleichbedeutend mit
    \(\mathit{out}(\lambda)=N(S)\), und \(S\cap N(S)=\emptyset\) ist genau die
    Konfliktfreiheit von \(S\). Ferner ist \(a\) genau dann \(\mathit{in}\),
    wenn alle Angreifer \(\mathit{out}\) sind, also genau dann, wenn
    \(N(a)\subseteq N(S)\). Für \(a\in S\) ist diese Bedingung wegen der
    Symmetrie automatisch erfüllt; die verbleibende Forderung ist, dass kein
    \(a\) außerhalb von \(S\) sie erfüllt.
\end{proof}

Insbesondere ist das \textit{grounded}-Labelling durch
\(\mathit{in}(\lambda_{\mathrm{gr}})=\{a\in A\mid N(a)=\emptyset\}\) gegeben:
Isolierte Argumente sind unangegriffen, und ein nicht isoliertes Argument \(a\)
besitzt einen Nachbarn, der von einem isolierten Argument nicht angegriffen
werden kann.

Damit scheitert die Reduktion aus Abschnitt~\ref{sec:sym-hauptresultat} an
Lemma~\ref{lem:sym-wertknoten}: Wählt man \(S=\emptyset\), so ist
\(N(S)=\emptyset\), und nach Proposition~\ref{prop:sym-co-struktur} ist das
Labelling, das alle nicht isolierten Argumente mit \(\mathit{und}\) versieht,
\textit{complete}. In diesem Labelling tragen sämtliche Wertknoten
\(u_i,\overline{u_i}\) das Label \(\mathit{und}\); die Zuordnung
\enquote{\(\lambda_2\) kodiert eine totale Belegung} bricht zusammen, und der
äußere Allquantor läuft über wesentlich mehr Labellings als über die
Belegungen von \(X^{+}\). Da zusätzliche \(\lambda_2\) zusätzliche
Kombinationsforderungen erzeugen, kann die Richtung
\enquote{\(\Psi\) gültig \(\Rightarrow\) unabhängig} verletzt werden.

In gerichteten Konstruktionen wird dieses Problem durch Detektorargumente
gelöst, die partielle Labellings erkennen und die eigentliche Formelauswertung
neutralisieren --- so etwa die Argumente \(d_i\), \(s\) und \(g\) im Beweis von
Proposition~\ref{prop:ocf-af-ci-co-pi^2-c}. Ein solcher Detektor benötigt einen
gerichteten Informationsfluss \(\{x_i,\overline{x_i}\}\to d_i\to s\to g\), der
sich in symmetrischen AFs nicht nachbilden lässt: Jeder Detektor griffe seine
Quelle zurück an. Wir halten daher fest:

\begin{Proposition}
    \label{prop:sym-co-obere-schranke}
    \(\textit{Ind}^{\mathrm{sym}}_{\mathrm{co}}\in\Pi_2^\text{P}\).
\end{Proposition}

\begin{proof}
    Folgt aus Theorem~\ref{th:ind-co-is-pi2p-c}, da symmetrische irreflexive
    AFs eine Teilklasse aller AFs bilden.
\end{proof}

Ob \(\textit{Ind}^{\mathrm{sym}}_{\mathrm{co}}\) auch \(\Pi_2^\text{P}\)-hart
ist, bleibt offen. Es ist die einzige der drei betrachteten Semantiken, für die
in dieser Klasse keine vollständige Einordnung erreicht wurde.

\section{Zusammenfassung}
\label{sec:sym-zusammenfassung}

\begin{table}[htbp]
    \myfloatalign
    \centering
    \renewcommand{\arraystretch}{1.15}
    \setlength{\tabcolsep}{1.0em}
    \begin{tabular}{lccc}
        \toprule
        & \(co\) & \(st\) & \(pr\) \\
        \midrule
        allgemeine AFs
            & \(\Pi_2^\text{P}\)-c
            & \(\Pi_2^\text{P}\)-c
            & \(\Pi_3^\text{P}\)-c \\
        symmetrisch, irreflexiv
            & in \(\Pi_2^\text{P}\)
            & \(\Pi_2^\text{P}\)-c
            & \(\Pi_2^\text{P}\)-c \\
        \quad davon \(|X|=|Y|=1\), \(Z=\emptyset\)
            & ---
            & in P
            & in P \\
        \bottomrule
    \end{tabular}
    \caption[Komplexität der Unabhängigkeit auf symmetrischen irreflexiven
    AFs]{Einordnung von \(\sigma\)-Unabhängigkeit auf symmetrischen
    irreflexiven AFs im Vergleich zum allgemeinen Fall. Die Einträge der
    zweiten Zeile stammen aus Proposition~\ref{prop:sym-co-obere-schranke},
    Theorem~\ref{th:sym-st-pi2p} und Korollar~\ref{cor:sym-pr-pi2p}, der
    Eintrag der dritten Zeile aus Proposition~\ref{prop:sym-singleton-p} in
    Verbindung mit Korollar~\ref{cor:sym-pr-pi2p}.}
    \label{tab:sym-komplexitaet}
\end{table}

Die Antwort auf die Ausgangsfrage ist differenziert. Symmetrische irreflexive
AFs sind für die \textit{preferred}-Semantik echt einfacher, weil dort die
Maximalitätsbedingung kollabiert und ein Quantorenwechsel entfällt. Für die
\textit{stable}-Semantik sind sie \emph{nicht} einfacher: Trotz kohärenter
Struktur, garantierter Existenz von Labellings und durchgängig polynomieller
klassischer Entscheidungsprobleme bleibt die bedingte Unabhängigkeit
\(\Pi_2^\text{P}\)-vollständig. Eine Vereinfachung tritt erst ein, wenn man
zusätzlich die Größe der Anfragemengen beschränkt.

Methodisch ist das der interessanteste Befund dieses Kapitels: Die
Tractability der klassischen Aufgaben ist kein verlässlicher Indikator für die
Komplexität der bedingten Unabhängigkeit. Maßgeblich ist stattdessen, ob
vorgeschriebene Teillabellings \emph{simultan} realisierbar sind
(Proposition~\ref{prop:sym-real-npc}) und ob der äußere Allquantor über
hinreichend viele, strukturell unterscheidbare Labellings läuft
(Lemma~\ref{lem:sym-wertknoten}). Beide Bedingungen sind in symmetrischen
irreflexiven AFs erfüllt, obwohl jede einzelne Akzeptanzfrage trivial ist.
 in thesis.tex,
% nach \chapter{Berechnungstheoretische Vorteile bestimmter Graphklassen zur Bestimmung bedingter Unabhängigkeit}
\label{ch:graphclasses}

Die folgenden Abschnitte orientieren sich an \cite{handbook__of_formal_argumentation__vol__3:2024}.

\section{Azyklische AFs}
\label{sec:acyc-afs}

Für azyklische AFs fällt jedes existierende \emph{stable}-Labelling mit dem eindeutig bestimmten \emph{grounded}-Labelling zusammen. 
Folglich besitzt ein azyklisches AF höchstens ein \emph{stable}-Labelling. 
Dadurch ist die bedingte Unabhängigkeit bezüglich der \emph{stable}-Semantik für beliebige paarweise disjunkte Argumentmengen 
stets erfüllt.

\begin{Proposition}
\label{prop:acyc-af-ci-st-top}
    Sei F=(A,R) ein azyklisches AF. Dann ist
    \(\textit{Ind}_{\mathrm{st}}(F) \) immer wahr.
\end{Proposition}

\begin{proof}
    \label{pr:acyc-af-ci-st-top-true}
    Seien \( X,Y,Z \subseteq A\) paarweise disjunkt. Besitzt \(F\) kein
    \emph{stable}-Labelling, so ist die universell quantifizierte
    Unabhängigkeitsbedingung vakuos erfüllt.

    Andernfalls existiert aufgrund der Azyklizität genau ein
    \emph{stable}-Labelling \(\lambda\), das mit dem
    \emph{grounded}-Labelling übereinstimmt. Für beliebige
    \(\lambda_1,\lambda_2\in\Lambda_{\mathrm{st}}(F)\) gilt daher
    \(\lambda_1=\lambda_2=\lambda\). Mit \(\lambda_3:=\lambda\) können
    somit die Belegungen von \(X\), \(Y\) und \(Z\) wie in der Definition
    der bedingten Unabhängigkeit gefordert kombiniert werden. Daher gilt
    \(X\bot_{\mathrm{st}}Y\mid Z\) und folglich
    \(\operatorname{Ind}_{\mathrm{st}}(F)\).
\end{proof}

\section{Geradezykelfreie AFs}
\label{sec:even-cycle-free-afs}

Für AFs ohne gerade Zykel ist allein das \emph{grounded}-Labelling ein Kandidat, ein \emph{stable}-Labelling zu sein. 
Damit existiert ebenfalls maximal ein \emph{stable}-Labelling. 

\begin{Proposition}
\label{prop:even-cycle-free-af-ci-st-top}
    Sei F=(A,R) ein AF ohne gerade Zyklen. Dann ist
    \(\textit{Ind}_{\mathrm{st}}(F) \) immer wahr.
\end{Proposition}

\begin{proof}
    Analog zu \ref{pr:acyc-af-ci-st-top-true}.
\end{proof}


\section{Ungeradezykelfreie AFs}
\label{sec:odd-cycle-free-afs}

Für die Berechnung von \emph{stable}-Labellings gibt es in dieser Graphklasse keine berechnungstheoretischen Vorteile. 


\section{Bipartite AFs}
\label{sec:bi-afs}

Als Spezialfall eines ungeradezykelfreien AFs gibt es für die Berechnung von \emph{stable}-Labellings in bipartiten AFs keine
berechnungstheoretischen Vorteile. 


\section{Symmetrische AFs}
\label{sec:sym-afs}

Beobachtung: Je mehr Angriffe im AF, desto häufiger nicht bedingt unabhängig.

.
% ****************************************************************************************************

\chapter{Analyse des Härtebeweises für bipartite AFs}
\label{ch:bip-analyse}

\section{Gegenstand und Vorgehen}
\label{sec:bip-analyse-gegenstand}

Dieses Kapitel untersucht Proposition~\ref{prop:bip-st-ind-hard} aus
Kapitel~\ref{ch:graphclasses}, also die Behauptung, dass
\textit{stable}-Unabhängigkeit auch auf bipartiten AFs
\(\Pi_2^\text{P}\)-vollständig ist. Die Untersuchung erfolgt in drei Schritten.
Zuerst wird die Konstruktion so rekonstruiert, dass ihre Semantik durch ein
explizites Strukturlemma erfasst wird (Abschnitt~\ref{sec:bip-analyse-struktur}).
Anschließend wird die Korrektheit vollständig nachgerechnet
(Abschnitt~\ref{sec:bip-analyse-korrektheit}) und die
Graphklassenzugehörigkeit verifiziert
(Abschnitt~\ref{sec:bip-analyse-bipartit}). Erst danach werden die gefundenen
Mängel gesammelt (Abschnitt~\ref{sec:bip-analyse-kritik}) und behoben
(Abschnitt~\ref{sec:bip-analyse-reparatur}).

Das Ergebnis vorweg: \emph{Die Reduktion ist inhaltlich korrekt.} Die
Beweisführung weist jedoch eine Reihe formaler Mängel auf, von denen einer
(die unbelegte Ausgangsprobleminstanz) den Beweis in seiner jetzigen Form
angreifbar macht, während die übrigen die Nachvollziehbarkeit betreffen.

\section{Rekonstruktion der Konstruktion}
\label{sec:bip-analyse-struktur}

Wir wiederholen die Konstruktion in fixierter Notation. Sei
\(\Psi=\forall X\exists Y\,\varphi\) mit
\(\varphi=\bigwedge_{C\in\mathcal{C}}C\) in konjunktiver Normalform, wobei jede
Klausel \emph{monoton} ist, also entweder ausschließlich positive Literale
(\(C\in\mathcal{C}^{+}\)) oder ausschließlich negative Literale
(\(C\in\mathcal{C}^{-}\)) enthält. Sei \(V:=X\cup Y\cup\{p,n\}\) mit zwei
frischen Variablen \(p,n\). Das AF \(F_\Psi=(A_\Psi,R_\Psi)\) ist gegeben durch
\[
\begin{aligned}
    A_{\Psi}
    ={}&
    \{v,\overline{v}\mid v\in V\}
    \cup
    \{c_C\mid C\in\mathcal{C}\}
    \cup
    \{\varphi^{+},\varphi^{-},t^{+},t^{-}\}
\end{aligned}
\]
und der in Proposition~\ref{prop:bip-st-ind-hard} angegebenen Angriffsmenge
\(R_\Psi\). Zentral ist die Beobachtung, dass die Argumente \(p\) und
\(\overline{n}\) genau so wirken, als käme das Literal \(p\) in jeder positiven
und das Literal \(\neg n\) in jeder negativen Klausel vor. Wir machen das
explizit und setzen
\[
    \widehat{\varphi}
    :=
    \bigwedge_{C\in\mathcal{C}^{+}}(p\lor C)
    \;\land\;
    \bigwedge_{C\in\mathcal{C}^{-}}(\neg n\lor C).
\]

\begin{Lemma}
    \label{lem:bip-qbf-aequiv}
    Es gilt
    \(
        \Psi\text{ ist gültig}
        \iff
        \forall (X\cup\{p,n\})\,\exists Y\,\widehat{\varphi}
        \text{ ist gültig.}
    \)
\end{Lemma}

\begin{proof}
    \enquote{\(\Rightarrow\)}: Seien \(\alpha\) eine Belegung von \(X\) und
    \(\pi,\nu\) beliebige Wahrheitswerte für \(p\) und \(n\). Da \(\Psi\)
    gültig ist, existiert \(\beta\) mit \(\varphi(\alpha,\beta)=\top\). Dann ist
    jede Klausel \(C\in\mathcal{C}\) unter \((\alpha,\beta)\) erfüllt, also erst
    recht \(p\lor C\) beziehungsweise \(\neg n\lor C\). Somit gilt
    \(\widehat{\varphi}=\top\).

    \enquote{\(\Leftarrow\)}: Sei \(\alpha\) eine Belegung von \(X\). Wir wählen
    die Belegung \(p=\bot\) und \(n=\top\). Dann gilt
    \(\widehat{\varphi}\equiv\varphi\), und die Voraussetzung liefert ein
    \(\beta\) mit \(\varphi(\alpha,\beta)=\top\).
\end{proof}

Der Beweis in Kapitel~\ref{ch:graphclasses} argumentiert an mehreren Stellen mit
Aussagen der Form \enquote{die Argumente \(t^{+}\) und \(t^{-}\) erlauben es,
die Labels von \(\varphi^{+}\) und \(\varphi^{-}\) unabhängig zu setzen}, ohne
diese zu beweisen. Das folgende Lemma stellt die vollständige Struktur von
\(\Lambda_{\mathrm{st}}(F_\Psi)\) bereit und macht alle späteren Schritte
elementar.

\begin{Lemma}[Strukturlemma]
    \label{lem:bip-struktur}
    Sei \(\lambda\) ein Labelling von \(F_\Psi\). Dann ist \(\lambda\) genau
    dann \textit{stable}, wenn es eine totale Belegung \(\tau:V\to\{\top,\bot\}\)
    gibt, sodass die folgenden vier Bedingungen gelten.
    \begin{enumerate}
        \item Für alle \(v\in V\) ist
        \(\lambda(v)=\mathit{in}\iff\tau(v)=\top\)
        und
        \(\lambda(\overline{v})=\mathit{in}\iff\tau(v)=\bot\).
        \item Für alle \(C\in\mathcal{C}^{+}\) ist
        \(\lambda(c_C)=\mathit{out}\) genau dann, wenn \(\tau\models p\lor C\);
        andernfalls \(\lambda(c_C)=\mathit{in}\). Analog ist für
        \(C\in\mathcal{C}^{-}\) genau dann \(\lambda(c_C)=\mathit{out}\), wenn
        \(\tau\models\neg n\lor C\).
        \item \(\lambda(t^{+})=\mathit{in}\iff\lambda(\varphi^{+})=\mathit{out}\)
        und \(\lambda(t^{-})=\mathit{in}\iff\lambda(\varphi^{-})=\mathit{out}\).
        \item \(\lambda(\varphi^{+})=\mathit{in}\) impliziert
        \(\lambda(c_C)=\mathit{out}\) für alle \(C\in\mathcal{C}^{+}\);
        \(\lambda(\varphi^{-})=\mathit{in}\) impliziert
        \(\lambda(c_C)=\mathit{out}\) für alle \(C\in\mathcal{C}^{-}\).
    \end{enumerate}
    Umgekehrt liefert jede Wahl von \(\tau\) und jede mit (3) und (4)
    verträgliche Wahl von \(\lambda(\varphi^{+}),\lambda(\varphi^{-})\)
    ein \textit{stable}-Labelling. Insbesondere ist
    \(\lambda(\varphi^{+})=\mathit{out}\) stets realisierbar, und
    \(\lambda(\varphi^{+})=\mathit{in}\) ist genau dann realisierbar, wenn
    \(\tau\) alle Klauseln \(p\lor C\) mit \(C\in\mathcal{C}^{+}\) erfüllt.
\end{Lemma}

\begin{proof}
    Sei \(\lambda\in\Lambda_{\mathrm{st}}(F_\Psi)\); insbesondere ist
    \(\mathit{und}(\lambda)=\emptyset\).

    Zu (1): Die einzigen Angreifer von \(v\) sind \(\overline{v}\) und
    umgekehrt. Wären beide \(\mathit{in}\), so griffe ein
    \(\mathit{in}\)-Argument ein \(\mathit{in}\)-Argument an. Wären beide
    \(\mathit{out}\), so besäße \(v\) keinen \(\mathit{in}\)-Angreifer. Also ist
    genau eines der beiden \(\mathit{in}\), und wir setzen
    \(\tau(v):=\top\) genau dann, wenn \(\lambda(v)=\mathit{in}\).

    Zu (2): Die Angreifer von \(c_C\) mit \(C\in\mathcal{C}^{+}\) sind genau
    die Argumente \(v\) mit \(v\in C\) sowie \(p\). Nach (1) ist ein solcher
    Angreifer genau dann \(\mathit{in}\), wenn das zugehörige Literal unter
    \(\tau\) wahr ist. Da \(\lambda\) \textit{stable} ist, gilt
    \(\lambda(c_C)=\mathit{out}\) genau dann, wenn mindestens ein Angreifer
    \(\mathit{in}\) ist, also genau dann, wenn \(\tau\models p\lor C\). Der Fall
    \(C\in\mathcal{C}^{-}\) verläuft analog mit den Angreifern \(\overline{v}\)
    für \(\neg v\in C\) und \(\overline{n}\).

    Zu (3): Der einzige Angreifer von \(t^{+}\) ist \(\varphi^{+}\). Also ist
    \(t^{+}\) genau dann \(\mathit{out}\), wenn \(\varphi^{+}\) \(\mathit{in}\)
    ist, und wegen \(\mathit{und}(\lambda)=\emptyset\) genau dann
    \(\mathit{in}\), wenn \(\varphi^{+}\) \(\mathit{out}\) ist.

    Zu (4): Die Angreifer von \(\varphi^{+}\) sind die \(c_C\) mit
    \(C\in\mathcal{C}^{+}\) sowie \(t^{+}\). Ist \(\varphi^{+}\)
    \(\mathit{in}\), so sind alle Angreifer \(\mathit{out}\).

    Für die Umkehrung sei \(\tau\) beliebig und seien
    \(\lambda(\varphi^{\pm})\) gemäß (3) und (4) gewählt; die übrigen Labels
    seien durch (1)--(3) festgelegt. Man verifiziert unmittelbar für jedes
    Argument die \textit{stable}-Bedingung: Argumente \(v,\overline{v}\) sind
    genau dann \(\mathit{out}\), wenn ihr Partner \(\mathit{in}\) ist;
    Klauselargumente genau dann, wenn ein Literalargument \(\mathit{in}\) ist;
    \(t^{\pm}\) genau dann, wenn \(\varphi^{\pm}\) \(\mathit{in}\) ist; und
    \(\varphi^{+}\) ist entweder \(\mathit{in}\) (dann sind nach (4) alle
    \(c_C\) und nach (3) auch \(t^{+}\) \(\mathit{out}\)) oder \(\mathit{out}\)
    (dann ist nach (3) \(t^{+}\) \(\mathit{in}\) und greift \(\varphi^{+}\) an).
    Es tritt kein \(\mathit{und}\) auf.

    Der Zusatz folgt: \(\lambda(\varphi^{+})=\mathit{out}\) ist wegen
    \(\lambda(t^{+})=\mathit{in}\) immer zulässig, und
    \(\lambda(\varphi^{+})=\mathit{in}\) ist nach (4) und (2) genau dann
    zulässig, wenn \(\tau\) jede Klausel \(p\lor C\), \(C\in\mathcal{C}^{+}\),
    erfüllt.
\end{proof}

Lemma~\ref{lem:bip-struktur} zeigt insbesondere: Die Labels von
\(\varphi^{+}\) und \(\varphi^{-}\) sind \emph{unabhängig voneinander} frei
wählbar, sobald die jeweiligen Klauselargumente \(\mathit{out}\) sind, da
\(\varphi^{+},t^{+}\) und \(\varphi^{-},t^{-}\) disjunkte Angreifermengen
besitzen. Genau diese Aussage benutzt der Originalbeweis, ohne sie zu begründen.

\section{Korrektheit der Reduktion}
\label{sec:bip-analyse-korrektheit}

\begin{Proposition}
    \label{prop:bip-korrekt}
    Es gilt
    \[
        \Psi\text{ ist gültig}
        \quad\Longleftrightarrow\quad
        \{\varphi^{+},\varphi^{-}\}\bot_{\mathrm{st}}(X\cup\{p,n\})\mid\emptyset
        \text{ in }F_\Psi.
    \]
\end{Proposition}

\begin{proof}
    Wir schreiben \(\mathcal{X}:=\{\varphi^{+},\varphi^{-}\}\) und
    \(\mathcal{Y}:=X\cup\{p,n\}\); beide Mengen sind disjunkt, und
    \(Z=\emptyset\), sodass die Bedingung \(\lambda_1|_Z=\lambda_2|_Z\) aus
    Definition~\ref{def:BU:AFs} für jedes Paar erfüllt ist. Gesucht ist also
    stets ein \(\lambda_3\in\Lambda_{\mathrm{st}}(F_\Psi)\) mit
    \(\lambda_3|_{\mathcal{X}}=\lambda_1|_{\mathcal{X}}\) und
    \(\lambda_3|_{\mathcal{Y}}=\lambda_2|_{\mathcal{Y}}\).

    \enquote{\(\Rightarrow\)}: Sei \(\Psi\) gültig und seien
    \(\lambda_1,\lambda_2\in\Lambda_{\mathrm{st}}(F_\Psi)\). Nach
    Lemma~\ref{lem:bip-struktur} induziert \(\lambda_2\) eine totale Belegung
    \(\tau_2\) von \(V\); insbesondere legt \(\lambda_2|_{\mathcal{Y}}\) die
    Werte \(\tau_2|_{X\cup\{p,n\}}\) fest. Nach Lemma~\ref{lem:bip-qbf-aequiv}
    existiert eine Belegung \(\beta\) von \(Y\), sodass
    \(\tau_3:=\tau_2|_{X\cup\{p,n\}}\cup\beta\) die Formel
    \(\widehat{\varphi}\) erfüllt. Wir wählen \(\lambda_3\) als das durch
    \(\tau_3\) und
    \(\lambda_3(\varphi^{+}):=\lambda_1(\varphi^{+})\),
    \(\lambda_3(\varphi^{-}):=\lambda_1(\varphi^{-})\)
    bestimmte Labelling. Da \(\tau_3\models\widehat{\varphi}\), sind nach
    Lemma~\ref{lem:bip-struktur}(2) sämtliche Klauselargumente
    \(\mathit{out}\); damit sind beide Wahlen für \(\varphi^{+}\) und
    \(\varphi^{-}\) mit (4) verträglich, und \(\lambda_3\) ist
    \textit{stable}. Nach Konstruktion gilt
    \(\lambda_3|_{\mathcal{X}}=\lambda_1|_{\mathcal{X}}\) und
    \(\lambda_3|_{\mathcal{Y}}=\lambda_2|_{\mathcal{Y}}\).

    \enquote{\(\Leftarrow\)}: Sei \(\Psi\) nicht gültig. Dann existiert eine
    Belegung \(\alpha\) von \(X\), sodass \(\varphi(\alpha,\beta)=\bot\) für
    jede Belegung \(\beta\) von \(Y\) gilt. Wir wählen:
    \begin{itemize}
        \item \(\lambda_1\) als ein beliebiges \textit{stable}-Labelling mit
        \(\lambda_1(\varphi^{+})=\lambda_1(\varphi^{-})=\mathit{in}\). Ein
        solches existiert: Man setze \(\tau_1(p)=\top\) und \(\tau_1(n)=\bot\)
        sowie \(\tau_1\) auf \(X\cup Y\) beliebig. Dann sind alle Klauseln
        \(p\lor C\) und \(\neg n\lor C\) erfüllt, also nach
        Lemma~\ref{lem:bip-struktur} alle Klauselargumente \(\mathit{out}\)
        und beide \(\varphi\)-Argumente \(\mathit{in}\) realisierbar.
        \item \(\lambda_2\) als ein \textit{stable}-Labelling zur Belegung
        \(\tau_2:=\alpha\cup\{p\mapsto\bot,\;n\mapsto\top\}\cup\beta_0\) mit
        beliebigem \(\beta_0\).
    \end{itemize}
    Angenommen, es gäbe ein \(\lambda_3\in\Lambda_{\mathrm{st}}(F_\Psi)\) mit
    \(\lambda_3|_{\mathcal{X}}=\lambda_1|_{\mathcal{X}}\) und
    \(\lambda_3|_{\mathcal{Y}}=\lambda_2|_{\mathcal{Y}}\). Sei \(\tau_3\) die
    zugehörige Belegung. Aus
    \(\lambda_3|_{\mathcal{Y}}=\lambda_2|_{\mathcal{Y}}\) folgt
    \(\tau_3|_X=\alpha\), \(\tau_3(p)=\bot\) und \(\tau_3(n)=\top\). Aus
    \(\lambda_3(\varphi^{+})=\lambda_3(\varphi^{-})=\mathit{in}\) folgt mit
    Lemma~\ref{lem:bip-struktur}(4) und (2), dass \(\tau_3\) sämtliche Klauseln
    \(p\lor C\) und \(\neg n\lor C\) erfüllt, also
    \(\tau_3\models\widehat{\varphi}\). Wegen \(\tau_3(p)=\bot\) und
    \(\tau_3(n)=\top\) gilt aber \(\widehat{\varphi}\equiv\varphi\) unter
    \(\tau_3\), also \(\varphi(\alpha,\tau_3|_Y)=\top\) im Widerspruch zur Wahl
    von \(\alpha\). Also existiert kein solches \(\lambda_3\), und es gilt
    \(\mathcal{X}\not\bot_{\mathrm{st}}\mathcal{Y}\mid\emptyset\).
\end{proof}

Die Konstruktion ist offensichtlich in Polynomialzeit berechenbar: Sie erzeugt
\(2|V|+|\mathcal{C}|+4\) Argumente und
\(\mathcal{O}(|V|+\sum_{C\in\mathcal{C}}|C|)\) Angriffe.

\section{Graphklassenzugehörigkeit}
\label{sec:bip-analyse-bipartit}

\begin{Proposition}
    \label{prop:bip-korrekt-bipartit}
    \(F_\Psi\) ist bipartit.
\end{Proposition}

\begin{proof}
    Wir setzen
    \[
    \begin{aligned}
        L:={}& \{v\mid v\in V\}\cup\{c_C\mid C\in\mathcal{C}^{-}\}
              \cup\{\varphi^{+},t^{-}\},\\
        R:={}& \{\overline{v}\mid v\in V\}\cup\{c_C\mid C\in\mathcal{C}^{+}\}
              \cup\{\varphi^{-},t^{+}\}.
    \end{aligned}
    \]
    Offensichtlich gilt \(L\cap R=\emptyset\) und \(L\cup R=A_\Psi\). Für jede
    Art von Angriff prüft man die Seitenwechsel-Eigenschaft:
    \((v,\overline{v})\) und \((\overline{v},v)\) verlaufen zwischen \(L\) und
    \(R\); \((v,c_C)\) mit \(C\in\mathcal{C}^{+}\) sowie \((p,c_C)\) verlaufen
    von \(L\) nach \(R\); \((\overline{v},c_C)\) mit \(C\in\mathcal{C}^{-}\)
    sowie \((\overline{n},c_C)\) verlaufen von \(R\) nach \(L\);
    \((c_C,\varphi^{+})\) mit \(C\in\mathcal{C}^{+}\) verläuft von \(R\) nach
    \(L\); \((c_C,\varphi^{-})\) mit \(C\in\mathcal{C}^{-}\) von \(L\) nach
    \(R\); \((\varphi^{+},t^{+})\) von \(L\) nach \(R\) und
    \((\varphi^{-},t^{-})\) von \(R\) nach \(L\), jeweils mit den
    Gegenangriffen. Also gilt
    \(R_\Psi\subseteq(L\times R)\cup(R\times L)\).
\end{proof}

Hier zeigt sich der eigentliche Grund für die Monotonieforderung: Nur weil
positive Klauseln ausschließlich von unquerstrichenen und negative Klauseln
ausschließlich von querstrichenen Literalargumenten angegriffen werden, lassen
sich die Klauselargumente überhaupt widerspruchsfrei auf die beiden Seiten
verteilen. Bei gemischten Klauseln entstünde ein Klauselargument mit Angreifern
aus \(L\) \emph{und} \(R\), und die Bipartition bräche zusammen.

\section{Kritikpunkte}
\label{sec:bip-analyse-kritik}

Die Reduktion ist nach den Abschnitten~\ref{sec:bip-analyse-korrektheit}
und~\ref{sec:bip-analyse-bipartit} korrekt. Der Beweis in
Kapitel~\ref{ch:graphclasses} weist jedoch die folgenden Mängel auf. Wir
ordnen sie nach abnehmender Schwere.

\begin{enumerate}
    \item \textbf{Unbelegtes Ausgangsproblem.} Der Beweis reduziert von
    \textsc{Balanced Monotone \(\forall\exists\)-3-SAT-\((1,1,2,2)\)} und
    behauptet ohne Quellenangabe und ohne Definition, dieses Problem sei
    \(\Pi_2^\text{P}\)-vollständig. Weder wird gesagt, worauf sich
    \enquote{balanced} bezieht, noch was das Tupel \((1,1,2,2)\) zählt. Das ist
    der einzige Mangel, der die \emph{Gültigkeit} der Aussage betrifft: Eine
    Reduktion von einem Problem unbekannter Komplexität beweist nichts.
    Vorkommensbeschränkungen sind zudem heikel. Nach dem Satz von Tovey ist
    jede \(k\)-CNF-Formel, in der jede Klausel \(k\) verschiedene Literale
    besitzt und jede Variable in höchstens \(k\) Klauseln vorkommt, erfüllbar;
    zu enge Vorkommensschranken können ein Problem also trivialisieren.
    \item \textbf{Die Zusatzrestriktionen werden nicht benötigt.} Die
    Konstruktion verwendet ausschließlich die Monotonie der Klauseln (siehe
    Abschnitt~\ref{sec:bip-analyse-bipartit}); weder die Ausgeglichenheit von
    \(|\mathcal{C}^{+}|\) und \(|\mathcal{C}^{-}|\) noch irgendeine
    Vorkommensschranke noch die Klauselbreite \(3\) geht in Korrektheit oder
    Bipartitheit ein. Der Beweis macht sich also ohne Not von einer stärkeren
    und schwerer zu belegenden Voraussetzung abhängig.
    Abschnitt~\ref{sec:bip-analyse-reparatur} zeigt, dass die Monotonie
    elementar und selbsttragend erreicht werden kann.
    \item \textbf{Vertauschte Rollen von \(\lambda_1\) und \(\lambda_2\).} In
    der Richtung \enquote{\(\Psi\) gültig} übernimmt \(\lambda_1\) die Labels
    von \(\varphi^{+},\varphi^{-}\) und \(\lambda_2\) die Belegung von
    \(X\cup\{p,n\}\), passend zur Reihenfolge der Mengen in
    \(\{\varphi^{+},\varphi^{-}\}\bot_{\mathrm{st}}X\cup\{p,n\}\mid\emptyset\)
    und zu Definition~\ref{def:BU:AFs}. In der Gegenrichtung sind die Rollen
    stillschweigend vertauscht: Dort trägt \(\lambda_1\) die Belegung und
    \(\lambda_2\) die \(\varphi\)-Labels. Inhaltlich ist das folgenlos, weil
    die bedingte Unabhängigkeit symmetrisch ist, formal ist es jedoch ein
    Bruch mit der eigenen Definition. Entweder man benennt konsistent um (wie
    in Proposition~\ref{prop:bip-korrekt}) oder man beruft sich explizit auf
    das Symmetrieaxiom.
    \item \textbf{Undefinierte Mengen \(A\) und \(B\).} Der Satz
    \enquote{[\ldots] das die Einschränkung von \(\lambda_1\) auf \(A\) und die
    Einschränkung von \(\lambda_2\) auf \(B\) kombiniert} verwendet zwei nie
    eingeführte Symbole. Erschwerend kommt hinzu, dass \(A\) in der gesamten
    Arbeit die Argumentmenge eines AFs bezeichnet; gemeint sind offenbar
    \(\{\varphi^{+},\varphi^{-}\}\) und \(X\cup\{p,n\}\).
    \item \textbf{Namenskollision \(L,R\).} Die Bipartitionsklassen heißen
    \(L\) und \(R\), während \(R\) beziehungsweise \(R_\Psi\) die
    Angriffsrelation bezeichnet. Die Schlussformel
    \(R_\Psi\subseteq(L\times R)\cup(R\times L)\) mischt beide Bedeutungen in
    einer Zeile. Vorzuziehen wären etwa \(A_1,A_2\) oder \(L_1,L_2\).
    \item \textbf{Fehlendes Strukturlemma.} Der Beweis benutzt mehrfach
    Eigenschaften von \(\Lambda_{\mathrm{st}}(F_\Psi)\), die er nicht beweist:
    dass in jedem \textit{stable}-Labelling genau eines von \(v,\overline{v}\)
    das Label \(\mathit{in}\) trägt; dass \(c_C\) genau dann \(\mathit{out}\)
    ist, wenn ein Literalargument \(\mathit{in}\) ist; und dass \(t^{+},t^{-}\)
    die freie und \emph{voneinander unabhängige} Wahl der \(\varphi\)-Labels
    erlauben. Lemma~\ref{lem:bip-struktur} schließt diese Lücke.
    \item \textbf{Unvollständige Fallunterscheidung über \(p\) und \(n\).} Der
    Abschnitt \enquote{Konstruktion} diskutiert nur die Kombinationen
    \((p,n)=(\mathit{out},\mathit{in})\) und
    \((p,n)=(\mathit{in},\mathit{out})\). Die ebenfalls in
    \textit{stable}-Labellings auftretenden Kombinationen
    \((\mathit{in},\mathit{in})\) und \((\mathit{out},\mathit{out})\) bleiben
    unerwähnt. Der Korrektheitsbeweis benötigt diese Fallunterscheidung
    tatsächlich nicht --- er läuft über die Formeläquivalenz aus
    Lemma~\ref{lem:bip-qbf-aequiv} ---, der Text suggeriert jedoch eine
    erschöpfende Analyse, die er nicht liefert.
    \item \textbf{Zugehörigkeit fehlt.} Die Proposition behauptet
    \(\Pi_2^\text{P}\)-\emph{Vollständigkeit}, der Beweis zeigt nur die Härte.
    Ein Verweis auf Theorem~\ref{th:ind-st-is-pi2p-c} (die obere Schranke gilt
    für alle AFs, also insbesondere für bipartite) fehlt, anders als beim
    ungeradezykelfreien Fall.
    \item \textbf{Formulierungs- und Notationsmängel.} Die Instanz wird als
    Behauptung formuliert (\enquote{wählen wir
    \(\{\varphi^{+},\varphi^{-}\}\perp_{st}\ldots\)}) statt als Anfrage; die
    Polynomialität der Konstruktion wird nicht erwähnt; \(\perp\) wird sowohl
    für die Unabhängigkeitsrelation als auch (in \enquote{\(p=\perp\)}) für den
    Wahrheitswert \emph{falsch} verwendet, während an anderer Stelle
    \(\bot_{\mathrm{st}}\) für die Unabhängigkeit steht; die Argumentmenge
    heißt hier \(\mathcal{A}_\Psi\), im ungeradezykelfreien Fall dagegen
    \(A_\Psi\); und der Satz \enquote{dann werden die Klauselargumente
    ausschließlich durch die ursprünglichen Literale kontrolliert werden}
    enthält ein doppeltes Verb.
\end{enumerate}

\begin{Corollary}
    \label{cor:bip-subsumiert-ocf}
    Proposition~\ref{prop:bip-st-ind-hard} impliziert
    Proposition~\ref{prop:ocf-af-ci-st-pi^2}.
\end{Corollary}

\begin{proof}
    Ist \(F\) bipartit, so wechselt jeder gerichtete Zyklus in jedem Schritt die
    Bipartitionsklasse und hat daher gerade Länge. Bipartite AFs sind also
    ungeradezykelfrei, und die Härte auf der Teilklasse überträgt sich auf die
    Oberklasse.
\end{proof}

Das ist ein struktureller Hinweis, kein Fehler: Der eigenständige Beweis für
ungeradezykelfreie AFs für \(\sigma=\mathrm{st}\) und \(\sigma=\mathrm{pr}\) ist
entbehrlich, sobald das bipartite Resultat gesichert ist. Umgekehrt bleibt der
gesonderte Beweis für \(\sigma=\mathrm{co}\)
(Proposition~\ref{prop:ocf-af-ci-co-pi^2-c}) nötig, denn das dort verwendete AF
ist wegen der Argumente \(d_i\), die sowohl von \(x_i\) als auch von
\(\overline{x_i}\) angegriffen werden, nicht bipartit.

Schließlich sei angemerkt, dass die Herleitung von
Proposition~\ref{prop:bip-prf-ind-hard} aus dem \textit{stable}-Fall die
Kohärenz ungeradezykelfreier AFs voraussetzt, also
\(\Lambda_{\mathrm{st}}(F)=\Lambda_{\mathrm{pr}}(F)\). Diese Tatsache stammt von
\textcite{dung:1995:af} und sollte dort zitiert werden.

\section{Reparatur: Monotonie ohne fremde Voraussetzung}
\label{sec:bip-analyse-reparatur}

Der schwerwiegendste Kritikpunkt lässt sich vollständig beheben. Statt sich auf
eine exotische Restriktion zu berufen, genügt eine elementare Normalform, die
sich in wenigen Zeilen beweisen lässt.

\begin{Lemma}[Monotone Normalform]
    \label{lem:monotone-normalform}
    Zu jeder Formel \(\Psi=\forall X\exists Y\,\varphi\) mit \(\varphi\) in KNF
    lässt sich in Polynomialzeit eine Formel
    \(\Psi'=\forall X\exists Y'\,\varphi'\) berechnen, sodass jede Klausel von
    \(\varphi'\) entweder nur positive oder nur negative Literale enthält und
    \(\Psi\) genau dann gültig ist, wenn \(\Psi'\) gültig ist.
\end{Lemma}

\begin{proof}
    Für jede Variable \(v\in X\cup Y\) sei \(v'\) eine frische Variable und
    \(Y':=Y\cup\{v'\mid v\in X\cup Y\}\). Die Formel \(\varphi'\) entstehe aus
    \(\varphi\), indem jedes negative Literal \(\neg v\) durch \(v'\) ersetzt
    wird; zusätzlich nehmen wir für jedes \(v\in X\cup Y\) die beiden Klauseln
    \(
        (v\lor v')
    \) und \(
        (\neg v\lor\neg v')
    \)
    auf. Alle aus \(\varphi\) entstandenen Klauseln sind rein positiv, die
    ersten Zusatzklauseln ebenfalls, die zweiten rein negativ. Also ist
    \(\varphi'\) monoton. Die beiden Zusatzklauseln sind zusammen äquivalent zu
    \(v'\leftrightarrow\neg v\).

    Sei \(\Psi\) gültig und \(\alpha\) eine Belegung von \(X\). Wähle \(\beta\)
    mit \(\varphi(\alpha,\beta)=\top\) und setze \(v'\) auf den zu \(v\)
    komplementären Wert. Die Zusatzklauseln sind dann erfüllt. Ist eine
    ursprüngliche Klausel \(C\) durch ein positives Literal \(v\) erfüllt, so
    auch ihr Bild; ist sie durch \(\neg v\) erfüllt, so ist \(v\) falsch, also
    \(v'\) wahr, und das Bild ist ebenfalls erfüllt. Somit ist \(\Psi'\) gültig.

    Sei umgekehrt \(\Psi'\) gültig und \(\alpha\) eine Belegung von \(X\). Es
    existiert eine Belegung von \(Y'\), die \(\varphi'\) erfüllt; wegen der
    Zusatzklauseln gilt dort \(v'=\neg v\) für alle \(v\in X\cup Y\). Sei
    \(\beta\) die Einschränkung auf \(Y\). Jede ursprüngliche Klausel \(C\) hat
    ein erfülltes Bildliteral; ist dieses ein \(v\), so erfüllt \(v\) auch
    \(C\); ist es ein \(v'\), so ist \(v\) falsch und \(\neg v\in C\) erfüllt.
    Also gilt \(\varphi(\alpha,\beta)=\top\) und \(\Psi\) ist gültig.

    Die Konstruktion verdoppelt die Variablenzahl und fügt \(2|X\cup Y|\)
    Klauseln hinzu, ist also polynomiell.
\end{proof}

Damit ist \textsc{Monotone \(\forall\exists\)-SAT} \(\Pi_2^\text{P}\)-hart,
sobald \(\forall\exists\)-SAT es ist, und
Proposition~\ref{prop:bip-st-ind-hard} steht auf einer belegten Grundlage. Wird
zusätzlich Klauselbreite \(3\) gewünscht, füllt man die zweielementigen
Zusatzklauseln durch Wiederholung eines Literals auf; die Monotonie bleibt
erhalten. Alle übrigen Kritikpunkte sind redaktioneller Natur und durch die
Fassung in den Abschnitten~\ref{sec:bip-analyse-struktur}
bis~\ref{sec:bip-analyse-bipartit} bereits ausgeräumt.

\section{Fazit}
\label{sec:bip-analyse-fazit}

Die Aussage von Proposition~\ref{prop:bip-st-ind-hard} ist richtig, und der
Kern der Konstruktion --- die Aufspaltung des Formelarguments in \(\varphi^{+}\)
und \(\varphi^{-}\) sowie die Steuerung durch die frischen Variablen \(p\) und
\(n\) --- ist elegant und trägt die Bipartitheit. Der Beweis sollte jedoch um
Lemma~\ref{lem:bip-qbf-aequiv}, Lemma~\ref{lem:bip-struktur} und
Lemma~\ref{lem:monotone-normalform} ergänzt, in den Rollen von
\(\lambda_1,\lambda_2\) vereinheitlicht und um den Zugehörigkeitsverweis
ergänzt werden. Danach ist er lückenlos.

\chapter{Symmetrische irreflexive AFs}
\label{ch:sym-analyse}

\section{Fragestellung}
\label{sec:sym-frage}

Abschnitt~\ref{sec:sym-afs} hat symmetrische AFs bereits als Extremfall hoher
Angriffsdichte betrachtet und für vollständig gerichtete AFs gezeigt, dass
\textit{stable}-Unabhängigkeit dort stets verletzt ist
(Proposition~\ref{prop:cdg-afs-not-st-ci}). Offen blieb die
komplexitätstheoretische Einordnung der gesamten Klasse. Dieses Kapitel
schließt diese Lücke.

Der Ausgangsverdacht ist, dass die Klasse einfach sein müsste. Symmetrische
irreflexive AFs sind kohärent \cite{Coste:2005:symmetric-af-irflx-coherent},
besitzen stets \textit{stable}-Labellings und ihre klassischen
Entscheidungsprobleme --- Existenz, Verifikation, kredulöse und skeptische
Akzeptanz --- sind sämtlich in Polynomialzeit lösbar
(Proposition~\ref{prop:sym-klassisch-p}). Für alle bisher betrachteten
Graphklassen (Abschnitt~\ref{sec:acyc-afs} bis
Abschnitt~\ref{sec:odd-cycle-free-afs}) war eine solche strukturelle
Vereinfachung stets mit einer Vereinfachung der Unabhängigkeit einhergegangen
oder zumindest verträglich.

Das Hauptresultat dieses Kapitels widerlegt diesen Verdacht für die
\textit{stable}-Semantik: \(\textit{Ind}_{\mathrm{st}}\) bleibt auf
symmetrischen irreflexiven AFs \(\Pi_2^\text{P}\)-vollständig
(Theorem~\ref{th:sym-st-pi2p}). Für die \textit{preferred}-Semantik dagegen
tritt eine echte Vereinfachung ein: Das Problem fällt von
\(\Pi_3^\text{P}\)-Vollständigkeit auf \(\Pi_2^\text{P}\)-Vollständigkeit
(Korollar~\ref{cor:sym-pr-pi2p}). Ergänzend wird ein Fragment identifiziert, in
dem die Unabhängigkeit tatsächlich in Polynomialzeit entscheidbar wird
(Proposition~\ref{prop:sym-singleton-p}).

Wir betrachten für \(\sigma\in\{co,st,pr\}\) das Entscheidungsproblem
\[
\begin{array}{ll}
\textit{Ind}^{\mathrm{sym}}_{\sigma}\\
\text{Input:}  & \text{symmetrisches, irreflexives AF } F=(A,R),\
                 \text{disjunkte } X,Y,Z\subseteq A \\
\text{Output:} & \text{TRUE} \Longleftrightarrow X \bot_{\sigma} Y \mid Z.
\end{array}
\]

\section{Struktur der Labellingmengen}
\label{sec:sym-struktur}

\begin{Definition}
    \label{def:sym-af}
    Ein AF \(F=(A,R)\) heißt \emph{symmetrisch}, falls
    \((a,b)\in R\Rightarrow(b,a)\in R\) für alle \(a,b\in A\) gilt, und
    \emph{irreflexiv}, falls \((a,a)\notin R\) für alle \(a\in A\) gilt. Für
    \(a\in A\) schreiben wir \(N(a):=\{b\in A\mid (a,b)\in R\}\) für die
    \emph{Nachbarschaft} von \(a\) und \(N[a]:=N(a)\cup\{a\}\). Für
    \(S\subseteq A\) sei \(N(S):=\bigcup_{a\in S}N(a)\).
\end{Definition}

In einem symmetrischen irreflexiven AF ist \(R\) die Kantenrelation eines
schlichten ungerichteten Graphen. Eine Menge \(S\subseteq A\) ist genau dann
konfliktfrei, wenn sie in diesem Graphen unabhängig ist, und wir nennen \(S\)
\emph{naiv}, falls \(S\) \(\subseteq\)-maximal konfliktfrei ist. Für
\(S\subseteq A\) bezeichne \(\lambda_S\) das Labelling mit
\(\mathit{in}(\lambda_S)=S\) und \(\mathit{out}(\lambda_S)=A\setminus S\).

\begin{Lemma}
    \label{lem:sym-cf-adm}
    Sei \(F=(A,R)\) symmetrisch und irreflexiv. Dann ist jede konfliktfreie
    Menge \(S\subseteq A\) admissibel.
\end{Lemma}

\begin{proof}
    Sei \(a\in S\) und sei \(b\) ein Angreifer von \(a\), also \((b,a)\in R\).
    Wegen der Konfliktfreiheit von \(S\) gilt \(b\notin S\). Aus der Symmetrie
    folgt \((a,b)\in R\), also greift \(a\in S\) den Angreifer \(b\) an. Damit
    verteidigt \(S\) jedes seiner Elemente.
\end{proof}

\begin{Proposition}
    \label{prop:sym-naiv-st-pr}
    Sei \(F=(A,R)\) symmetrisch und irreflexiv. Dann gilt
    \[
        \Lambda_{\mathrm{st}}(F)
        =\Lambda_{\mathrm{pr}}(F)
        =\{\lambda_S\mid S\subseteq A\text{ naiv}\}.
    \]
    Insbesondere ist \(\Lambda_{\mathrm{st}}(F)\neq\emptyset\).
\end{Proposition}

\begin{proof}
    Sei \(S\) naiv und \(a\in A\setminus S\). Dann ist \(S\cup\{a\}\) nicht
    konfliktfrei. Da \(S\) konfliktfrei und \(F\) irreflexiv ist, muss der
    Konflikt zwischen \(a\) und einem \(b\in S\) bestehen; wegen der Symmetrie
    gilt \((b,a)\in R\). Also greift \(S\) jedes Argument außerhalb von \(S\)
    an, das heißt \(\lambda_S\) ist ein \textit{stable}-Labelling.

    Umgekehrt ist jede \textit{stable}-Extension konfliktfrei und maximal
    konfliktfrei, denn ein hinzunehmbares Argument wäre unangegriffen. Also
    stimmen die \textit{stable}-Extensionen genau mit den naiven Mengen
    überein.

    Für die \textit{preferred}-Semantik gilt allgemein
    \(\Lambda_{\mathrm{st}}(F)\subseteq\Lambda_{\mathrm{pr}}(F)\). Sei
    umgekehrt \(E\) eine \textit{preferred}-Extension, also
    \(\subseteq\)-maximal admissibel. Wäre \(E\) nicht maximal konfliktfrei, so
    gäbe es eine konfliktfreie Menge \(E\subsetneq T\); nach
    Lemma~\ref{lem:sym-cf-adm} wäre \(T\) admissibel, im Widerspruch zur
    Maximalität von \(E\). Also ist \(E\) naiv und damit \textit{stable}. Da
    eine \textit{stable}-Extension \(E\) die Labellingdarstellung
    \(\mathit{in}=E\), \(\mathit{out}=A\setminus E\), \(\mathit{und}=\emptyset\)
    besitzt und dies zugleich das zu \(E\) gehörige
    \textit{preferred}-Labelling ist, stimmen auch die Labellingmengen überein.

    Da \(A\) endlich ist, besitzt jede konfliktfreie Menge eine maximale
    Obermenge; insbesondere existiert mindestens eine naive Menge.
\end{proof}

\begin{Corollary}
    \label{cor:sym-extend}
    Sei \(F=(A,R)\) symmetrisch und irreflexiv und sei \(T\subseteq A\)
    konfliktfrei. Dann existiert ein \(\lambda\in\Lambda_{\mathrm{st}}(F)\) mit
    \(T\subseteq\mathit{in}(\lambda)\). Ist zusätzlich \(a\in A\setminus T\) und
    besitzt \(a\) einen Nachbarn in \(T\), so gilt
    \(\lambda(a)=\mathit{out}\) für jedes solche \(\lambda\).
\end{Corollary}

\begin{proof}
    Man erweitere \(T\) zu einer naiven Menge \(S\) und wende
    Proposition~\ref{prop:sym-naiv-st-pr} an. Besitzt \(a\) einen Nachbarn in
    \(T\subseteq S\), so verhindert die Konfliktfreiheit \(a\in S\).
\end{proof}

Korollar~\ref{cor:sym-extend} ist das Arbeitspferd aller folgenden Beweise: Es
erlaubt, ein Labelling durch Angabe eines konfliktfreien \enquote{Kerns} zu
spezifizieren, sofern jedes Argument, das ausgeschlossen werden soll, einen
Nachbarn im Kern besitzt.

\begin{Proposition}
    \label{prop:sym-klassisch-p}
    Sei \(F=(A,R)\) symmetrisch und irreflexiv und sei \(\sigma\in\{st,pr\}\).
    Dann gilt:
    \begin{enumerate}
        \item \(\Lambda_\sigma(F)\neq\emptyset\).
        \item Die Verifikation, ob ein gegebenes Labelling in
        \(\Lambda_\sigma(F)\) liegt, ist in Polynomialzeit möglich.
        \item Jedes \(a\in A\) ist kredulös akzeptiert.
        \item Ein Argument \(a\in A\) ist genau dann skeptisch akzeptiert, wenn
        \(N(a)=\emptyset\).
    \end{enumerate}
\end{Proposition}

\begin{proof}
    (1) und (2) folgen aus Proposition~\ref{prop:sym-naiv-st-pr}: Zu prüfen ist
    lediglich, ob \(\mathit{in}(\lambda)\) konfliktfrei ist und ob jedes
    Argument außerhalb einen Nachbarn darin besitzt; beides ist in
    \(\mathcal{O}(|A|+|R|)\) möglich.

    (3) Wegen der Irreflexivität ist \(\{a\}\) konfliktfrei; nach
    Korollar~\ref{cor:sym-extend} existiert ein \(\lambda\in\Lambda_\sigma(F)\)
    mit \(\lambda(a)=\mathit{in}\).

    (4) Ist \(N(a)=\emptyset\), so ist \(a\) in jeder maximal konfliktfreien
    Menge enthalten, denn \(a\) steht mit keinem Argument in Konflikt. Ist
    dagegen \(b\in N(a)\), so liefert Korollar~\ref{cor:sym-extend} ein
    \(\lambda\) mit \(\lambda(b)=\mathit{in}\) und folglich
    \(\lambda(a)=\mathit{out}\).
\end{proof}

Proposition~\ref{prop:sym-klassisch-p} formalisiert den Befund von
\textcite{Coste:2005:symmetric-af-irflx-coherent}, dass in symmetrischen
irreflexiven AFs nur zwei sehr einfache und effizient prüfbare Formen der
Akzeptanz auftreten. Nach diesem Bild müsste auch die bedingte Unabhängigkeit
leicht sein: Sie ist nach Definition~\ref{def:BU:AFs} eine Aussage über
\(\Lambda_\sigma(F)\), und alle punktweisen Fragen über \(\Lambda_\sigma(F)\)
sind trivial.

\section{Die eigentliche Schwierigkeit: exakte Realisierbarkeit}
\label{sec:sym-realisierbarkeit}

Die bedingte Unabhängigkeit fragt nicht nach dem Label eines einzelnen
Arguments, sondern danach, ob ein \emph{vollständig vorgeschriebenes
Teillabelling} realisiert werden kann. Genau hier bricht die Einfachheit
zusammen.

\begin{Definition}
    \label{def:sym-real}
    Das Realisierbarkeitsproblem \(\textit{Real}^{\mathrm{sym}}_{\mathrm{st}}\)
    ist wie folgt definiert:
    \[
    \begin{array}{ll}
    \text{Input:}  & \text{symmetrisches, irreflexives AF } F=(A,R),\
                     W\subseteq A,\ \eta:W\to\{\mathit{in},\mathit{out}\}\\
    \text{Output:} & \text{TRUE} \Longleftrightarrow
                     \exists\lambda\in\Lambda_{\mathrm{st}}(F):\lambda|_W=\eta.
    \end{array}
    \]
\end{Definition}

Der Unterschied zur kredulösen Akzeptanz ist der Zwang, Argumente auch
\emph{ausschließen} zu müssen: Ein Argument \(a\) mit \(\eta(a)=\mathit{out}\)
verlangt einen Nachbarn im \(\mathit{in}\)-Teil, und mehrere solche Forderungen
müssen gleichzeitig durch eine einzige konfliktfreie Menge erfüllt werden.

\begin{Proposition}
    \label{prop:sym-real-npc}
    \(\textit{Real}^{\mathrm{sym}}_{\mathrm{st}}\) ist NP-vollständig.
\end{Proposition}

\begin{proof}
    Zugehörigkeit: Man rät \(\lambda\) und verifiziert in Polynomialzeit
    (Proposition~\ref{prop:sym-klassisch-p}(2)), dass
    \(\lambda\in\Lambda_{\mathrm{st}}(F)\) und \(\lambda|_W=\eta\).

    Härte: Wir reduzieren von 3-SAT. Sei \(\varphi=\bigwedge_{k=1}^{m}C_k\) eine
    KNF-Formel über den Variablen \(y_1,\dots,y_l\). Wir definieren ein
    symmetrisches irreflexives AF \(F_\varphi=(A,R)\) durch
    \[
        A:=\{w_j,\overline{w_j}\mid 1\leq j\leq l\}\cup\{c_k\mid 1\leq k\leq m\}
    \]
    und die (symmetrisch abgeschlossene) Angriffsrelation, die genau die
    folgenden Konflikte enthält:
    \begin{itemize}
        \item \(\{w_j,\overline{w_j}\}\) für jedes \(j\),
        \item \(\{w_j,c_k\}\), falls \(y_j\) positiv in \(C_k\) vorkommt,
        \item \(\{\overline{w_j},c_k\}\), falls \(y_j\) negativ in \(C_k\)
        vorkommt.
    \end{itemize}
    Wir setzen \(W:=\{c_1,\dots,c_m\}\) und \(\eta\equiv\mathit{out}\).

    Sei \(\lambda\in\Lambda_{\mathrm{st}}(F_\varphi)\) mit \(\lambda|_W=\eta\)
    und \(S:=\mathit{in}(\lambda)\). Da \(S\) konfliktfrei ist, gilt nicht
    gleichzeitig \(w_j\in S\) und \(\overline{w_j}\in S\); wir definieren
    \(\beta(y_j):=\top\) genau dann, wenn \(w_j\in S\). Nach
    Proposition~\ref{prop:sym-naiv-st-pr} ist \(S\) naiv, also besitzt jedes
    \(c_k\notin S\) einen Nachbarn in \(S\). Die Nachbarn von \(c_k\) sind
    ausschließlich Literalargumente, also enthält \(S\) ein Literalargument von
    \(C_k\). Ist dies \(w_j\), so gilt \(\beta(y_j)=\top\) und \(y_j\in C_k\);
    ist es \(\overline{w_j}\), so ist \(w_j\notin S\), also
    \(\beta(y_j)=\bot\) und \(\neg y_j\in C_k\). In beiden Fällen ist \(C_k\)
    unter \(\beta\) erfüllt. Somit ist \(\varphi\) erfüllbar.

    Sei umgekehrt \(\beta\) eine erfüllende Belegung und
    \(T:=\{w_j\mid\beta(y_j)=\top\}\cup\{\overline{w_j}\mid\beta(y_j)=\bot\}\).
    \(T\) ist konfliktfrei, denn \(T\) enthält aus jedem Paar genau ein
    Argument und die Paare stehen untereinander nicht in Konflikt. Jedes
    \(c_k\) besitzt ein unter \(\beta\) wahres Literal, dessen Literalargument
    in \(T\) liegt. Nach Korollar~\ref{cor:sym-extend} existiert ein
    \(\lambda\in\Lambda_{\mathrm{st}}(F_\varphi)\) mit
    \(T\subseteq\mathit{in}(\lambda)\), und alle \(c_k\) tragen dort
    \(\mathit{out}\). Also gilt \(\lambda|_W=\eta\).

    Die Konstruktion ist offensichtlich polynomiell.
\end{proof}

Proposition~\ref{prop:sym-real-npc} erklärt, warum die Tractability der
klassischen Aufgaben nicht auf die Unabhängigkeit durchschlägt: Der innere
Existenzquantor aus Definition~\ref{def:BU:AFs} ist bereits auf dieser Klasse
NP-vollständig. Was noch fehlt, ist ein Mechanismus, der den äußeren
Allquantor über \(\lambda_1,\lambda_2\) in einen echten \(\forall\)-Quantor
über Belegungen verwandelt. Dieser Mechanismus ist der Gegenstand des nächsten
Abschnitts.

\section{Hauptresultat}
\label{sec:sym-hauptresultat}

Die zentrale konstruktive Schwierigkeit in symmetrischen AFs ist, dass jeder
Angriff in beide Richtungen wirkt. Zwei Konsequenzen sind dabei entscheidend.

Erstens lässt sich die für Reduktionen übliche Invariante \enquote{aus jedem
Paar \(v,\overline{v}\) ist genau eines \(\mathit{in}\)} nicht mehr durch einen
bloßen Zweierzyklus erzwingen: Sobald \(v\) oder \(\overline{v}\) weitere
Nachbarn besitzen, können in einer naiven Menge beide fehlen, weil sie von
anderer Seite dominiert werden. Wir lösen das, indem wir Wert- und Signalknoten
trennen und zu einem Viererkreis verbinden.

Zweitens lässt sich die in gerichteten Konstruktionen übliche Aggregation
\enquote{\(\varphi\) ist \(\mathit{in}\), also sind alle Klauselargumente
\(\mathit{out}\), also ist jede Klausel erfüllt} nicht nachbauen: Ein
Aggregationsargument \(\varphi\), das alle Klauselargumente angreift,
dominiert diese selbst und macht damit jede Forderung an die Literale
gegenstandslos. Wir verzichten deshalb auf ein Aggregationsargument und
schreiben die Klauselargumente stattdessen \emph{direkt} als
\(\mathit{out}\)-Teil der Unabhängigkeitsanfrage vor. Damit die Menge der so
vorschreibbaren Muster klein und beherrschbar bleibt, verbinden wir die
Klauselargumente zu einer Clique.

\begin{Definition}
    \label{def:sym-konstruktion}
    Sei \(\Psi=\forall X\exists Y\,\varphi\) mit
    \(X=\{x_1,\dots,x_n\}\), \(Y=\{y_1,\dots,y_l\}\) und
    \(\varphi=\bigwedge_{k=1}^{m}C_k\) in KNF, \(m\geq 1\). Wir setzen
    \(X^{+}:=X\cup\{x_0\}\) mit einer frischen universellen Variablen \(x_0\)
    und
    \[
        \widehat{\varphi}:=\bigwedge_{k=1}^{m}\bigl(x_0\lor C_k\bigr).
    \]
    Das AF \(G_\Psi=(A_\Psi,R_\Psi)\) besitzt die Argumentmenge
    \[
    \begin{aligned}
        A_\Psi:={}&
        \{u_i,\overline{u_i},s_i,\overline{s_i}\mid x_i\in X^{+}\}
        \;\cup\;
        \{w_j,\overline{w_j}\mid y_j\in Y\}
        \;\cup\;
        \{c_1,\dots,c_m\},
    \end{aligned}
    \]
    und \(R_\Psi\) ist die symmetrische Hülle genau der folgenden Konflikte:
    \begin{enumerate}
        \item[(K1)] \(\{u_i,\overline{u_i}\}\),
        \(\{u_i,\overline{s_i}\}\),
        \(\{\overline{u_i},s_i\}\),
        \(\{s_i,\overline{s_i}\}\)
        \quad für alle \(x_i\in X^{+}\),
        \item[(K2)] \(\{w_j,\overline{w_j}\}\) \quad für alle \(y_j\in Y\),
        \item[(K3)] \(\{s_i,c_k\}\), falls \(x_i\) positiv in \(C_k\) vorkommt,
        und \(\{\overline{s_i},c_k\}\), falls \(x_i\) negativ in \(C_k\)
        vorkommt \quad(\(1\leq i\leq n\)); zusätzlich \(\{s_0,c_k\}\) für alle
        \(k\),
        \item[(K4)] \(\{w_j,c_k\}\), falls \(y_j\) positiv in \(C_k\) vorkommt,
        und \(\{\overline{w_j},c_k\}\), falls \(y_j\) negativ in \(C_k\)
        vorkommt,
        \item[(K5)] \(\{c_k,c_{k'}\}\) für alle \(k\neq k'\).
    \end{enumerate}
    Als Unabhängigkeitsanfrage setzen wir
    \[
        \mathcal{X}:=\{c_1,\dots,c_m\},\qquad
        \mathcal{Y}:=\{u_i,\overline{u_i}\mid x_i\in X^{+}\},\qquad
        Z:=\emptyset.
    \]
\end{Definition}

Das AF \(G_\Psi\) ist nach Konstruktion symmetrisch und irreflexiv, besitzt
\(4(n+1)+2l+m\) Argumente und ist in Polynomialzeit berechenbar. Die Blöcke
(K1) bilden für jede universelle Variable einen Viererkreis
\(u_i-\overline{u_i}-s_i-\overline{s_i}-u_i\); die Wertknoten
\(u_i,\overline{u_i}\) besitzen \emph{keine} weiteren Nachbarn, während die
Signalknoten \(s_i,\overline{s_i}\) die Verbindung zu den Klauselargumenten
herstellen.

\begin{Lemma}[Wertknoten]
    \label{lem:sym-wertknoten}
    Sei \(\lambda\in\Lambda_{\mathrm{st}}(G_\Psi)\) und
    \(S:=\mathit{in}(\lambda)\). Dann gilt für jedes \(x_i\in X^{+}\):
    \(|S\cap\{u_i,\overline{u_i}\}|=1\).
\end{Lemma}

\begin{proof}
    Wegen \(\{u_i,\overline{u_i}\}\in R_\Psi\) und der Konfliktfreiheit von
    \(S\) liegen nicht beide in \(S\). Seien nun beide nicht in \(S\). Nach
    Proposition~\ref{prop:sym-naiv-st-pr} ist \(S\) naiv, also besitzt
    \(u_i\) einen Nachbarn in \(S\). Die Nachbarn von \(u_i\) sind nach (K1)
    genau \(\overline{u_i}\) und \(\overline{s_i}\); folglich gilt
    \(\overline{s_i}\in S\). Analog besitzt \(\overline{u_i}\) die Nachbarn
    \(u_i\) und \(s_i\), also \(s_i\in S\). Wegen
    \(\{s_i,\overline{s_i}\}\in R_\Psi\) widerspricht das der
    Konfliktfreiheit von \(S\).
\end{proof}

Für \(\lambda\in\Lambda_{\mathrm{st}}(G_\Psi)\) mit
\(S=\mathit{in}(\lambda)\) definieren wir daher die \emph{induzierte Belegung}
\[
    \alpha_\lambda:X^{+}\to\{\top,\bot\},\quad
    \alpha_\lambda(x_i)=\top:\iff u_i\in S,
\]
sowie
\[
    \beta_\lambda:Y\to\{\top,\bot\},\quad
    \beta_\lambda(y_j)=\top:\iff w_j\in S.
\]

\begin{Lemma}[Signaltreue]
    \label{lem:sym-signaltreue}
    Sei \(\lambda\in\Lambda_{\mathrm{st}}(G_\Psi)\), \(S=\mathit{in}(\lambda)\).
    Dann gilt:
    \begin{enumerate}
        \item \(s_i\in S\Rightarrow\alpha_\lambda(x_i)=\top\) und
        \(\overline{s_i}\in S\Rightarrow\alpha_\lambda(x_i)=\bot\).
        \item \(w_j\in S\Rightarrow\beta_\lambda(y_j)=\top\) und
        \(\overline{w_j}\in S\Rightarrow\beta_\lambda(y_j)=\bot\).
    \end{enumerate}
\end{Lemma}

\begin{proof}
    (1) Aus \(s_i\in S\) und \(\{\overline{u_i},s_i\}\in R_\Psi\) folgt
    \(\overline{u_i}\notin S\), mit Lemma~\ref{lem:sym-wertknoten} also
    \(u_i\in S\) und damit \(\alpha_\lambda(x_i)=\top\). Der zweite Fall
    verläuft symmetrisch über \(\{u_i,\overline{s_i}\}\in R_\Psi\).

    (2) \(w_j\in S\) liefert \(\beta_\lambda(y_j)=\top\) definitionsgemäß; aus
    \(\overline{w_j}\in S\) folgt \(w_j\notin S\) wegen
    \(\{w_j,\overline{w_j}\}\in R_\Psi\), also \(\beta_\lambda(y_j)=\bot\).
\end{proof}

\begin{Lemma}[Realisierbare Teillabellings]
    \label{lem:sym-teile}
    Sei \(\lambda\in\Lambda_{\mathrm{st}}(G_\Psi)\) und
    \(S=\mathit{in}(\lambda)\). Dann gilt:
    \begin{enumerate}
        \item \(|S\cap\mathcal{X}|\leq 1\).
        \item Zu jedem \(k\in\{1,\dots,m\}\) und jeder Belegung
        \(\alpha:X^{+}\to\{\top,\bot\}\) existiert
        \(\lambda'\in\Lambda_{\mathrm{st}}(G_\Psi)\) mit
        \(\mathit{in}(\lambda')\cap\mathcal{X}=\{c_k\}\) und
        \(\alpha_{\lambda'}=\alpha\).
        \item Zu jeder Belegung \(\alpha:X^{+}\to\{\top,\bot\}\) existiert
        genau dann ein \(\lambda'\in\Lambda_{\mathrm{st}}(G_\Psi)\) mit
        \(\mathit{in}(\lambda')\cap\mathcal{X}=\emptyset\) und
        \(\alpha_{\lambda'}=\alpha\), wenn eine Belegung \(\beta\) von \(Y\)
        mit \(\widehat{\varphi}(\alpha,\beta)=\top\) existiert.
        \item Zu jeder Belegung \(\alpha:X^{+}\to\{\top,\bot\}\) existiert
        \(\lambda'\in\Lambda_{\mathrm{st}}(G_\Psi)\) mit
        \(\alpha_{\lambda'}=\alpha\).
    \end{enumerate}
\end{Lemma}

\begin{proof}
    Wir schreiben \(U_\alpha:=\{u_i\mid\alpha(x_i)=\top\}
    \cup\{\overline{u_i}\mid\alpha(x_i)=\bot\}\) und
    \(S_\alpha:=\{s_i\mid\alpha(x_i)=\top\}
    \cup\{\overline{s_i}\mid\alpha(x_i)=\bot\}\).
    Nach (K1) ist \(U_\alpha\cup S_\alpha\) konfliktfrei: Innerhalb eines
    Viererkreises werden nur gegenüberliegende Knoten gewählt
    (\(u_i\) mit \(s_i\) beziehungsweise \(\overline{u_i}\) mit
    \(\overline{s_i}\)), und zwischen verschiedenen Blöcken bestehen keine
    Konflikte. Ferner besitzt jedes Argument aus
    \(\{u_i,\overline{u_i}\}\setminus U_\alpha\) seinen Partner in
    \(U_\alpha\) als Nachbarn. Nach Korollar~\ref{cor:sym-extend} gilt
    daher für jede Erweiterung eines \(U_\alpha\) enthaltenden konfliktfreien
    Kerns zu einem \textit{stable}-Labelling \(\lambda'\) bereits
    \(\alpha_{\lambda'}=\alpha\).

    (1) Die Klauselargumente bilden nach (K5) eine Clique; eine konfliktfreie
    Menge enthält daher höchstens eines von ihnen.

    (2) Der Kern \(T:=U_\alpha\cup\{c_k\}\) ist konfliktfrei, da Wertknoten
    nach (K1) nur Wert- und Signalknoten desselben Blocks als Nachbarn haben
    und insbesondere zu keinem Klauselargument benachbart sind. Sei
    \(\lambda'\) eine Erweiterung gemäß Korollar~\ref{cor:sym-extend}. Dann gilt
    \(\alpha_{\lambda'}=\alpha\) nach der Vorbemerkung, und jedes \(c_{k'}\)
    mit \(k'\neq k\) besitzt den Nachbarn \(c_k\in T\), ist also
    \(\mathit{out}\). Somit ist
    \(\mathit{in}(\lambda')\cap\mathcal{X}=\{c_k\}\).

    (3) \enquote{\(\Leftarrow\)}: Sei \(\beta\) mit
    \(\widehat{\varphi}(\alpha,\beta)=\top\) und sei
    \[
        T:=U_\alpha\cup S_\alpha\cup
        \{w_j\mid\beta(y_j)=\top\}\cup\{\overline{w_j}\mid\beta(y_j)=\bot\}.
    \]
    \(T\) ist konfliktfrei: Konflikte bestehen nach (K3), (K4) nur zwischen
    Signal- beziehungsweise Literalargumenten und Klauselargumenten, und
    \(T\) enthält keine Klauselargumente; innerhalb der Blöcke wurde die
    Konfliktfreiheit bereits festgestellt. Da \(\widehat{\varphi}\) unter
    \((\alpha,\beta)\) erfüllt ist, besitzt jede Klausel \(x_0\lor C_k\) ein
    wahres Literal, dessen zugehöriges Signal- oder Literalargument nach (K3),
    (K4) in \(T\) liegt und zu \(c_k\) benachbart ist. Nach
    Korollar~\ref{cor:sym-extend} existiert ein \textit{stable}-Labelling
    \(\lambda'\) mit \(T\subseteq\mathit{in}(\lambda')\); dort sind alle
    \(c_k\) \(\mathit{out}\) und es gilt \(\alpha_{\lambda'}=\alpha\).

    \enquote{\(\Rightarrow\)}: Sei \(\lambda'\) mit
    \(\mathit{in}(\lambda')\cap\mathcal{X}=\emptyset\) und
    \(\alpha_{\lambda'}=\alpha\), und sei \(S'=\mathit{in}(\lambda')\). Da
    \(S'\) naiv ist, besitzt jedes \(c_k\) einen Nachbarn in \(S'\). Die
    Nachbarn von \(c_k\) sind nach (K3)--(K5) Signalargumente,
    Literalargumente und andere Klauselargumente; letztere liegen nicht in
    \(S'\). Also enthält \(S'\) ein Signal- oder Literalargument zu einem
    Literal von \(x_0\lor C_k\). Nach Lemma~\ref{lem:sym-signaltreue} ist
    dieses Literal unter \((\alpha,\beta_{\lambda'})\) wahr. Da \(k\) beliebig
    war, gilt \(\widehat{\varphi}(\alpha,\beta_{\lambda'})=\top\).

    (4) Man erweitere \(U_\alpha\) nach Korollar~\ref{cor:sym-extend}.
\end{proof}

\begin{Theorem}
    \label{th:sym-st-pi2p}
    \(\textit{Ind}^{\mathrm{sym}}_{\mathrm{st}}\) ist
    \(\Pi_2^\text{P}\)-vollständig. Die Härte gilt bereits für
    \(Z=\emptyset\).
\end{Theorem}

\begin{proof}
    \emph{Zugehörigkeit.} Symmetrische irreflexive AFs bilden eine Teilklasse
    aller AFs; die obere Schranke folgt daher aus
    Theorem~\ref{th:ind-st-is-pi2p-c}.

    \emph{Härte.} Sei \(\Psi=\forall X\exists Y\,\varphi\) eine Instanz des
    \(\Pi_2^\text{P}\)-vollständigen Problems \(\mathrm{QBF}^2_\forall\) und sei
    \(G_\Psi\) mit \(\mathcal{X},\mathcal{Y},Z=\emptyset\) wie in
    Definition~\ref{def:sym-konstruktion}. Die Mengen \(\mathcal{X}\),
    \(\mathcal{Y}\) und \(Z\) sind paarweise disjunkt, und \(G_\Psi\) ist
    symmetrisch und irreflexiv. Wegen \(Z=\emptyset\) ist die Voraussetzung
    \(\lambda_1|_Z=\lambda_2|_Z\) aus Definition~\ref{def:BU:AFs} für jedes
    Paar erfüllt. Wir zeigen
    \[
        \Psi\text{ ist gültig}
        \quad\Longleftrightarrow\quad
        \mathcal{X}\bot_{\mathrm{st}}\mathcal{Y}\mid\emptyset
        \text{ in }G_\Psi.
    \]

    Vorab halten wir fest, dass wegen
    \(\widehat{\varphi}=\bigwedge_k(x_0\lor C_k)\) gilt:
    \(\Psi\) ist genau dann gültig, wenn
    \(\forall X^{+}\exists Y\,\widehat{\varphi}\) gültig ist. Für
    \(\alpha(x_0)=\top\) ist \(\widehat{\varphi}\) unter jeder Belegung von
    \(Y\) erfüllt, für \(\alpha(x_0)=\bot\) gilt
    \(\widehat{\varphi}\equiv\varphi\).

    \enquote{\(\Rightarrow\)}: Sei \(\Psi\) gültig und seien
    \(\lambda_1,\lambda_2\in\Lambda_{\mathrm{st}}(G_\Psi)\). Sei
    \(\alpha:=\alpha_{\lambda_2}\); nach Lemma~\ref{lem:sym-wertknoten} ist
    \(\lambda_2|_{\mathcal{Y}}\) durch \(\alpha\) eindeutig bestimmt und
    umgekehrt. Nach Lemma~\ref{lem:sym-teile}(1) ist
    \(\mathit{in}(\lambda_1)\cap\mathcal{X}\) entweder leer oder von der Form
    \(\{c_k\}\).

    Im Fall \(\mathit{in}(\lambda_1)\cap\mathcal{X}=\{c_k\}\) liefert
    Lemma~\ref{lem:sym-teile}(2) ein \(\lambda_3\) mit
    \(\mathit{in}(\lambda_3)\cap\mathcal{X}=\{c_k\}\) und
    \(\alpha_{\lambda_3}=\alpha\), also
    \(\lambda_3|_{\mathcal{X}}=\lambda_1|_{\mathcal{X}}\) und
    \(\lambda_3|_{\mathcal{Y}}=\lambda_2|_{\mathcal{Y}}\).

    Im Fall \(\mathit{in}(\lambda_1)\cap\mathcal{X}=\emptyset\) existiert wegen
    der Gültigkeit von \(\forall X^{+}\exists Y\,\widehat{\varphi}\) eine
    Belegung \(\beta\) mit \(\widehat{\varphi}(\alpha,\beta)=\top\);
    Lemma~\ref{lem:sym-teile}(3) liefert das gesuchte \(\lambda_3\).

    Damit gilt \(\mathcal{X}\bot_{\mathrm{st}}\mathcal{Y}\mid\emptyset\).

    \enquote{\(\Leftarrow\)}: Sei \(\Psi\) nicht gültig. Dann existiert eine
    Belegung \(\alpha\) von \(X\), sodass \(\varphi(\alpha,\beta)=\bot\) für
    alle \(\beta\) gilt. Wir setzen
    \(\alpha^{+}:=\alpha\cup\{x_0\mapsto\bot\}\); dann gilt
    \(\widehat{\varphi}(\alpha^{+},\beta)=\bot\) für alle \(\beta\).

    Wähle \(\lambda_2\) nach Lemma~\ref{lem:sym-teile}(4) mit
    \(\alpha_{\lambda_2}=\alpha^{+}\). Wähle \(\lambda_1\) nach
    Lemma~\ref{lem:sym-teile}(3), angewandt auf eine beliebige Belegung
    \(\alpha'\) mit \(\alpha'(x_0)=\top\); wegen
    \(\widehat{\varphi}(\alpha',\beta)=\top\) für jedes \(\beta\) existiert ein
    solches \(\lambda_1\) mit
    \(\mathit{in}(\lambda_1)\cap\mathcal{X}=\emptyset\).

    Angenommen, es existierte \(\lambda_3\in\Lambda_{\mathrm{st}}(G_\Psi)\) mit
    \(\lambda_3|_{\mathcal{X}}=\lambda_1|_{\mathcal{X}}\) und
    \(\lambda_3|_{\mathcal{Y}}=\lambda_2|_{\mathcal{Y}}\). Dann gilt
    \(\mathit{in}(\lambda_3)\cap\mathcal{X}=\emptyset\) und
    \(\alpha_{\lambda_3}=\alpha^{+}\). Nach Lemma~\ref{lem:sym-teile}(3)
    existierte dann ein \(\beta\) mit
    \(\widehat{\varphi}(\alpha^{+},\beta)=\top\), im Widerspruch zur Wahl von
    \(\alpha\). Also gilt
    \(\mathcal{X}\not\bot_{\mathrm{st}}\mathcal{Y}\mid\emptyset\).

    Die Konstruktion ist in Polynomialzeit berechenbar, womit die Härte gezeigt
    ist.
\end{proof}

\begin{Corollary}
    \label{cor:sym-nicht-ocf}
    Theorem~\ref{th:sym-st-pi2p} folgt nicht aus
    Proposition~\ref{prop:ocf-af-ci-st-pi^2}.
\end{Corollary}

\begin{proof}
    Für \(m\geq 3\) enthält \(G_\Psi\) wegen (K5) den Dreierkreis
    \(c_1\to c_2\to c_3\to c_1\) und ist damit nicht ungeradezykelfrei. Ebenso
    ist \(G_\Psi\) wegen dieses Dreiecks nicht bipartit. Die Klasse der
    symmetrischen irreflexiven AFs ist somit weder in der Klasse der
    ungeradezykelfreien noch in der der bipartiten AFs enthalten.
\end{proof}

Bemerkenswert ist, dass symmetrische irreflexive AFs dennoch kohärent sind: Die
Kohärenz wird hier nicht über die Abwesenheit ungerader Zyklen erreicht, sondern
über die Symmetrie selbst (Lemma~\ref{lem:sym-cf-adm}). Die Zyklenfreiheit ist
also hinreichend, aber nicht notwendig für Kohärenz, und für die Komplexität
der Unabhängigkeit ist offenkundig nicht die Kohärenz entscheidend.

\section{Konsequenz für die \textit{preferred}-Semantik}
\label{sec:sym-preferred}

\begin{Lemma}[Transferlemma]
    \label{lem:sym-transfer}
    Sei \(F=(A,R)\) ein AF und seien \(\sigma,\tau\) Semantiken mit
    \(\Lambda_\sigma(F)=\Lambda_\tau(F)\). Dann gilt für alle paarweise
    disjunkten \(X,Y,Z\subseteq A\):
    \(X\bot_\sigma Y\mid Z\) genau dann, wenn \(X\bot_\tau Y\mid Z\).
\end{Lemma}

\begin{proof}
    Die definierende Bedingung aus Definition~\ref{def:BU:AFs} quantifiziert
    ausschließlich über Elemente von \(\Lambda_\sigma(F)\). Ersetzt man diese
    Menge durch die nach Voraussetzung identische Menge \(\Lambda_\tau(F)\), so
    geht die Bedingung Wort für Wort in die entsprechende Bedingung für
    \(\tau\) über.
\end{proof}

\begin{Corollary}
    \label{cor:sym-pr-pi2p}
    \(\textit{Ind}^{\mathrm{sym}}_{\mathrm{pr}}\) ist
    \(\Pi_2^\text{P}\)-vollständig.
\end{Corollary}

\begin{proof}
    Nach Proposition~\ref{prop:sym-naiv-st-pr} gilt
    \(\Lambda_{\mathrm{pr}}(F)=\Lambda_{\mathrm{st}}(F)\) für jedes
    symmetrische irreflexive AF \(F\). Mit Lemma~\ref{lem:sym-transfer} sind
    \(\textit{Ind}^{\mathrm{sym}}_{\mathrm{pr}}\) und
    \(\textit{Ind}^{\mathrm{sym}}_{\mathrm{st}}\) dieselbe Sprache; die
    Behauptung folgt aus Theorem~\ref{th:sym-st-pi2p}.
\end{proof}

Für \(\sigma=\mathrm{pr}\) ist die Klasse also \emph{echt} einfacher: Das
allgemeine Problem ist nach Theorem~\ref{th:ind-pr-is-pi3p-c}
\(\Pi_3^\text{P}\)-vollständig, auf symmetrischen irreflexiven AFs sinkt es um
eine Stufe der polynomiellen Hierarchie. Der Grund ist strukturell: Die
Maximalitätsbedingung der \textit{preferred}-Semantik, die im allgemeinen Fall
den zusätzlichen Quantorenwechsel erzwingt, kollabiert hier zu einer in
Polynomialzeit prüfbaren Bedingung (Proposition~\ref{prop:sym-klassisch-p}(2)).
Unter der üblichen Annahme \(\Pi_2^\text{P}\neq\Pi_3^\text{P}\) ist dieser
Gewinn nicht nur formal.

Für \(\sigma=\mathrm{st}\) tritt dagegen kein Gewinn ein. Die Erklärung liefert
Proposition~\ref{prop:sym-real-npc}: Die Symmetrie vereinfacht alle
\emph{punktweisen} Fragen über \(\Lambda_{\mathrm{st}}(F)\), lässt aber die
\emph{simultane} Realisierbarkeit vorgeschriebener Teillabellings NP-hart. Da
die bedingte Unabhängigkeit genau von dieser simultanen Realisierbarkeit
handelt, überträgt sich die Vereinfachung nicht.

\section{Ein in Polynomialzeit entscheidbares Fragment}
\label{sec:sym-fragment}

Die Härte in Theorem~\ref{th:sym-st-pi2p} benötigt eine unbeschränkt große
Menge \(\mathcal{Y}\). Beschränkt man beide Mengen auf einzelne Argumente und
setzt \(Z=\emptyset\), so wird das Problem tatsächlich einfach. Damit ist die
Klasse in einem präzisen Sinn doch einfacher als der allgemeine Fall, in dem
schon die zugrundeliegenden Realisierbarkeitsfragen NP-vollständig sind.

\begin{Proposition}
    \label{prop:sym-singleton-p}
    Sei \(F=(A,R)\) symmetrisch und irreflexiv und seien \(x,y\in A\) mit
    \(x\neq y\). Dann ist \(\{x\}\bot_{\mathrm{st}}\{y\}\mid\emptyset\) in Zeit
    \(\mathcal{O}(|A|^2+|A|\cdot|R|)\) entscheidbar.
\end{Proposition}

\begin{proof}
    Wegen \(Z=\emptyset\) ist die Bedingung \(\lambda_1|_Z=\lambda_2|_Z\)
    stets erfüllt. Setze
    \[
        \mathcal{R}:=\bigl\{(\lambda(x),\lambda(y))\;\big|\;
        \lambda\in\Lambda_{\mathrm{st}}(F)\bigr\}
        \subseteq\{\mathit{in},\mathit{out}\}^2 .
    \]
    Nach Definition~\ref{def:BU:AFs} gilt
    \(\{x\}\bot_{\mathrm{st}}\{y\}\mid\emptyset\) genau dann, wenn aus
    \((\alpha_1,\beta_1),(\alpha_2,\beta_2)\in\mathcal{R}\) stets
    \((\alpha_1,\beta_2)\in\mathcal{R}\) folgt. Da \(\mathcal{R}\) höchstens
    vier Elemente besitzt, genügt es, die vier Zugehörigkeiten zu bestimmen.
    Wir benutzen durchgehend Proposition~\ref{prop:sym-naiv-st-pr} und
    Korollar~\ref{cor:sym-extend}.

    \emph{\((\mathit{in},\mathit{in})\).} Genau dann realisierbar, wenn
    \((x,y)\notin R\): Ist \(\{x,y\}\) konfliktfrei, erweitere man es; andernfalls
    verbietet die Konfliktfreiheit beide gemeinsam.

    \emph{\((\mathit{in},\mathit{out})\).} Ist \((x,y)\in R\), so erweitere man
    \(\{x\}\); dann ist \(y\) \(\mathit{out}\). Ist \((x,y)\notin R\), so ist
    das Paar genau dann realisierbar, wenn \(N(y)\setminus N[x]\neq\emptyset\).
    Existiert nämlich ein solches \(w\), so ist \(\{x,w\}\) konfliktfrei und
    jede Erweiterung enthält \(x\), aber wegen \(w\in N(y)\) nicht \(y\).
    Existiert umgekehrt \(\lambda\) mit \(\lambda(x)=\mathit{in}\) und
    \(\lambda(y)=\mathit{out}\), so besitzt \(y\) einen Nachbarn
    \(w\in\mathit{in}(\lambda)\); wegen \(x\notin N(y)\) ist \(w\neq x\), und
    wegen der Konfliktfreiheit gilt \(w\notin N(x)\).

    \emph{\((\mathit{out},\mathit{in})\).} Symmetrisch mit vertauschten Rollen.

    \emph{\((\mathit{out},\mathit{out})\).} Genau dann realisierbar, wenn
    \(w_1\in N(x)\setminus\{y\}\) und \(w_2\in N(y)\setminus\{x\}\) mit
    \(w_1=w_2\) oder \((w_1,w_2)\notin R\) existieren. Sind solche \(w_1,w_2\)
    gegeben, so ist \(\{w_1,w_2\}\) konfliktfrei, und jede Erweiterung schließt
    \(x\) und \(y\) aus. Existiert umgekehrt \(\lambda\) mit
    \(\lambda(x)=\lambda(y)=\mathit{out}\), so besitzen \(x\) und \(y\) je einen
    Nachbarn \(w_1,w_2\in\mathit{in}(\lambda)\); diese liegen wegen
    \(\lambda(x)=\lambda(y)=\mathit{out}\) nicht in \(\{x,y\}\) und stehen
    untereinander nicht in Konflikt.

    Alle vier Tests sind durch Absuchen von \(N(x)\), \(N(y)\) beziehungsweise
    aller Paare \((w_1,w_2)\in N(x)\times N(y)\) in der angegebenen Zeit
    durchführbar. Anschließend prüft man die Rechteckbedingung auf der höchstens
    vierelementigen Menge \(\mathcal{R}\) in konstanter Zeit.
\end{proof}

Im allgemeinen Fall sind bereits die vier Realisierbarkeitstests
NP-vollständig, da \((\mathit{in},\cdot)\) die kredulöse Akzeptanz unter
\textit{stable}-Semantik enthält. Proposition~\ref{prop:sym-singleton-p} ist
also eine echte Vereinfachung; sie zeigt zugleich, dass die Härte aus
Theorem~\ref{th:sym-st-pi2p} nicht aus einzelnen Argumentpaaren stammt, sondern
aus der \emph{gemeinsamen} Vorschrift vieler Labels.

Als Konsistenzprüfung sei angemerkt, dass
Proposition~\ref{prop:cdg-afs-not-st-ci} ein Spezialfall dieser Analyse ist: Ein
vollständig gerichtetes AF ist symmetrisch und irreflexiv, seine naiven Mengen
sind genau die einelementigen Teilmengen von \(A\), und für \(x\neq y\) ist
\((\mathit{in},\mathit{in})\) wegen \((x,y)\in R\) nicht realisierbar, während
\((\mathit{in},\mathit{out})\) und \((\mathit{out},\mathit{in})\) realisierbar
sind. Die Rechteckbedingung ist damit verletzt.

\section{Die \textit{complete}-Semantik bleibt offen}
\label{sec:sym-complete}

Für \(\sigma=\mathrm{co}\) lässt sich die Konstruktion aus
Definition~\ref{def:sym-konstruktion} nicht übernehmen. Wir halten zunächst die
Struktur der \textit{complete}-Labellings fest.

\begin{Proposition}
    \label{prop:sym-co-struktur}
    Sei \(F=(A,R)\) symmetrisch und irreflexiv und sei \(\lambda\) ein
    Labelling mit \(S:=\mathit{in}(\lambda)\). Dann ist \(\lambda\) genau dann
    \textit{complete}, wenn \(S\) konfliktfrei ist,
    \(\mathit{out}(\lambda)=N(S)\) gilt und kein \(a\in A\setminus S\) mit
    \(N(a)\subseteq N(S)\) existiert.
\end{Proposition}

\begin{proof}
    Nach Definition~\ref{def:complete-labelling} ist \(a\) genau dann
    \(\mathit{out}\), wenn \(a\) einen \(\mathit{in}\)-Angreifer besitzt, also
    genau dann, wenn \(a\in N(S)\); das ist gleichbedeutend mit
    \(\mathit{out}(\lambda)=N(S)\), und \(S\cap N(S)=\emptyset\) ist genau die
    Konfliktfreiheit von \(S\). Ferner ist \(a\) genau dann \(\mathit{in}\),
    wenn alle Angreifer \(\mathit{out}\) sind, also genau dann, wenn
    \(N(a)\subseteq N(S)\). Für \(a\in S\) ist diese Bedingung wegen der
    Symmetrie automatisch erfüllt; die verbleibende Forderung ist, dass kein
    \(a\) außerhalb von \(S\) sie erfüllt.
\end{proof}

Insbesondere ist das \textit{grounded}-Labelling durch
\(\mathit{in}(\lambda_{\mathrm{gr}})=\{a\in A\mid N(a)=\emptyset\}\) gegeben:
Isolierte Argumente sind unangegriffen, und ein nicht isoliertes Argument \(a\)
besitzt einen Nachbarn, der von einem isolierten Argument nicht angegriffen
werden kann.

Damit scheitert die Reduktion aus Abschnitt~\ref{sec:sym-hauptresultat} an
Lemma~\ref{lem:sym-wertknoten}: Wählt man \(S=\emptyset\), so ist
\(N(S)=\emptyset\), und nach Proposition~\ref{prop:sym-co-struktur} ist das
Labelling, das alle nicht isolierten Argumente mit \(\mathit{und}\) versieht,
\textit{complete}. In diesem Labelling tragen sämtliche Wertknoten
\(u_i,\overline{u_i}\) das Label \(\mathit{und}\); die Zuordnung
\enquote{\(\lambda_2\) kodiert eine totale Belegung} bricht zusammen, und der
äußere Allquantor läuft über wesentlich mehr Labellings als über die
Belegungen von \(X^{+}\). Da zusätzliche \(\lambda_2\) zusätzliche
Kombinationsforderungen erzeugen, kann die Richtung
\enquote{\(\Psi\) gültig \(\Rightarrow\) unabhängig} verletzt werden.

In gerichteten Konstruktionen wird dieses Problem durch Detektorargumente
gelöst, die partielle Labellings erkennen und die eigentliche Formelauswertung
neutralisieren --- so etwa die Argumente \(d_i\), \(s\) und \(g\) im Beweis von
Proposition~\ref{prop:ocf-af-ci-co-pi^2-c}. Ein solcher Detektor benötigt einen
gerichteten Informationsfluss \(\{x_i,\overline{x_i}\}\to d_i\to s\to g\), der
sich in symmetrischen AFs nicht nachbilden lässt: Jeder Detektor griffe seine
Quelle zurück an. Wir halten daher fest:

\begin{Proposition}
    \label{prop:sym-co-obere-schranke}
    \(\textit{Ind}^{\mathrm{sym}}_{\mathrm{co}}\in\Pi_2^\text{P}\).
\end{Proposition}

\begin{proof}
    Folgt aus Theorem~\ref{th:ind-co-is-pi2p-c}, da symmetrische irreflexive
    AFs eine Teilklasse aller AFs bilden.
\end{proof}

Ob \(\textit{Ind}^{\mathrm{sym}}_{\mathrm{co}}\) auch \(\Pi_2^\text{P}\)-hart
ist, bleibt offen. Es ist die einzige der drei betrachteten Semantiken, für die
in dieser Klasse keine vollständige Einordnung erreicht wurde.

\section{Zusammenfassung}
\label{sec:sym-zusammenfassung}

\begin{table}[htbp]
    \myfloatalign
    \centering
    \renewcommand{\arraystretch}{1.15}
    \setlength{\tabcolsep}{1.0em}
    \begin{tabular}{lccc}
        \toprule
        & \(co\) & \(st\) & \(pr\) \\
        \midrule
        allgemeine AFs
            & \(\Pi_2^\text{P}\)-c
            & \(\Pi_2^\text{P}\)-c
            & \(\Pi_3^\text{P}\)-c \\
        symmetrisch, irreflexiv
            & in \(\Pi_2^\text{P}\)
            & \(\Pi_2^\text{P}\)-c
            & \(\Pi_2^\text{P}\)-c \\
        \quad davon \(|X|=|Y|=1\), \(Z=\emptyset\)
            & ---
            & in P
            & in P \\
        \bottomrule
    \end{tabular}
    \caption[Komplexität der Unabhängigkeit auf symmetrischen irreflexiven
    AFs]{Einordnung von \(\sigma\)-Unabhängigkeit auf symmetrischen
    irreflexiven AFs im Vergleich zum allgemeinen Fall. Die Einträge der
    zweiten Zeile stammen aus Proposition~\ref{prop:sym-co-obere-schranke},
    Theorem~\ref{th:sym-st-pi2p} und Korollar~\ref{cor:sym-pr-pi2p}, der
    Eintrag der dritten Zeile aus Proposition~\ref{prop:sym-singleton-p} in
    Verbindung mit Korollar~\ref{cor:sym-pr-pi2p}.}
    \label{tab:sym-komplexitaet}
\end{table}

Die Antwort auf die Ausgangsfrage ist differenziert. Symmetrische irreflexive
AFs sind für die \textit{preferred}-Semantik echt einfacher, weil dort die
Maximalitätsbedingung kollabiert und ein Quantorenwechsel entfällt. Für die
\textit{stable}-Semantik sind sie \emph{nicht} einfacher: Trotz kohärenter
Struktur, garantierter Existenz von Labellings und durchgängig polynomieller
klassischer Entscheidungsprobleme bleibt die bedingte Unabhängigkeit
\(\Pi_2^\text{P}\)-vollständig. Eine Vereinfachung tritt erst ein, wenn man
zusätzlich die Größe der Anfragemengen beschränkt.

Methodisch ist das der interessanteste Befund dieses Kapitels: Die
Tractability der klassischen Aufgaben ist kein verlässlicher Indikator für die
Komplexität der bedingten Unabhängigkeit. Maßgeblich ist stattdessen, ob
vorgeschriebene Teillabellings \emph{simultan} realisierbar sind
(Proposition~\ref{prop:sym-real-npc}) und ob der äußere Allquantor über
hinreichend viele, strukturell unterscheidbare Labellings läuft
(Lemma~\ref{lem:sym-wertknoten}). Beide Bedingungen sind in symmetrischen
irreflexiven AFs erfüllt, obwohl jede einzelne Akzeptanzfrage trivial ist.
 in thesis.tex,
% nach \chapter{Berechnungstheoretische Vorteile bestimmter Graphklassen zur Bestimmung bedingter Unabhängigkeit}
\label{ch:graphclasses}

Die folgenden Abschnitte orientieren sich an \cite{handbook__of_formal_argumentation__vol__3:2024}.

\section{Azyklische AFs}
\label{sec:acyc-afs}

Für azyklische AFs fällt jedes existierende \emph{stable}-Labelling mit dem eindeutig bestimmten \emph{grounded}-Labelling zusammen. 
Folglich besitzt ein azyklisches AF höchstens ein \emph{stable}-Labelling. 
Dadurch ist die bedingte Unabhängigkeit bezüglich der \emph{stable}-Semantik für beliebige paarweise disjunkte Argumentmengen 
stets erfüllt.

\begin{Proposition}
\label{prop:acyc-af-ci-st-top}
    Sei F=(A,R) ein azyklisches AF. Dann ist
    \(\textit{Ind}_{\mathrm{st}}(F) \) immer wahr.
\end{Proposition}

\begin{proof}
    \label{pr:acyc-af-ci-st-top-true}
    Seien \( X,Y,Z \subseteq A\) paarweise disjunkt. Besitzt \(F\) kein
    \emph{stable}-Labelling, so ist die universell quantifizierte
    Unabhängigkeitsbedingung vakuos erfüllt.

    Andernfalls existiert aufgrund der Azyklizität genau ein
    \emph{stable}-Labelling \(\lambda\), das mit dem
    \emph{grounded}-Labelling übereinstimmt. Für beliebige
    \(\lambda_1,\lambda_2\in\Lambda_{\mathrm{st}}(F)\) gilt daher
    \(\lambda_1=\lambda_2=\lambda\). Mit \(\lambda_3:=\lambda\) können
    somit die Belegungen von \(X\), \(Y\) und \(Z\) wie in der Definition
    der bedingten Unabhängigkeit gefordert kombiniert werden. Daher gilt
    \(X\bot_{\mathrm{st}}Y\mid Z\) und folglich
    \(\operatorname{Ind}_{\mathrm{st}}(F)\).
\end{proof}

\section{Geradezykelfreie AFs}
\label{sec:even-cycle-free-afs}

Für AFs ohne gerade Zykel ist allein das \emph{grounded}-Labelling ein Kandidat, ein \emph{stable}-Labelling zu sein. 
Damit existiert ebenfalls maximal ein \emph{stable}-Labelling. 

\begin{Proposition}
\label{prop:even-cycle-free-af-ci-st-top}
    Sei F=(A,R) ein AF ohne gerade Zyklen. Dann ist
    \(\textit{Ind}_{\mathrm{st}}(F) \) immer wahr.
\end{Proposition}

\begin{proof}
    Analog zu \ref{pr:acyc-af-ci-st-top-true}.
\end{proof}


\section{Ungeradezykelfreie AFs}
\label{sec:odd-cycle-free-afs}

Für die Berechnung von \emph{stable}-Labellings gibt es in dieser Graphklasse keine berechnungstheoretischen Vorteile. 


\section{Bipartite AFs}
\label{sec:bi-afs}

Als Spezialfall eines ungeradezykelfreien AFs gibt es für die Berechnung von \emph{stable}-Labellings in bipartiten AFs keine
berechnungstheoretischen Vorteile. 


\section{Symmetrische AFs}
\label{sec:sym-afs}

Beobachtung: Je mehr Angriffe im AF, desto häufiger nicht bedingt unabhängig.

.
% ****************************************************************************************************

\chapter{Analyse des Härtebeweises für bipartite AFs}
\label{ch:bip-analyse}

\section{Gegenstand und Vorgehen}
\label{sec:bip-analyse-gegenstand}

Dieses Kapitel untersucht Proposition~\ref{prop:bip-st-ind-hard} aus
Kapitel~\ref{ch:graphclasses}, also die Behauptung, dass
\textit{stable}-Unabhängigkeit auch auf bipartiten AFs
\(\Pi_2^\text{P}\)-vollständig ist. Die Untersuchung erfolgt in drei Schritten.
Zuerst wird die Konstruktion so rekonstruiert, dass ihre Semantik durch ein
explizites Strukturlemma erfasst wird (Abschnitt~\ref{sec:bip-analyse-struktur}).
Anschließend wird die Korrektheit vollständig nachgerechnet
(Abschnitt~\ref{sec:bip-analyse-korrektheit}) und die
Graphklassenzugehörigkeit verifiziert
(Abschnitt~\ref{sec:bip-analyse-bipartit}). Erst danach werden die gefundenen
Mängel gesammelt (Abschnitt~\ref{sec:bip-analyse-kritik}) und behoben
(Abschnitt~\ref{sec:bip-analyse-reparatur}).

Das Ergebnis vorweg: \emph{Die Reduktion ist inhaltlich korrekt.} Die
Beweisführung weist jedoch eine Reihe formaler Mängel auf, von denen einer
(die unbelegte Ausgangsprobleminstanz) den Beweis in seiner jetzigen Form
angreifbar macht, während die übrigen die Nachvollziehbarkeit betreffen.

\section{Rekonstruktion der Konstruktion}
\label{sec:bip-analyse-struktur}

Wir wiederholen die Konstruktion in fixierter Notation. Sei
\(\Psi=\forall X\exists Y\,\varphi\) mit
\(\varphi=\bigwedge_{C\in\mathcal{C}}C\) in konjunktiver Normalform, wobei jede
Klausel \emph{monoton} ist, also entweder ausschließlich positive Literale
(\(C\in\mathcal{C}^{+}\)) oder ausschließlich negative Literale
(\(C\in\mathcal{C}^{-}\)) enthält. Sei \(V:=X\cup Y\cup\{p,n\}\) mit zwei
frischen Variablen \(p,n\). Das AF \(F_\Psi=(A_\Psi,R_\Psi)\) ist gegeben durch
\[
\begin{aligned}
    A_{\Psi}
    ={}&
    \{v,\overline{v}\mid v\in V\}
    \cup
    \{c_C\mid C\in\mathcal{C}\}
    \cup
    \{\varphi^{+},\varphi^{-},t^{+},t^{-}\}
\end{aligned}
\]
und der in Proposition~\ref{prop:bip-st-ind-hard} angegebenen Angriffsmenge
\(R_\Psi\). Zentral ist die Beobachtung, dass die Argumente \(p\) und
\(\overline{n}\) genau so wirken, als käme das Literal \(p\) in jeder positiven
und das Literal \(\neg n\) in jeder negativen Klausel vor. Wir machen das
explizit und setzen
\[
    \widehat{\varphi}
    :=
    \bigwedge_{C\in\mathcal{C}^{+}}(p\lor C)
    \;\land\;
    \bigwedge_{C\in\mathcal{C}^{-}}(\neg n\lor C).
\]

\begin{Lemma}
    \label{lem:bip-qbf-aequiv}
    Es gilt
    \(
        \Psi\text{ ist gültig}
        \iff
        \forall (X\cup\{p,n\})\,\exists Y\,\widehat{\varphi}
        \text{ ist gültig.}
    \)
\end{Lemma}

\begin{proof}
    \enquote{\(\Rightarrow\)}: Seien \(\alpha\) eine Belegung von \(X\) und
    \(\pi,\nu\) beliebige Wahrheitswerte für \(p\) und \(n\). Da \(\Psi\)
    gültig ist, existiert \(\beta\) mit \(\varphi(\alpha,\beta)=\top\). Dann ist
    jede Klausel \(C\in\mathcal{C}\) unter \((\alpha,\beta)\) erfüllt, also erst
    recht \(p\lor C\) beziehungsweise \(\neg n\lor C\). Somit gilt
    \(\widehat{\varphi}=\top\).

    \enquote{\(\Leftarrow\)}: Sei \(\alpha\) eine Belegung von \(X\). Wir wählen
    die Belegung \(p=\bot\) und \(n=\top\). Dann gilt
    \(\widehat{\varphi}\equiv\varphi\), und die Voraussetzung liefert ein
    \(\beta\) mit \(\varphi(\alpha,\beta)=\top\).
\end{proof}

Der Beweis in Kapitel~\ref{ch:graphclasses} argumentiert an mehreren Stellen mit
Aussagen der Form \enquote{die Argumente \(t^{+}\) und \(t^{-}\) erlauben es,
die Labels von \(\varphi^{+}\) und \(\varphi^{-}\) unabhängig zu setzen}, ohne
diese zu beweisen. Das folgende Lemma stellt die vollständige Struktur von
\(\Lambda_{\mathrm{st}}(F_\Psi)\) bereit und macht alle späteren Schritte
elementar.

\begin{Lemma}[Strukturlemma]
    \label{lem:bip-struktur}
    Sei \(\lambda\) ein Labelling von \(F_\Psi\). Dann ist \(\lambda\) genau
    dann \textit{stable}, wenn es eine totale Belegung \(\tau:V\to\{\top,\bot\}\)
    gibt, sodass die folgenden vier Bedingungen gelten.
    \begin{enumerate}
        \item Für alle \(v\in V\) ist
        \(\lambda(v)=\mathit{in}\iff\tau(v)=\top\)
        und
        \(\lambda(\overline{v})=\mathit{in}\iff\tau(v)=\bot\).
        \item Für alle \(C\in\mathcal{C}^{+}\) ist
        \(\lambda(c_C)=\mathit{out}\) genau dann, wenn \(\tau\models p\lor C\);
        andernfalls \(\lambda(c_C)=\mathit{in}\). Analog ist für
        \(C\in\mathcal{C}^{-}\) genau dann \(\lambda(c_C)=\mathit{out}\), wenn
        \(\tau\models\neg n\lor C\).
        \item \(\lambda(t^{+})=\mathit{in}\iff\lambda(\varphi^{+})=\mathit{out}\)
        und \(\lambda(t^{-})=\mathit{in}\iff\lambda(\varphi^{-})=\mathit{out}\).
        \item \(\lambda(\varphi^{+})=\mathit{in}\) impliziert
        \(\lambda(c_C)=\mathit{out}\) für alle \(C\in\mathcal{C}^{+}\);
        \(\lambda(\varphi^{-})=\mathit{in}\) impliziert
        \(\lambda(c_C)=\mathit{out}\) für alle \(C\in\mathcal{C}^{-}\).
    \end{enumerate}
    Umgekehrt liefert jede Wahl von \(\tau\) und jede mit (3) und (4)
    verträgliche Wahl von \(\lambda(\varphi^{+}),\lambda(\varphi^{-})\)
    ein \textit{stable}-Labelling. Insbesondere ist
    \(\lambda(\varphi^{+})=\mathit{out}\) stets realisierbar, und
    \(\lambda(\varphi^{+})=\mathit{in}\) ist genau dann realisierbar, wenn
    \(\tau\) alle Klauseln \(p\lor C\) mit \(C\in\mathcal{C}^{+}\) erfüllt.
\end{Lemma}

\begin{proof}
    Sei \(\lambda\in\Lambda_{\mathrm{st}}(F_\Psi)\); insbesondere ist
    \(\mathit{und}(\lambda)=\emptyset\).

    Zu (1): Die einzigen Angreifer von \(v\) sind \(\overline{v}\) und
    umgekehrt. Wären beide \(\mathit{in}\), so griffe ein
    \(\mathit{in}\)-Argument ein \(\mathit{in}\)-Argument an. Wären beide
    \(\mathit{out}\), so besäße \(v\) keinen \(\mathit{in}\)-Angreifer. Also ist
    genau eines der beiden \(\mathit{in}\), und wir setzen
    \(\tau(v):=\top\) genau dann, wenn \(\lambda(v)=\mathit{in}\).

    Zu (2): Die Angreifer von \(c_C\) mit \(C\in\mathcal{C}^{+}\) sind genau
    die Argumente \(v\) mit \(v\in C\) sowie \(p\). Nach (1) ist ein solcher
    Angreifer genau dann \(\mathit{in}\), wenn das zugehörige Literal unter
    \(\tau\) wahr ist. Da \(\lambda\) \textit{stable} ist, gilt
    \(\lambda(c_C)=\mathit{out}\) genau dann, wenn mindestens ein Angreifer
    \(\mathit{in}\) ist, also genau dann, wenn \(\tau\models p\lor C\). Der Fall
    \(C\in\mathcal{C}^{-}\) verläuft analog mit den Angreifern \(\overline{v}\)
    für \(\neg v\in C\) und \(\overline{n}\).

    Zu (3): Der einzige Angreifer von \(t^{+}\) ist \(\varphi^{+}\). Also ist
    \(t^{+}\) genau dann \(\mathit{out}\), wenn \(\varphi^{+}\) \(\mathit{in}\)
    ist, und wegen \(\mathit{und}(\lambda)=\emptyset\) genau dann
    \(\mathit{in}\), wenn \(\varphi^{+}\) \(\mathit{out}\) ist.

    Zu (4): Die Angreifer von \(\varphi^{+}\) sind die \(c_C\) mit
    \(C\in\mathcal{C}^{+}\) sowie \(t^{+}\). Ist \(\varphi^{+}\)
    \(\mathit{in}\), so sind alle Angreifer \(\mathit{out}\).

    Für die Umkehrung sei \(\tau\) beliebig und seien
    \(\lambda(\varphi^{\pm})\) gemäß (3) und (4) gewählt; die übrigen Labels
    seien durch (1)--(3) festgelegt. Man verifiziert unmittelbar für jedes
    Argument die \textit{stable}-Bedingung: Argumente \(v,\overline{v}\) sind
    genau dann \(\mathit{out}\), wenn ihr Partner \(\mathit{in}\) ist;
    Klauselargumente genau dann, wenn ein Literalargument \(\mathit{in}\) ist;
    \(t^{\pm}\) genau dann, wenn \(\varphi^{\pm}\) \(\mathit{in}\) ist; und
    \(\varphi^{+}\) ist entweder \(\mathit{in}\) (dann sind nach (4) alle
    \(c_C\) und nach (3) auch \(t^{+}\) \(\mathit{out}\)) oder \(\mathit{out}\)
    (dann ist nach (3) \(t^{+}\) \(\mathit{in}\) und greift \(\varphi^{+}\) an).
    Es tritt kein \(\mathit{und}\) auf.

    Der Zusatz folgt: \(\lambda(\varphi^{+})=\mathit{out}\) ist wegen
    \(\lambda(t^{+})=\mathit{in}\) immer zulässig, und
    \(\lambda(\varphi^{+})=\mathit{in}\) ist nach (4) und (2) genau dann
    zulässig, wenn \(\tau\) jede Klausel \(p\lor C\), \(C\in\mathcal{C}^{+}\),
    erfüllt.
\end{proof}

Lemma~\ref{lem:bip-struktur} zeigt insbesondere: Die Labels von
\(\varphi^{+}\) und \(\varphi^{-}\) sind \emph{unabhängig voneinander} frei
wählbar, sobald die jeweiligen Klauselargumente \(\mathit{out}\) sind, da
\(\varphi^{+},t^{+}\) und \(\varphi^{-},t^{-}\) disjunkte Angreifermengen
besitzen. Genau diese Aussage benutzt der Originalbeweis, ohne sie zu begründen.

\section{Korrektheit der Reduktion}
\label{sec:bip-analyse-korrektheit}

\begin{Proposition}
    \label{prop:bip-korrekt}
    Es gilt
    \[
        \Psi\text{ ist gültig}
        \quad\Longleftrightarrow\quad
        \{\varphi^{+},\varphi^{-}\}\bot_{\mathrm{st}}(X\cup\{p,n\})\mid\emptyset
        \text{ in }F_\Psi.
    \]
\end{Proposition}

\begin{proof}
    Wir schreiben \(\mathcal{X}:=\{\varphi^{+},\varphi^{-}\}\) und
    \(\mathcal{Y}:=X\cup\{p,n\}\); beide Mengen sind disjunkt, und
    \(Z=\emptyset\), sodass die Bedingung \(\lambda_1|_Z=\lambda_2|_Z\) aus
    Definition~\ref{def:BU:AFs} für jedes Paar erfüllt ist. Gesucht ist also
    stets ein \(\lambda_3\in\Lambda_{\mathrm{st}}(F_\Psi)\) mit
    \(\lambda_3|_{\mathcal{X}}=\lambda_1|_{\mathcal{X}}\) und
    \(\lambda_3|_{\mathcal{Y}}=\lambda_2|_{\mathcal{Y}}\).

    \enquote{\(\Rightarrow\)}: Sei \(\Psi\) gültig und seien
    \(\lambda_1,\lambda_2\in\Lambda_{\mathrm{st}}(F_\Psi)\). Nach
    Lemma~\ref{lem:bip-struktur} induziert \(\lambda_2\) eine totale Belegung
    \(\tau_2\) von \(V\); insbesondere legt \(\lambda_2|_{\mathcal{Y}}\) die
    Werte \(\tau_2|_{X\cup\{p,n\}}\) fest. Nach Lemma~\ref{lem:bip-qbf-aequiv}
    existiert eine Belegung \(\beta\) von \(Y\), sodass
    \(\tau_3:=\tau_2|_{X\cup\{p,n\}}\cup\beta\) die Formel
    \(\widehat{\varphi}\) erfüllt. Wir wählen \(\lambda_3\) als das durch
    \(\tau_3\) und
    \(\lambda_3(\varphi^{+}):=\lambda_1(\varphi^{+})\),
    \(\lambda_3(\varphi^{-}):=\lambda_1(\varphi^{-})\)
    bestimmte Labelling. Da \(\tau_3\models\widehat{\varphi}\), sind nach
    Lemma~\ref{lem:bip-struktur}(2) sämtliche Klauselargumente
    \(\mathit{out}\); damit sind beide Wahlen für \(\varphi^{+}\) und
    \(\varphi^{-}\) mit (4) verträglich, und \(\lambda_3\) ist
    \textit{stable}. Nach Konstruktion gilt
    \(\lambda_3|_{\mathcal{X}}=\lambda_1|_{\mathcal{X}}\) und
    \(\lambda_3|_{\mathcal{Y}}=\lambda_2|_{\mathcal{Y}}\).

    \enquote{\(\Leftarrow\)}: Sei \(\Psi\) nicht gültig. Dann existiert eine
    Belegung \(\alpha\) von \(X\), sodass \(\varphi(\alpha,\beta)=\bot\) für
    jede Belegung \(\beta\) von \(Y\) gilt. Wir wählen:
    \begin{itemize}
        \item \(\lambda_1\) als ein beliebiges \textit{stable}-Labelling mit
        \(\lambda_1(\varphi^{+})=\lambda_1(\varphi^{-})=\mathit{in}\). Ein
        solches existiert: Man setze \(\tau_1(p)=\top\) und \(\tau_1(n)=\bot\)
        sowie \(\tau_1\) auf \(X\cup Y\) beliebig. Dann sind alle Klauseln
        \(p\lor C\) und \(\neg n\lor C\) erfüllt, also nach
        Lemma~\ref{lem:bip-struktur} alle Klauselargumente \(\mathit{out}\)
        und beide \(\varphi\)-Argumente \(\mathit{in}\) realisierbar.
        \item \(\lambda_2\) als ein \textit{stable}-Labelling zur Belegung
        \(\tau_2:=\alpha\cup\{p\mapsto\bot,\;n\mapsto\top\}\cup\beta_0\) mit
        beliebigem \(\beta_0\).
    \end{itemize}
    Angenommen, es gäbe ein \(\lambda_3\in\Lambda_{\mathrm{st}}(F_\Psi)\) mit
    \(\lambda_3|_{\mathcal{X}}=\lambda_1|_{\mathcal{X}}\) und
    \(\lambda_3|_{\mathcal{Y}}=\lambda_2|_{\mathcal{Y}}\). Sei \(\tau_3\) die
    zugehörige Belegung. Aus
    \(\lambda_3|_{\mathcal{Y}}=\lambda_2|_{\mathcal{Y}}\) folgt
    \(\tau_3|_X=\alpha\), \(\tau_3(p)=\bot\) und \(\tau_3(n)=\top\). Aus
    \(\lambda_3(\varphi^{+})=\lambda_3(\varphi^{-})=\mathit{in}\) folgt mit
    Lemma~\ref{lem:bip-struktur}(4) und (2), dass \(\tau_3\) sämtliche Klauseln
    \(p\lor C\) und \(\neg n\lor C\) erfüllt, also
    \(\tau_3\models\widehat{\varphi}\). Wegen \(\tau_3(p)=\bot\) und
    \(\tau_3(n)=\top\) gilt aber \(\widehat{\varphi}\equiv\varphi\) unter
    \(\tau_3\), also \(\varphi(\alpha,\tau_3|_Y)=\top\) im Widerspruch zur Wahl
    von \(\alpha\). Also existiert kein solches \(\lambda_3\), und es gilt
    \(\mathcal{X}\not\bot_{\mathrm{st}}\mathcal{Y}\mid\emptyset\).
\end{proof}

Die Konstruktion ist offensichtlich in Polynomialzeit berechenbar: Sie erzeugt
\(2|V|+|\mathcal{C}|+4\) Argumente und
\(\mathcal{O}(|V|+\sum_{C\in\mathcal{C}}|C|)\) Angriffe.

\section{Graphklassenzugehörigkeit}
\label{sec:bip-analyse-bipartit}

\begin{Proposition}
    \label{prop:bip-korrekt-bipartit}
    \(F_\Psi\) ist bipartit.
\end{Proposition}

\begin{proof}
    Wir setzen
    \[
    \begin{aligned}
        L:={}& \{v\mid v\in V\}\cup\{c_C\mid C\in\mathcal{C}^{-}\}
              \cup\{\varphi^{+},t^{-}\},\\
        R:={}& \{\overline{v}\mid v\in V\}\cup\{c_C\mid C\in\mathcal{C}^{+}\}
              \cup\{\varphi^{-},t^{+}\}.
    \end{aligned}
    \]
    Offensichtlich gilt \(L\cap R=\emptyset\) und \(L\cup R=A_\Psi\). Für jede
    Art von Angriff prüft man die Seitenwechsel-Eigenschaft:
    \((v,\overline{v})\) und \((\overline{v},v)\) verlaufen zwischen \(L\) und
    \(R\); \((v,c_C)\) mit \(C\in\mathcal{C}^{+}\) sowie \((p,c_C)\) verlaufen
    von \(L\) nach \(R\); \((\overline{v},c_C)\) mit \(C\in\mathcal{C}^{-}\)
    sowie \((\overline{n},c_C)\) verlaufen von \(R\) nach \(L\);
    \((c_C,\varphi^{+})\) mit \(C\in\mathcal{C}^{+}\) verläuft von \(R\) nach
    \(L\); \((c_C,\varphi^{-})\) mit \(C\in\mathcal{C}^{-}\) von \(L\) nach
    \(R\); \((\varphi^{+},t^{+})\) von \(L\) nach \(R\) und
    \((\varphi^{-},t^{-})\) von \(R\) nach \(L\), jeweils mit den
    Gegenangriffen. Also gilt
    \(R_\Psi\subseteq(L\times R)\cup(R\times L)\).
\end{proof}

Hier zeigt sich der eigentliche Grund für die Monotonieforderung: Nur weil
positive Klauseln ausschließlich von unquerstrichenen und negative Klauseln
ausschließlich von querstrichenen Literalargumenten angegriffen werden, lassen
sich die Klauselargumente überhaupt widerspruchsfrei auf die beiden Seiten
verteilen. Bei gemischten Klauseln entstünde ein Klauselargument mit Angreifern
aus \(L\) \emph{und} \(R\), und die Bipartition bräche zusammen.

\section{Kritikpunkte}
\label{sec:bip-analyse-kritik}

Die Reduktion ist nach den Abschnitten~\ref{sec:bip-analyse-korrektheit}
und~\ref{sec:bip-analyse-bipartit} korrekt. Der Beweis in
Kapitel~\ref{ch:graphclasses} weist jedoch die folgenden Mängel auf. Wir
ordnen sie nach abnehmender Schwere.

\begin{enumerate}
    \item \textbf{Unbelegtes Ausgangsproblem.} Der Beweis reduziert von
    \textsc{Balanced Monotone \(\forall\exists\)-3-SAT-\((1,1,2,2)\)} und
    behauptet ohne Quellenangabe und ohne Definition, dieses Problem sei
    \(\Pi_2^\text{P}\)-vollständig. Weder wird gesagt, worauf sich
    \enquote{balanced} bezieht, noch was das Tupel \((1,1,2,2)\) zählt. Das ist
    der einzige Mangel, der die \emph{Gültigkeit} der Aussage betrifft: Eine
    Reduktion von einem Problem unbekannter Komplexität beweist nichts.
    Vorkommensbeschränkungen sind zudem heikel. Nach dem Satz von Tovey ist
    jede \(k\)-CNF-Formel, in der jede Klausel \(k\) verschiedene Literale
    besitzt und jede Variable in höchstens \(k\) Klauseln vorkommt, erfüllbar;
    zu enge Vorkommensschranken können ein Problem also trivialisieren.
    \item \textbf{Die Zusatzrestriktionen werden nicht benötigt.} Die
    Konstruktion verwendet ausschließlich die Monotonie der Klauseln (siehe
    Abschnitt~\ref{sec:bip-analyse-bipartit}); weder die Ausgeglichenheit von
    \(|\mathcal{C}^{+}|\) und \(|\mathcal{C}^{-}|\) noch irgendeine
    Vorkommensschranke noch die Klauselbreite \(3\) geht in Korrektheit oder
    Bipartitheit ein. Der Beweis macht sich also ohne Not von einer stärkeren
    und schwerer zu belegenden Voraussetzung abhängig.
    Abschnitt~\ref{sec:bip-analyse-reparatur} zeigt, dass die Monotonie
    elementar und selbsttragend erreicht werden kann.
    \item \textbf{Vertauschte Rollen von \(\lambda_1\) und \(\lambda_2\).} In
    der Richtung \enquote{\(\Psi\) gültig} übernimmt \(\lambda_1\) die Labels
    von \(\varphi^{+},\varphi^{-}\) und \(\lambda_2\) die Belegung von
    \(X\cup\{p,n\}\), passend zur Reihenfolge der Mengen in
    \(\{\varphi^{+},\varphi^{-}\}\bot_{\mathrm{st}}X\cup\{p,n\}\mid\emptyset\)
    und zu Definition~\ref{def:BU:AFs}. In der Gegenrichtung sind die Rollen
    stillschweigend vertauscht: Dort trägt \(\lambda_1\) die Belegung und
    \(\lambda_2\) die \(\varphi\)-Labels. Inhaltlich ist das folgenlos, weil
    die bedingte Unabhängigkeit symmetrisch ist, formal ist es jedoch ein
    Bruch mit der eigenen Definition. Entweder man benennt konsistent um (wie
    in Proposition~\ref{prop:bip-korrekt}) oder man beruft sich explizit auf
    das Symmetrieaxiom.
    \item \textbf{Undefinierte Mengen \(A\) und \(B\).} Der Satz
    \enquote{[\ldots] das die Einschränkung von \(\lambda_1\) auf \(A\) und die
    Einschränkung von \(\lambda_2\) auf \(B\) kombiniert} verwendet zwei nie
    eingeführte Symbole. Erschwerend kommt hinzu, dass \(A\) in der gesamten
    Arbeit die Argumentmenge eines AFs bezeichnet; gemeint sind offenbar
    \(\{\varphi^{+},\varphi^{-}\}\) und \(X\cup\{p,n\}\).
    \item \textbf{Namenskollision \(L,R\).} Die Bipartitionsklassen heißen
    \(L\) und \(R\), während \(R\) beziehungsweise \(R_\Psi\) die
    Angriffsrelation bezeichnet. Die Schlussformel
    \(R_\Psi\subseteq(L\times R)\cup(R\times L)\) mischt beide Bedeutungen in
    einer Zeile. Vorzuziehen wären etwa \(A_1,A_2\) oder \(L_1,L_2\).
    \item \textbf{Fehlendes Strukturlemma.} Der Beweis benutzt mehrfach
    Eigenschaften von \(\Lambda_{\mathrm{st}}(F_\Psi)\), die er nicht beweist:
    dass in jedem \textit{stable}-Labelling genau eines von \(v,\overline{v}\)
    das Label \(\mathit{in}\) trägt; dass \(c_C\) genau dann \(\mathit{out}\)
    ist, wenn ein Literalargument \(\mathit{in}\) ist; und dass \(t^{+},t^{-}\)
    die freie und \emph{voneinander unabhängige} Wahl der \(\varphi\)-Labels
    erlauben. Lemma~\ref{lem:bip-struktur} schließt diese Lücke.
    \item \textbf{Unvollständige Fallunterscheidung über \(p\) und \(n\).} Der
    Abschnitt \enquote{Konstruktion} diskutiert nur die Kombinationen
    \((p,n)=(\mathit{out},\mathit{in})\) und
    \((p,n)=(\mathit{in},\mathit{out})\). Die ebenfalls in
    \textit{stable}-Labellings auftretenden Kombinationen
    \((\mathit{in},\mathit{in})\) und \((\mathit{out},\mathit{out})\) bleiben
    unerwähnt. Der Korrektheitsbeweis benötigt diese Fallunterscheidung
    tatsächlich nicht --- er läuft über die Formeläquivalenz aus
    Lemma~\ref{lem:bip-qbf-aequiv} ---, der Text suggeriert jedoch eine
    erschöpfende Analyse, die er nicht liefert.
    \item \textbf{Zugehörigkeit fehlt.} Die Proposition behauptet
    \(\Pi_2^\text{P}\)-\emph{Vollständigkeit}, der Beweis zeigt nur die Härte.
    Ein Verweis auf Theorem~\ref{th:ind-st-is-pi2p-c} (die obere Schranke gilt
    für alle AFs, also insbesondere für bipartite) fehlt, anders als beim
    ungeradezykelfreien Fall.
    \item \textbf{Formulierungs- und Notationsmängel.} Die Instanz wird als
    Behauptung formuliert (\enquote{wählen wir
    \(\{\varphi^{+},\varphi^{-}\}\perp_{st}\ldots\)}) statt als Anfrage; die
    Polynomialität der Konstruktion wird nicht erwähnt; \(\perp\) wird sowohl
    für die Unabhängigkeitsrelation als auch (in \enquote{\(p=\perp\)}) für den
    Wahrheitswert \emph{falsch} verwendet, während an anderer Stelle
    \(\bot_{\mathrm{st}}\) für die Unabhängigkeit steht; die Argumentmenge
    heißt hier \(\mathcal{A}_\Psi\), im ungeradezykelfreien Fall dagegen
    \(A_\Psi\); und der Satz \enquote{dann werden die Klauselargumente
    ausschließlich durch die ursprünglichen Literale kontrolliert werden}
    enthält ein doppeltes Verb.
\end{enumerate}

\begin{Corollary}
    \label{cor:bip-subsumiert-ocf}
    Proposition~\ref{prop:bip-st-ind-hard} impliziert
    Proposition~\ref{prop:ocf-af-ci-st-pi^2}.
\end{Corollary}

\begin{proof}
    Ist \(F\) bipartit, so wechselt jeder gerichtete Zyklus in jedem Schritt die
    Bipartitionsklasse und hat daher gerade Länge. Bipartite AFs sind also
    ungeradezykelfrei, und die Härte auf der Teilklasse überträgt sich auf die
    Oberklasse.
\end{proof}

Das ist ein struktureller Hinweis, kein Fehler: Der eigenständige Beweis für
ungeradezykelfreie AFs für \(\sigma=\mathrm{st}\) und \(\sigma=\mathrm{pr}\) ist
entbehrlich, sobald das bipartite Resultat gesichert ist. Umgekehrt bleibt der
gesonderte Beweis für \(\sigma=\mathrm{co}\)
(Proposition~\ref{prop:ocf-af-ci-co-pi^2-c}) nötig, denn das dort verwendete AF
ist wegen der Argumente \(d_i\), die sowohl von \(x_i\) als auch von
\(\overline{x_i}\) angegriffen werden, nicht bipartit.

Schließlich sei angemerkt, dass die Herleitung von
Proposition~\ref{prop:bip-prf-ind-hard} aus dem \textit{stable}-Fall die
Kohärenz ungeradezykelfreier AFs voraussetzt, also
\(\Lambda_{\mathrm{st}}(F)=\Lambda_{\mathrm{pr}}(F)\). Diese Tatsache stammt von
\textcite{dung:1995:af} und sollte dort zitiert werden.

\section{Reparatur: Monotonie ohne fremde Voraussetzung}
\label{sec:bip-analyse-reparatur}

Der schwerwiegendste Kritikpunkt lässt sich vollständig beheben. Statt sich auf
eine exotische Restriktion zu berufen, genügt eine elementare Normalform, die
sich in wenigen Zeilen beweisen lässt.

\begin{Lemma}[Monotone Normalform]
    \label{lem:monotone-normalform}
    Zu jeder Formel \(\Psi=\forall X\exists Y\,\varphi\) mit \(\varphi\) in KNF
    lässt sich in Polynomialzeit eine Formel
    \(\Psi'=\forall X\exists Y'\,\varphi'\) berechnen, sodass jede Klausel von
    \(\varphi'\) entweder nur positive oder nur negative Literale enthält und
    \(\Psi\) genau dann gültig ist, wenn \(\Psi'\) gültig ist.
\end{Lemma}

\begin{proof}
    Für jede Variable \(v\in X\cup Y\) sei \(v'\) eine frische Variable und
    \(Y':=Y\cup\{v'\mid v\in X\cup Y\}\). Die Formel \(\varphi'\) entstehe aus
    \(\varphi\), indem jedes negative Literal \(\neg v\) durch \(v'\) ersetzt
    wird; zusätzlich nehmen wir für jedes \(v\in X\cup Y\) die beiden Klauseln
    \(
        (v\lor v')
    \) und \(
        (\neg v\lor\neg v')
    \)
    auf. Alle aus \(\varphi\) entstandenen Klauseln sind rein positiv, die
    ersten Zusatzklauseln ebenfalls, die zweiten rein negativ. Also ist
    \(\varphi'\) monoton. Die beiden Zusatzklauseln sind zusammen äquivalent zu
    \(v'\leftrightarrow\neg v\).

    Sei \(\Psi\) gültig und \(\alpha\) eine Belegung von \(X\). Wähle \(\beta\)
    mit \(\varphi(\alpha,\beta)=\top\) und setze \(v'\) auf den zu \(v\)
    komplementären Wert. Die Zusatzklauseln sind dann erfüllt. Ist eine
    ursprüngliche Klausel \(C\) durch ein positives Literal \(v\) erfüllt, so
    auch ihr Bild; ist sie durch \(\neg v\) erfüllt, so ist \(v\) falsch, also
    \(v'\) wahr, und das Bild ist ebenfalls erfüllt. Somit ist \(\Psi'\) gültig.

    Sei umgekehrt \(\Psi'\) gültig und \(\alpha\) eine Belegung von \(X\). Es
    existiert eine Belegung von \(Y'\), die \(\varphi'\) erfüllt; wegen der
    Zusatzklauseln gilt dort \(v'=\neg v\) für alle \(v\in X\cup Y\). Sei
    \(\beta\) die Einschränkung auf \(Y\). Jede ursprüngliche Klausel \(C\) hat
    ein erfülltes Bildliteral; ist dieses ein \(v\), so erfüllt \(v\) auch
    \(C\); ist es ein \(v'\), so ist \(v\) falsch und \(\neg v\in C\) erfüllt.
    Also gilt \(\varphi(\alpha,\beta)=\top\) und \(\Psi\) ist gültig.

    Die Konstruktion verdoppelt die Variablenzahl und fügt \(2|X\cup Y|\)
    Klauseln hinzu, ist also polynomiell.
\end{proof}

Damit ist \textsc{Monotone \(\forall\exists\)-SAT} \(\Pi_2^\text{P}\)-hart,
sobald \(\forall\exists\)-SAT es ist, und
Proposition~\ref{prop:bip-st-ind-hard} steht auf einer belegten Grundlage. Wird
zusätzlich Klauselbreite \(3\) gewünscht, füllt man die zweielementigen
Zusatzklauseln durch Wiederholung eines Literals auf; die Monotonie bleibt
erhalten. Alle übrigen Kritikpunkte sind redaktioneller Natur und durch die
Fassung in den Abschnitten~\ref{sec:bip-analyse-struktur}
bis~\ref{sec:bip-analyse-bipartit} bereits ausgeräumt.

\section{Fazit}
\label{sec:bip-analyse-fazit}

Die Aussage von Proposition~\ref{prop:bip-st-ind-hard} ist richtig, und der
Kern der Konstruktion --- die Aufspaltung des Formelarguments in \(\varphi^{+}\)
und \(\varphi^{-}\) sowie die Steuerung durch die frischen Variablen \(p\) und
\(n\) --- ist elegant und trägt die Bipartitheit. Der Beweis sollte jedoch um
Lemma~\ref{lem:bip-qbf-aequiv}, Lemma~\ref{lem:bip-struktur} und
Lemma~\ref{lem:monotone-normalform} ergänzt, in den Rollen von
\(\lambda_1,\lambda_2\) vereinheitlicht und um den Zugehörigkeitsverweis
ergänzt werden. Danach ist er lückenlos.

\chapter{Symmetrische irreflexive AFs}
\label{ch:sym-analyse}

\section{Fragestellung}
\label{sec:sym-frage}

Abschnitt~\ref{sec:sym-afs} hat symmetrische AFs bereits als Extremfall hoher
Angriffsdichte betrachtet und für vollständig gerichtete AFs gezeigt, dass
\textit{stable}-Unabhängigkeit dort stets verletzt ist
(Proposition~\ref{prop:cdg-afs-not-st-ci}). Offen blieb die
komplexitätstheoretische Einordnung der gesamten Klasse. Dieses Kapitel
schließt diese Lücke.

Der Ausgangsverdacht ist, dass die Klasse einfach sein müsste. Symmetrische
irreflexive AFs sind kohärent \cite{Coste:2005:symmetric-af-irflx-coherent},
besitzen stets \textit{stable}-Labellings und ihre klassischen
Entscheidungsprobleme --- Existenz, Verifikation, kredulöse und skeptische
Akzeptanz --- sind sämtlich in Polynomialzeit lösbar
(Proposition~\ref{prop:sym-klassisch-p}). Für alle bisher betrachteten
Graphklassen (Abschnitt~\ref{sec:acyc-afs} bis
Abschnitt~\ref{sec:odd-cycle-free-afs}) war eine solche strukturelle
Vereinfachung stets mit einer Vereinfachung der Unabhängigkeit einhergegangen
oder zumindest verträglich.

Das Hauptresultat dieses Kapitels widerlegt diesen Verdacht für die
\textit{stable}-Semantik: \(\textit{Ind}_{\mathrm{st}}\) bleibt auf
symmetrischen irreflexiven AFs \(\Pi_2^\text{P}\)-vollständig
(Theorem~\ref{th:sym-st-pi2p}). Für die \textit{preferred}-Semantik dagegen
tritt eine echte Vereinfachung ein: Das Problem fällt von
\(\Pi_3^\text{P}\)-Vollständigkeit auf \(\Pi_2^\text{P}\)-Vollständigkeit
(Korollar~\ref{cor:sym-pr-pi2p}). Ergänzend wird ein Fragment identifiziert, in
dem die Unabhängigkeit tatsächlich in Polynomialzeit entscheidbar wird
(Proposition~\ref{prop:sym-singleton-p}).

Wir betrachten für \(\sigma\in\{co,st,pr\}\) das Entscheidungsproblem
\[
\begin{array}{ll}
\textit{Ind}^{\mathrm{sym}}_{\sigma}\\
\text{Input:}  & \text{symmetrisches, irreflexives AF } F=(A,R),\
                 \text{disjunkte } X,Y,Z\subseteq A \\
\text{Output:} & \text{TRUE} \Longleftrightarrow X \bot_{\sigma} Y \mid Z.
\end{array}
\]

\section{Struktur der Labellingmengen}
\label{sec:sym-struktur}

\begin{Definition}
    \label{def:sym-af}
    Ein AF \(F=(A,R)\) heißt \emph{symmetrisch}, falls
    \((a,b)\in R\Rightarrow(b,a)\in R\) für alle \(a,b\in A\) gilt, und
    \emph{irreflexiv}, falls \((a,a)\notin R\) für alle \(a\in A\) gilt. Für
    \(a\in A\) schreiben wir \(N(a):=\{b\in A\mid (a,b)\in R\}\) für die
    \emph{Nachbarschaft} von \(a\) und \(N[a]:=N(a)\cup\{a\}\). Für
    \(S\subseteq A\) sei \(N(S):=\bigcup_{a\in S}N(a)\).
\end{Definition}

In einem symmetrischen irreflexiven AF ist \(R\) die Kantenrelation eines
schlichten ungerichteten Graphen. Eine Menge \(S\subseteq A\) ist genau dann
konfliktfrei, wenn sie in diesem Graphen unabhängig ist, und wir nennen \(S\)
\emph{naiv}, falls \(S\) \(\subseteq\)-maximal konfliktfrei ist. Für
\(S\subseteq A\) bezeichne \(\lambda_S\) das Labelling mit
\(\mathit{in}(\lambda_S)=S\) und \(\mathit{out}(\lambda_S)=A\setminus S\).

\begin{Lemma}
    \label{lem:sym-cf-adm}
    Sei \(F=(A,R)\) symmetrisch und irreflexiv. Dann ist jede konfliktfreie
    Menge \(S\subseteq A\) admissibel.
\end{Lemma}

\begin{proof}
    Sei \(a\in S\) und sei \(b\) ein Angreifer von \(a\), also \((b,a)\in R\).
    Wegen der Konfliktfreiheit von \(S\) gilt \(b\notin S\). Aus der Symmetrie
    folgt \((a,b)\in R\), also greift \(a\in S\) den Angreifer \(b\) an. Damit
    verteidigt \(S\) jedes seiner Elemente.
\end{proof}

\begin{Proposition}
    \label{prop:sym-naiv-st-pr}
    Sei \(F=(A,R)\) symmetrisch und irreflexiv. Dann gilt
    \[
        \Lambda_{\mathrm{st}}(F)
        =\Lambda_{\mathrm{pr}}(F)
        =\{\lambda_S\mid S\subseteq A\text{ naiv}\}.
    \]
    Insbesondere ist \(\Lambda_{\mathrm{st}}(F)\neq\emptyset\).
\end{Proposition}

\begin{proof}
    Sei \(S\) naiv und \(a\in A\setminus S\). Dann ist \(S\cup\{a\}\) nicht
    konfliktfrei. Da \(S\) konfliktfrei und \(F\) irreflexiv ist, muss der
    Konflikt zwischen \(a\) und einem \(b\in S\) bestehen; wegen der Symmetrie
    gilt \((b,a)\in R\). Also greift \(S\) jedes Argument außerhalb von \(S\)
    an, das heißt \(\lambda_S\) ist ein \textit{stable}-Labelling.

    Umgekehrt ist jede \textit{stable}-Extension konfliktfrei und maximal
    konfliktfrei, denn ein hinzunehmbares Argument wäre unangegriffen. Also
    stimmen die \textit{stable}-Extensionen genau mit den naiven Mengen
    überein.

    Für die \textit{preferred}-Semantik gilt allgemein
    \(\Lambda_{\mathrm{st}}(F)\subseteq\Lambda_{\mathrm{pr}}(F)\). Sei
    umgekehrt \(E\) eine \textit{preferred}-Extension, also
    \(\subseteq\)-maximal admissibel. Wäre \(E\) nicht maximal konfliktfrei, so
    gäbe es eine konfliktfreie Menge \(E\subsetneq T\); nach
    Lemma~\ref{lem:sym-cf-adm} wäre \(T\) admissibel, im Widerspruch zur
    Maximalität von \(E\). Also ist \(E\) naiv und damit \textit{stable}. Da
    eine \textit{stable}-Extension \(E\) die Labellingdarstellung
    \(\mathit{in}=E\), \(\mathit{out}=A\setminus E\), \(\mathit{und}=\emptyset\)
    besitzt und dies zugleich das zu \(E\) gehörige
    \textit{preferred}-Labelling ist, stimmen auch die Labellingmengen überein.

    Da \(A\) endlich ist, besitzt jede konfliktfreie Menge eine maximale
    Obermenge; insbesondere existiert mindestens eine naive Menge.
\end{proof}

\begin{Corollary}
    \label{cor:sym-extend}
    Sei \(F=(A,R)\) symmetrisch und irreflexiv und sei \(T\subseteq A\)
    konfliktfrei. Dann existiert ein \(\lambda\in\Lambda_{\mathrm{st}}(F)\) mit
    \(T\subseteq\mathit{in}(\lambda)\). Ist zusätzlich \(a\in A\setminus T\) und
    besitzt \(a\) einen Nachbarn in \(T\), so gilt
    \(\lambda(a)=\mathit{out}\) für jedes solche \(\lambda\).
\end{Corollary}

\begin{proof}
    Man erweitere \(T\) zu einer naiven Menge \(S\) und wende
    Proposition~\ref{prop:sym-naiv-st-pr} an. Besitzt \(a\) einen Nachbarn in
    \(T\subseteq S\), so verhindert die Konfliktfreiheit \(a\in S\).
\end{proof}

Korollar~\ref{cor:sym-extend} ist das Arbeitspferd aller folgenden Beweise: Es
erlaubt, ein Labelling durch Angabe eines konfliktfreien \enquote{Kerns} zu
spezifizieren, sofern jedes Argument, das ausgeschlossen werden soll, einen
Nachbarn im Kern besitzt.

\begin{Proposition}
    \label{prop:sym-klassisch-p}
    Sei \(F=(A,R)\) symmetrisch und irreflexiv und sei \(\sigma\in\{st,pr\}\).
    Dann gilt:
    \begin{enumerate}
        \item \(\Lambda_\sigma(F)\neq\emptyset\).
        \item Die Verifikation, ob ein gegebenes Labelling in
        \(\Lambda_\sigma(F)\) liegt, ist in Polynomialzeit möglich.
        \item Jedes \(a\in A\) ist kredulös akzeptiert.
        \item Ein Argument \(a\in A\) ist genau dann skeptisch akzeptiert, wenn
        \(N(a)=\emptyset\).
    \end{enumerate}
\end{Proposition}

\begin{proof}
    (1) und (2) folgen aus Proposition~\ref{prop:sym-naiv-st-pr}: Zu prüfen ist
    lediglich, ob \(\mathit{in}(\lambda)\) konfliktfrei ist und ob jedes
    Argument außerhalb einen Nachbarn darin besitzt; beides ist in
    \(\mathcal{O}(|A|+|R|)\) möglich.

    (3) Wegen der Irreflexivität ist \(\{a\}\) konfliktfrei; nach
    Korollar~\ref{cor:sym-extend} existiert ein \(\lambda\in\Lambda_\sigma(F)\)
    mit \(\lambda(a)=\mathit{in}\).

    (4) Ist \(N(a)=\emptyset\), so ist \(a\) in jeder maximal konfliktfreien
    Menge enthalten, denn \(a\) steht mit keinem Argument in Konflikt. Ist
    dagegen \(b\in N(a)\), so liefert Korollar~\ref{cor:sym-extend} ein
    \(\lambda\) mit \(\lambda(b)=\mathit{in}\) und folglich
    \(\lambda(a)=\mathit{out}\).
\end{proof}

Proposition~\ref{prop:sym-klassisch-p} formalisiert den Befund von
\textcite{Coste:2005:symmetric-af-irflx-coherent}, dass in symmetrischen
irreflexiven AFs nur zwei sehr einfache und effizient prüfbare Formen der
Akzeptanz auftreten. Nach diesem Bild müsste auch die bedingte Unabhängigkeit
leicht sein: Sie ist nach Definition~\ref{def:BU:AFs} eine Aussage über
\(\Lambda_\sigma(F)\), und alle punktweisen Fragen über \(\Lambda_\sigma(F)\)
sind trivial.

\section{Die eigentliche Schwierigkeit: exakte Realisierbarkeit}
\label{sec:sym-realisierbarkeit}

Die bedingte Unabhängigkeit fragt nicht nach dem Label eines einzelnen
Arguments, sondern danach, ob ein \emph{vollständig vorgeschriebenes
Teillabelling} realisiert werden kann. Genau hier bricht die Einfachheit
zusammen.

\begin{Definition}
    \label{def:sym-real}
    Das Realisierbarkeitsproblem \(\textit{Real}^{\mathrm{sym}}_{\mathrm{st}}\)
    ist wie folgt definiert:
    \[
    \begin{array}{ll}
    \text{Input:}  & \text{symmetrisches, irreflexives AF } F=(A,R),\
                     W\subseteq A,\ \eta:W\to\{\mathit{in},\mathit{out}\}\\
    \text{Output:} & \text{TRUE} \Longleftrightarrow
                     \exists\lambda\in\Lambda_{\mathrm{st}}(F):\lambda|_W=\eta.
    \end{array}
    \]
\end{Definition}

Der Unterschied zur kredulösen Akzeptanz ist der Zwang, Argumente auch
\emph{ausschließen} zu müssen: Ein Argument \(a\) mit \(\eta(a)=\mathit{out}\)
verlangt einen Nachbarn im \(\mathit{in}\)-Teil, und mehrere solche Forderungen
müssen gleichzeitig durch eine einzige konfliktfreie Menge erfüllt werden.

\begin{Proposition}
    \label{prop:sym-real-npc}
    \(\textit{Real}^{\mathrm{sym}}_{\mathrm{st}}\) ist NP-vollständig.
\end{Proposition}

\begin{proof}
    Zugehörigkeit: Man rät \(\lambda\) und verifiziert in Polynomialzeit
    (Proposition~\ref{prop:sym-klassisch-p}(2)), dass
    \(\lambda\in\Lambda_{\mathrm{st}}(F)\) und \(\lambda|_W=\eta\).

    Härte: Wir reduzieren von 3-SAT. Sei \(\varphi=\bigwedge_{k=1}^{m}C_k\) eine
    KNF-Formel über den Variablen \(y_1,\dots,y_l\). Wir definieren ein
    symmetrisches irreflexives AF \(F_\varphi=(A,R)\) durch
    \[
        A:=\{w_j,\overline{w_j}\mid 1\leq j\leq l\}\cup\{c_k\mid 1\leq k\leq m\}
    \]
    und die (symmetrisch abgeschlossene) Angriffsrelation, die genau die
    folgenden Konflikte enthält:
    \begin{itemize}
        \item \(\{w_j,\overline{w_j}\}\) für jedes \(j\),
        \item \(\{w_j,c_k\}\), falls \(y_j\) positiv in \(C_k\) vorkommt,
        \item \(\{\overline{w_j},c_k\}\), falls \(y_j\) negativ in \(C_k\)
        vorkommt.
    \end{itemize}
    Wir setzen \(W:=\{c_1,\dots,c_m\}\) und \(\eta\equiv\mathit{out}\).

    Sei \(\lambda\in\Lambda_{\mathrm{st}}(F_\varphi)\) mit \(\lambda|_W=\eta\)
    und \(S:=\mathit{in}(\lambda)\). Da \(S\) konfliktfrei ist, gilt nicht
    gleichzeitig \(w_j\in S\) und \(\overline{w_j}\in S\); wir definieren
    \(\beta(y_j):=\top\) genau dann, wenn \(w_j\in S\). Nach
    Proposition~\ref{prop:sym-naiv-st-pr} ist \(S\) naiv, also besitzt jedes
    \(c_k\notin S\) einen Nachbarn in \(S\). Die Nachbarn von \(c_k\) sind
    ausschließlich Literalargumente, also enthält \(S\) ein Literalargument von
    \(C_k\). Ist dies \(w_j\), so gilt \(\beta(y_j)=\top\) und \(y_j\in C_k\);
    ist es \(\overline{w_j}\), so ist \(w_j\notin S\), also
    \(\beta(y_j)=\bot\) und \(\neg y_j\in C_k\). In beiden Fällen ist \(C_k\)
    unter \(\beta\) erfüllt. Somit ist \(\varphi\) erfüllbar.

    Sei umgekehrt \(\beta\) eine erfüllende Belegung und
    \(T:=\{w_j\mid\beta(y_j)=\top\}\cup\{\overline{w_j}\mid\beta(y_j)=\bot\}\).
    \(T\) ist konfliktfrei, denn \(T\) enthält aus jedem Paar genau ein
    Argument und die Paare stehen untereinander nicht in Konflikt. Jedes
    \(c_k\) besitzt ein unter \(\beta\) wahres Literal, dessen Literalargument
    in \(T\) liegt. Nach Korollar~\ref{cor:sym-extend} existiert ein
    \(\lambda\in\Lambda_{\mathrm{st}}(F_\varphi)\) mit
    \(T\subseteq\mathit{in}(\lambda)\), und alle \(c_k\) tragen dort
    \(\mathit{out}\). Also gilt \(\lambda|_W=\eta\).

    Die Konstruktion ist offensichtlich polynomiell.
\end{proof}

Proposition~\ref{prop:sym-real-npc} erklärt, warum die Tractability der
klassischen Aufgaben nicht auf die Unabhängigkeit durchschlägt: Der innere
Existenzquantor aus Definition~\ref{def:BU:AFs} ist bereits auf dieser Klasse
NP-vollständig. Was noch fehlt, ist ein Mechanismus, der den äußeren
Allquantor über \(\lambda_1,\lambda_2\) in einen echten \(\forall\)-Quantor
über Belegungen verwandelt. Dieser Mechanismus ist der Gegenstand des nächsten
Abschnitts.

\section{Hauptresultat}
\label{sec:sym-hauptresultat}

Die zentrale konstruktive Schwierigkeit in symmetrischen AFs ist, dass jeder
Angriff in beide Richtungen wirkt. Zwei Konsequenzen sind dabei entscheidend.

Erstens lässt sich die für Reduktionen übliche Invariante \enquote{aus jedem
Paar \(v,\overline{v}\) ist genau eines \(\mathit{in}\)} nicht mehr durch einen
bloßen Zweierzyklus erzwingen: Sobald \(v\) oder \(\overline{v}\) weitere
Nachbarn besitzen, können in einer naiven Menge beide fehlen, weil sie von
anderer Seite dominiert werden. Wir lösen das, indem wir Wert- und Signalknoten
trennen und zu einem Viererkreis verbinden.

Zweitens lässt sich die in gerichteten Konstruktionen übliche Aggregation
\enquote{\(\varphi\) ist \(\mathit{in}\), also sind alle Klauselargumente
\(\mathit{out}\), also ist jede Klausel erfüllt} nicht nachbauen: Ein
Aggregationsargument \(\varphi\), das alle Klauselargumente angreift,
dominiert diese selbst und macht damit jede Forderung an die Literale
gegenstandslos. Wir verzichten deshalb auf ein Aggregationsargument und
schreiben die Klauselargumente stattdessen \emph{direkt} als
\(\mathit{out}\)-Teil der Unabhängigkeitsanfrage vor. Damit die Menge der so
vorschreibbaren Muster klein und beherrschbar bleibt, verbinden wir die
Klauselargumente zu einer Clique.

\begin{Definition}
    \label{def:sym-konstruktion}
    Sei \(\Psi=\forall X\exists Y\,\varphi\) mit
    \(X=\{x_1,\dots,x_n\}\), \(Y=\{y_1,\dots,y_l\}\) und
    \(\varphi=\bigwedge_{k=1}^{m}C_k\) in KNF, \(m\geq 1\). Wir setzen
    \(X^{+}:=X\cup\{x_0\}\) mit einer frischen universellen Variablen \(x_0\)
    und
    \[
        \widehat{\varphi}:=\bigwedge_{k=1}^{m}\bigl(x_0\lor C_k\bigr).
    \]
    Das AF \(G_\Psi=(A_\Psi,R_\Psi)\) besitzt die Argumentmenge
    \[
    \begin{aligned}
        A_\Psi:={}&
        \{u_i,\overline{u_i},s_i,\overline{s_i}\mid x_i\in X^{+}\}
        \;\cup\;
        \{w_j,\overline{w_j}\mid y_j\in Y\}
        \;\cup\;
        \{c_1,\dots,c_m\},
    \end{aligned}
    \]
    und \(R_\Psi\) ist die symmetrische Hülle genau der folgenden Konflikte:
    \begin{enumerate}
        \item[(K1)] \(\{u_i,\overline{u_i}\}\),
        \(\{u_i,\overline{s_i}\}\),
        \(\{\overline{u_i},s_i\}\),
        \(\{s_i,\overline{s_i}\}\)
        \quad für alle \(x_i\in X^{+}\),
        \item[(K2)] \(\{w_j,\overline{w_j}\}\) \quad für alle \(y_j\in Y\),
        \item[(K3)] \(\{s_i,c_k\}\), falls \(x_i\) positiv in \(C_k\) vorkommt,
        und \(\{\overline{s_i},c_k\}\), falls \(x_i\) negativ in \(C_k\)
        vorkommt \quad(\(1\leq i\leq n\)); zusätzlich \(\{s_0,c_k\}\) für alle
        \(k\),
        \item[(K4)] \(\{w_j,c_k\}\), falls \(y_j\) positiv in \(C_k\) vorkommt,
        und \(\{\overline{w_j},c_k\}\), falls \(y_j\) negativ in \(C_k\)
        vorkommt,
        \item[(K5)] \(\{c_k,c_{k'}\}\) für alle \(k\neq k'\).
    \end{enumerate}
    Als Unabhängigkeitsanfrage setzen wir
    \[
        \mathcal{X}:=\{c_1,\dots,c_m\},\qquad
        \mathcal{Y}:=\{u_i,\overline{u_i}\mid x_i\in X^{+}\},\qquad
        Z:=\emptyset.
    \]
\end{Definition}

Das AF \(G_\Psi\) ist nach Konstruktion symmetrisch und irreflexiv, besitzt
\(4(n+1)+2l+m\) Argumente und ist in Polynomialzeit berechenbar. Die Blöcke
(K1) bilden für jede universelle Variable einen Viererkreis
\(u_i-\overline{u_i}-s_i-\overline{s_i}-u_i\); die Wertknoten
\(u_i,\overline{u_i}\) besitzen \emph{keine} weiteren Nachbarn, während die
Signalknoten \(s_i,\overline{s_i}\) die Verbindung zu den Klauselargumenten
herstellen.

\begin{Lemma}[Wertknoten]
    \label{lem:sym-wertknoten}
    Sei \(\lambda\in\Lambda_{\mathrm{st}}(G_\Psi)\) und
    \(S:=\mathit{in}(\lambda)\). Dann gilt für jedes \(x_i\in X^{+}\):
    \(|S\cap\{u_i,\overline{u_i}\}|=1\).
\end{Lemma}

\begin{proof}
    Wegen \(\{u_i,\overline{u_i}\}\in R_\Psi\) und der Konfliktfreiheit von
    \(S\) liegen nicht beide in \(S\). Seien nun beide nicht in \(S\). Nach
    Proposition~\ref{prop:sym-naiv-st-pr} ist \(S\) naiv, also besitzt
    \(u_i\) einen Nachbarn in \(S\). Die Nachbarn von \(u_i\) sind nach (K1)
    genau \(\overline{u_i}\) und \(\overline{s_i}\); folglich gilt
    \(\overline{s_i}\in S\). Analog besitzt \(\overline{u_i}\) die Nachbarn
    \(u_i\) und \(s_i\), also \(s_i\in S\). Wegen
    \(\{s_i,\overline{s_i}\}\in R_\Psi\) widerspricht das der
    Konfliktfreiheit von \(S\).
\end{proof}

Für \(\lambda\in\Lambda_{\mathrm{st}}(G_\Psi)\) mit
\(S=\mathit{in}(\lambda)\) definieren wir daher die \emph{induzierte Belegung}
\[
    \alpha_\lambda:X^{+}\to\{\top,\bot\},\quad
    \alpha_\lambda(x_i)=\top:\iff u_i\in S,
\]
sowie
\[
    \beta_\lambda:Y\to\{\top,\bot\},\quad
    \beta_\lambda(y_j)=\top:\iff w_j\in S.
\]

\begin{Lemma}[Signaltreue]
    \label{lem:sym-signaltreue}
    Sei \(\lambda\in\Lambda_{\mathrm{st}}(G_\Psi)\), \(S=\mathit{in}(\lambda)\).
    Dann gilt:
    \begin{enumerate}
        \item \(s_i\in S\Rightarrow\alpha_\lambda(x_i)=\top\) und
        \(\overline{s_i}\in S\Rightarrow\alpha_\lambda(x_i)=\bot\).
        \item \(w_j\in S\Rightarrow\beta_\lambda(y_j)=\top\) und
        \(\overline{w_j}\in S\Rightarrow\beta_\lambda(y_j)=\bot\).
    \end{enumerate}
\end{Lemma}

\begin{proof}
    (1) Aus \(s_i\in S\) und \(\{\overline{u_i},s_i\}\in R_\Psi\) folgt
    \(\overline{u_i}\notin S\), mit Lemma~\ref{lem:sym-wertknoten} also
    \(u_i\in S\) und damit \(\alpha_\lambda(x_i)=\top\). Der zweite Fall
    verläuft symmetrisch über \(\{u_i,\overline{s_i}\}\in R_\Psi\).

    (2) \(w_j\in S\) liefert \(\beta_\lambda(y_j)=\top\) definitionsgemäß; aus
    \(\overline{w_j}\in S\) folgt \(w_j\notin S\) wegen
    \(\{w_j,\overline{w_j}\}\in R_\Psi\), also \(\beta_\lambda(y_j)=\bot\).
\end{proof}

\begin{Lemma}[Realisierbare Teillabellings]
    \label{lem:sym-teile}
    Sei \(\lambda\in\Lambda_{\mathrm{st}}(G_\Psi)\) und
    \(S=\mathit{in}(\lambda)\). Dann gilt:
    \begin{enumerate}
        \item \(|S\cap\mathcal{X}|\leq 1\).
        \item Zu jedem \(k\in\{1,\dots,m\}\) und jeder Belegung
        \(\alpha:X^{+}\to\{\top,\bot\}\) existiert
        \(\lambda'\in\Lambda_{\mathrm{st}}(G_\Psi)\) mit
        \(\mathit{in}(\lambda')\cap\mathcal{X}=\{c_k\}\) und
        \(\alpha_{\lambda'}=\alpha\).
        \item Zu jeder Belegung \(\alpha:X^{+}\to\{\top,\bot\}\) existiert
        genau dann ein \(\lambda'\in\Lambda_{\mathrm{st}}(G_\Psi)\) mit
        \(\mathit{in}(\lambda')\cap\mathcal{X}=\emptyset\) und
        \(\alpha_{\lambda'}=\alpha\), wenn eine Belegung \(\beta\) von \(Y\)
        mit \(\widehat{\varphi}(\alpha,\beta)=\top\) existiert.
        \item Zu jeder Belegung \(\alpha:X^{+}\to\{\top,\bot\}\) existiert
        \(\lambda'\in\Lambda_{\mathrm{st}}(G_\Psi)\) mit
        \(\alpha_{\lambda'}=\alpha\).
    \end{enumerate}
\end{Lemma}

\begin{proof}
    Wir schreiben \(U_\alpha:=\{u_i\mid\alpha(x_i)=\top\}
    \cup\{\overline{u_i}\mid\alpha(x_i)=\bot\}\) und
    \(S_\alpha:=\{s_i\mid\alpha(x_i)=\top\}
    \cup\{\overline{s_i}\mid\alpha(x_i)=\bot\}\).
    Nach (K1) ist \(U_\alpha\cup S_\alpha\) konfliktfrei: Innerhalb eines
    Viererkreises werden nur gegenüberliegende Knoten gewählt
    (\(u_i\) mit \(s_i\) beziehungsweise \(\overline{u_i}\) mit
    \(\overline{s_i}\)), und zwischen verschiedenen Blöcken bestehen keine
    Konflikte. Ferner besitzt jedes Argument aus
    \(\{u_i,\overline{u_i}\}\setminus U_\alpha\) seinen Partner in
    \(U_\alpha\) als Nachbarn. Nach Korollar~\ref{cor:sym-extend} gilt
    daher für jede Erweiterung eines \(U_\alpha\) enthaltenden konfliktfreien
    Kerns zu einem \textit{stable}-Labelling \(\lambda'\) bereits
    \(\alpha_{\lambda'}=\alpha\).

    (1) Die Klauselargumente bilden nach (K5) eine Clique; eine konfliktfreie
    Menge enthält daher höchstens eines von ihnen.

    (2) Der Kern \(T:=U_\alpha\cup\{c_k\}\) ist konfliktfrei, da Wertknoten
    nach (K1) nur Wert- und Signalknoten desselben Blocks als Nachbarn haben
    und insbesondere zu keinem Klauselargument benachbart sind. Sei
    \(\lambda'\) eine Erweiterung gemäß Korollar~\ref{cor:sym-extend}. Dann gilt
    \(\alpha_{\lambda'}=\alpha\) nach der Vorbemerkung, und jedes \(c_{k'}\)
    mit \(k'\neq k\) besitzt den Nachbarn \(c_k\in T\), ist also
    \(\mathit{out}\). Somit ist
    \(\mathit{in}(\lambda')\cap\mathcal{X}=\{c_k\}\).

    (3) \enquote{\(\Leftarrow\)}: Sei \(\beta\) mit
    \(\widehat{\varphi}(\alpha,\beta)=\top\) und sei
    \[
        T:=U_\alpha\cup S_\alpha\cup
        \{w_j\mid\beta(y_j)=\top\}\cup\{\overline{w_j}\mid\beta(y_j)=\bot\}.
    \]
    \(T\) ist konfliktfrei: Konflikte bestehen nach (K3), (K4) nur zwischen
    Signal- beziehungsweise Literalargumenten und Klauselargumenten, und
    \(T\) enthält keine Klauselargumente; innerhalb der Blöcke wurde die
    Konfliktfreiheit bereits festgestellt. Da \(\widehat{\varphi}\) unter
    \((\alpha,\beta)\) erfüllt ist, besitzt jede Klausel \(x_0\lor C_k\) ein
    wahres Literal, dessen zugehöriges Signal- oder Literalargument nach (K3),
    (K4) in \(T\) liegt und zu \(c_k\) benachbart ist. Nach
    Korollar~\ref{cor:sym-extend} existiert ein \textit{stable}-Labelling
    \(\lambda'\) mit \(T\subseteq\mathit{in}(\lambda')\); dort sind alle
    \(c_k\) \(\mathit{out}\) und es gilt \(\alpha_{\lambda'}=\alpha\).

    \enquote{\(\Rightarrow\)}: Sei \(\lambda'\) mit
    \(\mathit{in}(\lambda')\cap\mathcal{X}=\emptyset\) und
    \(\alpha_{\lambda'}=\alpha\), und sei \(S'=\mathit{in}(\lambda')\). Da
    \(S'\) naiv ist, besitzt jedes \(c_k\) einen Nachbarn in \(S'\). Die
    Nachbarn von \(c_k\) sind nach (K3)--(K5) Signalargumente,
    Literalargumente und andere Klauselargumente; letztere liegen nicht in
    \(S'\). Also enthält \(S'\) ein Signal- oder Literalargument zu einem
    Literal von \(x_0\lor C_k\). Nach Lemma~\ref{lem:sym-signaltreue} ist
    dieses Literal unter \((\alpha,\beta_{\lambda'})\) wahr. Da \(k\) beliebig
    war, gilt \(\widehat{\varphi}(\alpha,\beta_{\lambda'})=\top\).

    (4) Man erweitere \(U_\alpha\) nach Korollar~\ref{cor:sym-extend}.
\end{proof}

\begin{Theorem}
    \label{th:sym-st-pi2p}
    \(\textit{Ind}^{\mathrm{sym}}_{\mathrm{st}}\) ist
    \(\Pi_2^\text{P}\)-vollständig. Die Härte gilt bereits für
    \(Z=\emptyset\).
\end{Theorem}

\begin{proof}
    \emph{Zugehörigkeit.} Symmetrische irreflexive AFs bilden eine Teilklasse
    aller AFs; die obere Schranke folgt daher aus
    Theorem~\ref{th:ind-st-is-pi2p-c}.

    \emph{Härte.} Sei \(\Psi=\forall X\exists Y\,\varphi\) eine Instanz des
    \(\Pi_2^\text{P}\)-vollständigen Problems \(\mathrm{QBF}^2_\forall\) und sei
    \(G_\Psi\) mit \(\mathcal{X},\mathcal{Y},Z=\emptyset\) wie in
    Definition~\ref{def:sym-konstruktion}. Die Mengen \(\mathcal{X}\),
    \(\mathcal{Y}\) und \(Z\) sind paarweise disjunkt, und \(G_\Psi\) ist
    symmetrisch und irreflexiv. Wegen \(Z=\emptyset\) ist die Voraussetzung
    \(\lambda_1|_Z=\lambda_2|_Z\) aus Definition~\ref{def:BU:AFs} für jedes
    Paar erfüllt. Wir zeigen
    \[
        \Psi\text{ ist gültig}
        \quad\Longleftrightarrow\quad
        \mathcal{X}\bot_{\mathrm{st}}\mathcal{Y}\mid\emptyset
        \text{ in }G_\Psi.
    \]

    Vorab halten wir fest, dass wegen
    \(\widehat{\varphi}=\bigwedge_k(x_0\lor C_k)\) gilt:
    \(\Psi\) ist genau dann gültig, wenn
    \(\forall X^{+}\exists Y\,\widehat{\varphi}\) gültig ist. Für
    \(\alpha(x_0)=\top\) ist \(\widehat{\varphi}\) unter jeder Belegung von
    \(Y\) erfüllt, für \(\alpha(x_0)=\bot\) gilt
    \(\widehat{\varphi}\equiv\varphi\).

    \enquote{\(\Rightarrow\)}: Sei \(\Psi\) gültig und seien
    \(\lambda_1,\lambda_2\in\Lambda_{\mathrm{st}}(G_\Psi)\). Sei
    \(\alpha:=\alpha_{\lambda_2}\); nach Lemma~\ref{lem:sym-wertknoten} ist
    \(\lambda_2|_{\mathcal{Y}}\) durch \(\alpha\) eindeutig bestimmt und
    umgekehrt. Nach Lemma~\ref{lem:sym-teile}(1) ist
    \(\mathit{in}(\lambda_1)\cap\mathcal{X}\) entweder leer oder von der Form
    \(\{c_k\}\).

    Im Fall \(\mathit{in}(\lambda_1)\cap\mathcal{X}=\{c_k\}\) liefert
    Lemma~\ref{lem:sym-teile}(2) ein \(\lambda_3\) mit
    \(\mathit{in}(\lambda_3)\cap\mathcal{X}=\{c_k\}\) und
    \(\alpha_{\lambda_3}=\alpha\), also
    \(\lambda_3|_{\mathcal{X}}=\lambda_1|_{\mathcal{X}}\) und
    \(\lambda_3|_{\mathcal{Y}}=\lambda_2|_{\mathcal{Y}}\).

    Im Fall \(\mathit{in}(\lambda_1)\cap\mathcal{X}=\emptyset\) existiert wegen
    der Gültigkeit von \(\forall X^{+}\exists Y\,\widehat{\varphi}\) eine
    Belegung \(\beta\) mit \(\widehat{\varphi}(\alpha,\beta)=\top\);
    Lemma~\ref{lem:sym-teile}(3) liefert das gesuchte \(\lambda_3\).

    Damit gilt \(\mathcal{X}\bot_{\mathrm{st}}\mathcal{Y}\mid\emptyset\).

    \enquote{\(\Leftarrow\)}: Sei \(\Psi\) nicht gültig. Dann existiert eine
    Belegung \(\alpha\) von \(X\), sodass \(\varphi(\alpha,\beta)=\bot\) für
    alle \(\beta\) gilt. Wir setzen
    \(\alpha^{+}:=\alpha\cup\{x_0\mapsto\bot\}\); dann gilt
    \(\widehat{\varphi}(\alpha^{+},\beta)=\bot\) für alle \(\beta\).

    Wähle \(\lambda_2\) nach Lemma~\ref{lem:sym-teile}(4) mit
    \(\alpha_{\lambda_2}=\alpha^{+}\). Wähle \(\lambda_1\) nach
    Lemma~\ref{lem:sym-teile}(3), angewandt auf eine beliebige Belegung
    \(\alpha'\) mit \(\alpha'(x_0)=\top\); wegen
    \(\widehat{\varphi}(\alpha',\beta)=\top\) für jedes \(\beta\) existiert ein
    solches \(\lambda_1\) mit
    \(\mathit{in}(\lambda_1)\cap\mathcal{X}=\emptyset\).

    Angenommen, es existierte \(\lambda_3\in\Lambda_{\mathrm{st}}(G_\Psi)\) mit
    \(\lambda_3|_{\mathcal{X}}=\lambda_1|_{\mathcal{X}}\) und
    \(\lambda_3|_{\mathcal{Y}}=\lambda_2|_{\mathcal{Y}}\). Dann gilt
    \(\mathit{in}(\lambda_3)\cap\mathcal{X}=\emptyset\) und
    \(\alpha_{\lambda_3}=\alpha^{+}\). Nach Lemma~\ref{lem:sym-teile}(3)
    existierte dann ein \(\beta\) mit
    \(\widehat{\varphi}(\alpha^{+},\beta)=\top\), im Widerspruch zur Wahl von
    \(\alpha\). Also gilt
    \(\mathcal{X}\not\bot_{\mathrm{st}}\mathcal{Y}\mid\emptyset\).

    Die Konstruktion ist in Polynomialzeit berechenbar, womit die Härte gezeigt
    ist.
\end{proof}

\begin{Corollary}
    \label{cor:sym-nicht-ocf}
    Theorem~\ref{th:sym-st-pi2p} folgt nicht aus
    Proposition~\ref{prop:ocf-af-ci-st-pi^2}.
\end{Corollary}

\begin{proof}
    Für \(m\geq 3\) enthält \(G_\Psi\) wegen (K5) den Dreierkreis
    \(c_1\to c_2\to c_3\to c_1\) und ist damit nicht ungeradezykelfrei. Ebenso
    ist \(G_\Psi\) wegen dieses Dreiecks nicht bipartit. Die Klasse der
    symmetrischen irreflexiven AFs ist somit weder in der Klasse der
    ungeradezykelfreien noch in der der bipartiten AFs enthalten.
\end{proof}

Bemerkenswert ist, dass symmetrische irreflexive AFs dennoch kohärent sind: Die
Kohärenz wird hier nicht über die Abwesenheit ungerader Zyklen erreicht, sondern
über die Symmetrie selbst (Lemma~\ref{lem:sym-cf-adm}). Die Zyklenfreiheit ist
also hinreichend, aber nicht notwendig für Kohärenz, und für die Komplexität
der Unabhängigkeit ist offenkundig nicht die Kohärenz entscheidend.

\section{Konsequenz für die \textit{preferred}-Semantik}
\label{sec:sym-preferred}

\begin{Lemma}[Transferlemma]
    \label{lem:sym-transfer}
    Sei \(F=(A,R)\) ein AF und seien \(\sigma,\tau\) Semantiken mit
    \(\Lambda_\sigma(F)=\Lambda_\tau(F)\). Dann gilt für alle paarweise
    disjunkten \(X,Y,Z\subseteq A\):
    \(X\bot_\sigma Y\mid Z\) genau dann, wenn \(X\bot_\tau Y\mid Z\).
\end{Lemma}

\begin{proof}
    Die definierende Bedingung aus Definition~\ref{def:BU:AFs} quantifiziert
    ausschließlich über Elemente von \(\Lambda_\sigma(F)\). Ersetzt man diese
    Menge durch die nach Voraussetzung identische Menge \(\Lambda_\tau(F)\), so
    geht die Bedingung Wort für Wort in die entsprechende Bedingung für
    \(\tau\) über.
\end{proof}

\begin{Corollary}
    \label{cor:sym-pr-pi2p}
    \(\textit{Ind}^{\mathrm{sym}}_{\mathrm{pr}}\) ist
    \(\Pi_2^\text{P}\)-vollständig.
\end{Corollary}

\begin{proof}
    Nach Proposition~\ref{prop:sym-naiv-st-pr} gilt
    \(\Lambda_{\mathrm{pr}}(F)=\Lambda_{\mathrm{st}}(F)\) für jedes
    symmetrische irreflexive AF \(F\). Mit Lemma~\ref{lem:sym-transfer} sind
    \(\textit{Ind}^{\mathrm{sym}}_{\mathrm{pr}}\) und
    \(\textit{Ind}^{\mathrm{sym}}_{\mathrm{st}}\) dieselbe Sprache; die
    Behauptung folgt aus Theorem~\ref{th:sym-st-pi2p}.
\end{proof}

Für \(\sigma=\mathrm{pr}\) ist die Klasse also \emph{echt} einfacher: Das
allgemeine Problem ist nach Theorem~\ref{th:ind-pr-is-pi3p-c}
\(\Pi_3^\text{P}\)-vollständig, auf symmetrischen irreflexiven AFs sinkt es um
eine Stufe der polynomiellen Hierarchie. Der Grund ist strukturell: Die
Maximalitätsbedingung der \textit{preferred}-Semantik, die im allgemeinen Fall
den zusätzlichen Quantorenwechsel erzwingt, kollabiert hier zu einer in
Polynomialzeit prüfbaren Bedingung (Proposition~\ref{prop:sym-klassisch-p}(2)).
Unter der üblichen Annahme \(\Pi_2^\text{P}\neq\Pi_3^\text{P}\) ist dieser
Gewinn nicht nur formal.

Für \(\sigma=\mathrm{st}\) tritt dagegen kein Gewinn ein. Die Erklärung liefert
Proposition~\ref{prop:sym-real-npc}: Die Symmetrie vereinfacht alle
\emph{punktweisen} Fragen über \(\Lambda_{\mathrm{st}}(F)\), lässt aber die
\emph{simultane} Realisierbarkeit vorgeschriebener Teillabellings NP-hart. Da
die bedingte Unabhängigkeit genau von dieser simultanen Realisierbarkeit
handelt, überträgt sich die Vereinfachung nicht.

\section{Ein in Polynomialzeit entscheidbares Fragment}
\label{sec:sym-fragment}

Die Härte in Theorem~\ref{th:sym-st-pi2p} benötigt eine unbeschränkt große
Menge \(\mathcal{Y}\). Beschränkt man beide Mengen auf einzelne Argumente und
setzt \(Z=\emptyset\), so wird das Problem tatsächlich einfach. Damit ist die
Klasse in einem präzisen Sinn doch einfacher als der allgemeine Fall, in dem
schon die zugrundeliegenden Realisierbarkeitsfragen NP-vollständig sind.

\begin{Proposition}
    \label{prop:sym-singleton-p}
    Sei \(F=(A,R)\) symmetrisch und irreflexiv und seien \(x,y\in A\) mit
    \(x\neq y\). Dann ist \(\{x\}\bot_{\mathrm{st}}\{y\}\mid\emptyset\) in Zeit
    \(\mathcal{O}(|A|^2+|A|\cdot|R|)\) entscheidbar.
\end{Proposition}

\begin{proof}
    Wegen \(Z=\emptyset\) ist die Bedingung \(\lambda_1|_Z=\lambda_2|_Z\)
    stets erfüllt. Setze
    \[
        \mathcal{R}:=\bigl\{(\lambda(x),\lambda(y))\;\big|\;
        \lambda\in\Lambda_{\mathrm{st}}(F)\bigr\}
        \subseteq\{\mathit{in},\mathit{out}\}^2 .
    \]
    Nach Definition~\ref{def:BU:AFs} gilt
    \(\{x\}\bot_{\mathrm{st}}\{y\}\mid\emptyset\) genau dann, wenn aus
    \((\alpha_1,\beta_1),(\alpha_2,\beta_2)\in\mathcal{R}\) stets
    \((\alpha_1,\beta_2)\in\mathcal{R}\) folgt. Da \(\mathcal{R}\) höchstens
    vier Elemente besitzt, genügt es, die vier Zugehörigkeiten zu bestimmen.
    Wir benutzen durchgehend Proposition~\ref{prop:sym-naiv-st-pr} und
    Korollar~\ref{cor:sym-extend}.

    \emph{\((\mathit{in},\mathit{in})\).} Genau dann realisierbar, wenn
    \((x,y)\notin R\): Ist \(\{x,y\}\) konfliktfrei, erweitere man es; andernfalls
    verbietet die Konfliktfreiheit beide gemeinsam.

    \emph{\((\mathit{in},\mathit{out})\).} Ist \((x,y)\in R\), so erweitere man
    \(\{x\}\); dann ist \(y\) \(\mathit{out}\). Ist \((x,y)\notin R\), so ist
    das Paar genau dann realisierbar, wenn \(N(y)\setminus N[x]\neq\emptyset\).
    Existiert nämlich ein solches \(w\), so ist \(\{x,w\}\) konfliktfrei und
    jede Erweiterung enthält \(x\), aber wegen \(w\in N(y)\) nicht \(y\).
    Existiert umgekehrt \(\lambda\) mit \(\lambda(x)=\mathit{in}\) und
    \(\lambda(y)=\mathit{out}\), so besitzt \(y\) einen Nachbarn
    \(w\in\mathit{in}(\lambda)\); wegen \(x\notin N(y)\) ist \(w\neq x\), und
    wegen der Konfliktfreiheit gilt \(w\notin N(x)\).

    \emph{\((\mathit{out},\mathit{in})\).} Symmetrisch mit vertauschten Rollen.

    \emph{\((\mathit{out},\mathit{out})\).} Genau dann realisierbar, wenn
    \(w_1\in N(x)\setminus\{y\}\) und \(w_2\in N(y)\setminus\{x\}\) mit
    \(w_1=w_2\) oder \((w_1,w_2)\notin R\) existieren. Sind solche \(w_1,w_2\)
    gegeben, so ist \(\{w_1,w_2\}\) konfliktfrei, und jede Erweiterung schließt
    \(x\) und \(y\) aus. Existiert umgekehrt \(\lambda\) mit
    \(\lambda(x)=\lambda(y)=\mathit{out}\), so besitzen \(x\) und \(y\) je einen
    Nachbarn \(w_1,w_2\in\mathit{in}(\lambda)\); diese liegen wegen
    \(\lambda(x)=\lambda(y)=\mathit{out}\) nicht in \(\{x,y\}\) und stehen
    untereinander nicht in Konflikt.

    Alle vier Tests sind durch Absuchen von \(N(x)\), \(N(y)\) beziehungsweise
    aller Paare \((w_1,w_2)\in N(x)\times N(y)\) in der angegebenen Zeit
    durchführbar. Anschließend prüft man die Rechteckbedingung auf der höchstens
    vierelementigen Menge \(\mathcal{R}\) in konstanter Zeit.
\end{proof}

Im allgemeinen Fall sind bereits die vier Realisierbarkeitstests
NP-vollständig, da \((\mathit{in},\cdot)\) die kredulöse Akzeptanz unter
\textit{stable}-Semantik enthält. Proposition~\ref{prop:sym-singleton-p} ist
also eine echte Vereinfachung; sie zeigt zugleich, dass die Härte aus
Theorem~\ref{th:sym-st-pi2p} nicht aus einzelnen Argumentpaaren stammt, sondern
aus der \emph{gemeinsamen} Vorschrift vieler Labels.

Als Konsistenzprüfung sei angemerkt, dass
Proposition~\ref{prop:cdg-afs-not-st-ci} ein Spezialfall dieser Analyse ist: Ein
vollständig gerichtetes AF ist symmetrisch und irreflexiv, seine naiven Mengen
sind genau die einelementigen Teilmengen von \(A\), und für \(x\neq y\) ist
\((\mathit{in},\mathit{in})\) wegen \((x,y)\in R\) nicht realisierbar, während
\((\mathit{in},\mathit{out})\) und \((\mathit{out},\mathit{in})\) realisierbar
sind. Die Rechteckbedingung ist damit verletzt.

\section{Die \textit{complete}-Semantik bleibt offen}
\label{sec:sym-complete}

Für \(\sigma=\mathrm{co}\) lässt sich die Konstruktion aus
Definition~\ref{def:sym-konstruktion} nicht übernehmen. Wir halten zunächst die
Struktur der \textit{complete}-Labellings fest.

\begin{Proposition}
    \label{prop:sym-co-struktur}
    Sei \(F=(A,R)\) symmetrisch und irreflexiv und sei \(\lambda\) ein
    Labelling mit \(S:=\mathit{in}(\lambda)\). Dann ist \(\lambda\) genau dann
    \textit{complete}, wenn \(S\) konfliktfrei ist,
    \(\mathit{out}(\lambda)=N(S)\) gilt und kein \(a\in A\setminus S\) mit
    \(N(a)\subseteq N(S)\) existiert.
\end{Proposition}

\begin{proof}
    Nach Definition~\ref{def:complete-labelling} ist \(a\) genau dann
    \(\mathit{out}\), wenn \(a\) einen \(\mathit{in}\)-Angreifer besitzt, also
    genau dann, wenn \(a\in N(S)\); das ist gleichbedeutend mit
    \(\mathit{out}(\lambda)=N(S)\), und \(S\cap N(S)=\emptyset\) ist genau die
    Konfliktfreiheit von \(S\). Ferner ist \(a\) genau dann \(\mathit{in}\),
    wenn alle Angreifer \(\mathit{out}\) sind, also genau dann, wenn
    \(N(a)\subseteq N(S)\). Für \(a\in S\) ist diese Bedingung wegen der
    Symmetrie automatisch erfüllt; die verbleibende Forderung ist, dass kein
    \(a\) außerhalb von \(S\) sie erfüllt.
\end{proof}

Insbesondere ist das \textit{grounded}-Labelling durch
\(\mathit{in}(\lambda_{\mathrm{gr}})=\{a\in A\mid N(a)=\emptyset\}\) gegeben:
Isolierte Argumente sind unangegriffen, und ein nicht isoliertes Argument \(a\)
besitzt einen Nachbarn, der von einem isolierten Argument nicht angegriffen
werden kann.

Damit scheitert die Reduktion aus Abschnitt~\ref{sec:sym-hauptresultat} an
Lemma~\ref{lem:sym-wertknoten}: Wählt man \(S=\emptyset\), so ist
\(N(S)=\emptyset\), und nach Proposition~\ref{prop:sym-co-struktur} ist das
Labelling, das alle nicht isolierten Argumente mit \(\mathit{und}\) versieht,
\textit{complete}. In diesem Labelling tragen sämtliche Wertknoten
\(u_i,\overline{u_i}\) das Label \(\mathit{und}\); die Zuordnung
\enquote{\(\lambda_2\) kodiert eine totale Belegung} bricht zusammen, und der
äußere Allquantor läuft über wesentlich mehr Labellings als über die
Belegungen von \(X^{+}\). Da zusätzliche \(\lambda_2\) zusätzliche
Kombinationsforderungen erzeugen, kann die Richtung
\enquote{\(\Psi\) gültig \(\Rightarrow\) unabhängig} verletzt werden.

In gerichteten Konstruktionen wird dieses Problem durch Detektorargumente
gelöst, die partielle Labellings erkennen und die eigentliche Formelauswertung
neutralisieren --- so etwa die Argumente \(d_i\), \(s\) und \(g\) im Beweis von
Proposition~\ref{prop:ocf-af-ci-co-pi^2-c}. Ein solcher Detektor benötigt einen
gerichteten Informationsfluss \(\{x_i,\overline{x_i}\}\to d_i\to s\to g\), der
sich in symmetrischen AFs nicht nachbilden lässt: Jeder Detektor griffe seine
Quelle zurück an. Wir halten daher fest:

\begin{Proposition}
    \label{prop:sym-co-obere-schranke}
    \(\textit{Ind}^{\mathrm{sym}}_{\mathrm{co}}\in\Pi_2^\text{P}\).
\end{Proposition}

\begin{proof}
    Folgt aus Theorem~\ref{th:ind-co-is-pi2p-c}, da symmetrische irreflexive
    AFs eine Teilklasse aller AFs bilden.
\end{proof}

Ob \(\textit{Ind}^{\mathrm{sym}}_{\mathrm{co}}\) auch \(\Pi_2^\text{P}\)-hart
ist, bleibt offen. Es ist die einzige der drei betrachteten Semantiken, für die
in dieser Klasse keine vollständige Einordnung erreicht wurde.

\section{Zusammenfassung}
\label{sec:sym-zusammenfassung}

\begin{table}[htbp]
    \myfloatalign
    \centering
    \renewcommand{\arraystretch}{1.15}
    \setlength{\tabcolsep}{1.0em}
    \begin{tabular}{lccc}
        \toprule
        & \(co\) & \(st\) & \(pr\) \\
        \midrule
        allgemeine AFs
            & \(\Pi_2^\text{P}\)-c
            & \(\Pi_2^\text{P}\)-c
            & \(\Pi_3^\text{P}\)-c \\
        symmetrisch, irreflexiv
            & in \(\Pi_2^\text{P}\)
            & \(\Pi_2^\text{P}\)-c
            & \(\Pi_2^\text{P}\)-c \\
        \quad davon \(|X|=|Y|=1\), \(Z=\emptyset\)
            & ---
            & in P
            & in P \\
        \bottomrule
    \end{tabular}
    \caption[Komplexität der Unabhängigkeit auf symmetrischen irreflexiven
    AFs]{Einordnung von \(\sigma\)-Unabhängigkeit auf symmetrischen
    irreflexiven AFs im Vergleich zum allgemeinen Fall. Die Einträge der
    zweiten Zeile stammen aus Proposition~\ref{prop:sym-co-obere-schranke},
    Theorem~\ref{th:sym-st-pi2p} und Korollar~\ref{cor:sym-pr-pi2p}, der
    Eintrag der dritten Zeile aus Proposition~\ref{prop:sym-singleton-p} in
    Verbindung mit Korollar~\ref{cor:sym-pr-pi2p}.}
    \label{tab:sym-komplexitaet}
\end{table}

Die Antwort auf die Ausgangsfrage ist differenziert. Symmetrische irreflexive
AFs sind für die \textit{preferred}-Semantik echt einfacher, weil dort die
Maximalitätsbedingung kollabiert und ein Quantorenwechsel entfällt. Für die
\textit{stable}-Semantik sind sie \emph{nicht} einfacher: Trotz kohärenter
Struktur, garantierter Existenz von Labellings und durchgängig polynomieller
klassischer Entscheidungsprobleme bleibt die bedingte Unabhängigkeit
\(\Pi_2^\text{P}\)-vollständig. Eine Vereinfachung tritt erst ein, wenn man
zusätzlich die Größe der Anfragemengen beschränkt.

Methodisch ist das der interessanteste Befund dieses Kapitels: Die
Tractability der klassischen Aufgaben ist kein verlässlicher Indikator für die
Komplexität der bedingten Unabhängigkeit. Maßgeblich ist stattdessen, ob
vorgeschriebene Teillabellings \emph{simultan} realisierbar sind
(Proposition~\ref{prop:sym-real-npc}) und ob der äußere Allquantor über
hinreichend viele, strukturell unterscheidbare Labellings läuft
(Lemma~\ref{lem:sym-wertknoten}). Beide Bedingungen sind in symmetrischen
irreflexiven AFs erfüllt, obwohl jede einzelne Akzeptanzfrage trivial ist.
 in thesis.tex,
% nach \chapter{Berechnungstheoretische Vorteile bestimmter Graphklassen zur Bestimmung bedingter Unabhängigkeit}
\label{ch:graphclasses}

Die folgenden Abschnitte orientieren sich an \cite{handbook__of_formal_argumentation__vol__3:2024}.

\section{Azyklische AFs}
\label{sec:acyc-afs}

Für azyklische AFs fällt jedes existierende \emph{stable}-Labelling mit dem eindeutig bestimmten \emph{grounded}-Labelling zusammen. 
Folglich besitzt ein azyklisches AF höchstens ein \emph{stable}-Labelling. 
Dadurch ist die bedingte Unabhängigkeit bezüglich der \emph{stable}-Semantik für beliebige paarweise disjunkte Argumentmengen 
stets erfüllt.

\begin{Proposition}
\label{prop:acyc-af-ci-st-top}
    Sei F=(A,R) ein azyklisches AF. Dann ist
    \(\textit{Ind}_{\mathrm{st}}(F) \) immer wahr.
\end{Proposition}

\begin{proof}
    \label{pr:acyc-af-ci-st-top-true}
    Seien \( X,Y,Z \subseteq A\) paarweise disjunkt. Besitzt \(F\) kein
    \emph{stable}-Labelling, so ist die universell quantifizierte
    Unabhängigkeitsbedingung vakuos erfüllt.

    Andernfalls existiert aufgrund der Azyklizität genau ein
    \emph{stable}-Labelling \(\lambda\), das mit dem
    \emph{grounded}-Labelling übereinstimmt. Für beliebige
    \(\lambda_1,\lambda_2\in\Lambda_{\mathrm{st}}(F)\) gilt daher
    \(\lambda_1=\lambda_2=\lambda\). Mit \(\lambda_3:=\lambda\) können
    somit die Belegungen von \(X\), \(Y\) und \(Z\) wie in der Definition
    der bedingten Unabhängigkeit gefordert kombiniert werden. Daher gilt
    \(X\bot_{\mathrm{st}}Y\mid Z\) und folglich
    \(\operatorname{Ind}_{\mathrm{st}}(F)\).
\end{proof}

\section{Geradezykelfreie AFs}
\label{sec:even-cycle-free-afs}

Für AFs ohne gerade Zykel ist allein das \emph{grounded}-Labelling ein Kandidat, ein \emph{stable}-Labelling zu sein. 
Damit existiert ebenfalls maximal ein \emph{stable}-Labelling. 

\begin{Proposition}
\label{prop:even-cycle-free-af-ci-st-top}
    Sei F=(A,R) ein AF ohne gerade Zyklen. Dann ist
    \(\textit{Ind}_{\mathrm{st}}(F) \) immer wahr.
\end{Proposition}

\begin{proof}
    Analog zu \ref{pr:acyc-af-ci-st-top-true}.
\end{proof}


\section{Ungeradezykelfreie AFs}
\label{sec:odd-cycle-free-afs}

Für die Berechnung von \emph{stable}-Labellings gibt es in dieser Graphklasse keine berechnungstheoretischen Vorteile. 


\section{Bipartite AFs}
\label{sec:bi-afs}

Als Spezialfall eines ungeradezykelfreien AFs gibt es für die Berechnung von \emph{stable}-Labellings in bipartiten AFs keine
berechnungstheoretischen Vorteile. 


\section{Symmetrische AFs}
\label{sec:sym-afs}

Beobachtung: Je mehr Angriffe im AF, desto häufiger nicht bedingt unabhängig.

.
% ****************************************************************************************************

\chapter{Analyse des Härtebeweises für bipartite AFs}
\label{ch:bip-analyse}

\section{Gegenstand und Vorgehen}
\label{sec:bip-analyse-gegenstand}

Dieses Kapitel untersucht Proposition~\ref{prop:bip-st-ind-hard} aus
Kapitel~\ref{ch:graphclasses}, also die Behauptung, dass
\textit{stable}-Unabhängigkeit auch auf bipartiten AFs
\(\Pi_2^\text{P}\)-vollständig ist. Die Untersuchung erfolgt in drei Schritten.
Zuerst wird die Konstruktion so rekonstruiert, dass ihre Semantik durch ein
explizites Strukturlemma erfasst wird (Abschnitt~\ref{sec:bip-analyse-struktur}).
Anschließend wird die Korrektheit vollständig nachgerechnet
(Abschnitt~\ref{sec:bip-analyse-korrektheit}) und die
Graphklassenzugehörigkeit verifiziert
(Abschnitt~\ref{sec:bip-analyse-bipartit}). Erst danach werden die gefundenen
Mängel gesammelt (Abschnitt~\ref{sec:bip-analyse-kritik}) und behoben
(Abschnitt~\ref{sec:bip-analyse-reparatur}).

Das Ergebnis vorweg: \emph{Die Reduktion ist inhaltlich korrekt.} Die
Beweisführung weist jedoch eine Reihe formaler Mängel auf, von denen einer
(die unbelegte Ausgangsprobleminstanz) den Beweis in seiner jetzigen Form
angreifbar macht, während die übrigen die Nachvollziehbarkeit betreffen.

\section{Rekonstruktion der Konstruktion}
\label{sec:bip-analyse-struktur}

Wir wiederholen die Konstruktion in fixierter Notation. Sei
\(\Psi=\forall X\exists Y\,\varphi\) mit
\(\varphi=\bigwedge_{C\in\mathcal{C}}C\) in konjunktiver Normalform, wobei jede
Klausel \emph{monoton} ist, also entweder ausschließlich positive Literale
(\(C\in\mathcal{C}^{+}\)) oder ausschließlich negative Literale
(\(C\in\mathcal{C}^{-}\)) enthält. Sei \(V:=X\cup Y\cup\{p,n\}\) mit zwei
frischen Variablen \(p,n\). Das AF \(F_\Psi=(A_\Psi,R_\Psi)\) ist gegeben durch
\[
\begin{aligned}
    A_{\Psi}
    ={}&
    \{v,\overline{v}\mid v\in V\}
    \cup
    \{c_C\mid C\in\mathcal{C}\}
    \cup
    \{\varphi^{+},\varphi^{-},t^{+},t^{-}\}
\end{aligned}
\]
und der in Proposition~\ref{prop:bip-st-ind-hard} angegebenen Angriffsmenge
\(R_\Psi\). Zentral ist die Beobachtung, dass die Argumente \(p\) und
\(\overline{n}\) genau so wirken, als käme das Literal \(p\) in jeder positiven
und das Literal \(\neg n\) in jeder negativen Klausel vor. Wir machen das
explizit und setzen
\[
    \widehat{\varphi}
    :=
    \bigwedge_{C\in\mathcal{C}^{+}}(p\lor C)
    \;\land\;
    \bigwedge_{C\in\mathcal{C}^{-}}(\neg n\lor C).
\]

\begin{Lemma}
    \label{lem:bip-qbf-aequiv}
    Es gilt
    \(
        \Psi\text{ ist gültig}
        \iff
        \forall (X\cup\{p,n\})\,\exists Y\,\widehat{\varphi}
        \text{ ist gültig.}
    \)
\end{Lemma}

\begin{proof}
    \enquote{\(\Rightarrow\)}: Seien \(\alpha\) eine Belegung von \(X\) und
    \(\pi,\nu\) beliebige Wahrheitswerte für \(p\) und \(n\). Da \(\Psi\)
    gültig ist, existiert \(\beta\) mit \(\varphi(\alpha,\beta)=\top\). Dann ist
    jede Klausel \(C\in\mathcal{C}\) unter \((\alpha,\beta)\) erfüllt, also erst
    recht \(p\lor C\) beziehungsweise \(\neg n\lor C\). Somit gilt
    \(\widehat{\varphi}=\top\).

    \enquote{\(\Leftarrow\)}: Sei \(\alpha\) eine Belegung von \(X\). Wir wählen
    die Belegung \(p=\bot\) und \(n=\top\). Dann gilt
    \(\widehat{\varphi}\equiv\varphi\), und die Voraussetzung liefert ein
    \(\beta\) mit \(\varphi(\alpha,\beta)=\top\).
\end{proof}

Der Beweis in Kapitel~\ref{ch:graphclasses} argumentiert an mehreren Stellen mit
Aussagen der Form \enquote{die Argumente \(t^{+}\) und \(t^{-}\) erlauben es,
die Labels von \(\varphi^{+}\) und \(\varphi^{-}\) unabhängig zu setzen}, ohne
diese zu beweisen. Das folgende Lemma stellt die vollständige Struktur von
\(\Lambda_{\mathrm{st}}(F_\Psi)\) bereit und macht alle späteren Schritte
elementar.

\begin{Lemma}[Strukturlemma]
    \label{lem:bip-struktur}
    Sei \(\lambda\) ein Labelling von \(F_\Psi\). Dann ist \(\lambda\) genau
    dann \textit{stable}, wenn es eine totale Belegung \(\tau:V\to\{\top,\bot\}\)
    gibt, sodass die folgenden vier Bedingungen gelten.
    \begin{enumerate}
        \item Für alle \(v\in V\) ist
        \(\lambda(v)=\mathit{in}\iff\tau(v)=\top\)
        und
        \(\lambda(\overline{v})=\mathit{in}\iff\tau(v)=\bot\).
        \item Für alle \(C\in\mathcal{C}^{+}\) ist
        \(\lambda(c_C)=\mathit{out}\) genau dann, wenn \(\tau\models p\lor C\);
        andernfalls \(\lambda(c_C)=\mathit{in}\). Analog ist für
        \(C\in\mathcal{C}^{-}\) genau dann \(\lambda(c_C)=\mathit{out}\), wenn
        \(\tau\models\neg n\lor C\).
        \item \(\lambda(t^{+})=\mathit{in}\iff\lambda(\varphi^{+})=\mathit{out}\)
        und \(\lambda(t^{-})=\mathit{in}\iff\lambda(\varphi^{-})=\mathit{out}\).
        \item \(\lambda(\varphi^{+})=\mathit{in}\) impliziert
        \(\lambda(c_C)=\mathit{out}\) für alle \(C\in\mathcal{C}^{+}\);
        \(\lambda(\varphi^{-})=\mathit{in}\) impliziert
        \(\lambda(c_C)=\mathit{out}\) für alle \(C\in\mathcal{C}^{-}\).
    \end{enumerate}
    Umgekehrt liefert jede Wahl von \(\tau\) und jede mit (3) und (4)
    verträgliche Wahl von \(\lambda(\varphi^{+}),\lambda(\varphi^{-})\)
    ein \textit{stable}-Labelling. Insbesondere ist
    \(\lambda(\varphi^{+})=\mathit{out}\) stets realisierbar, und
    \(\lambda(\varphi^{+})=\mathit{in}\) ist genau dann realisierbar, wenn
    \(\tau\) alle Klauseln \(p\lor C\) mit \(C\in\mathcal{C}^{+}\) erfüllt.
\end{Lemma}

\begin{proof}
    Sei \(\lambda\in\Lambda_{\mathrm{st}}(F_\Psi)\); insbesondere ist
    \(\mathit{und}(\lambda)=\emptyset\).

    Zu (1): Die einzigen Angreifer von \(v\) sind \(\overline{v}\) und
    umgekehrt. Wären beide \(\mathit{in}\), so griffe ein
    \(\mathit{in}\)-Argument ein \(\mathit{in}\)-Argument an. Wären beide
    \(\mathit{out}\), so besäße \(v\) keinen \(\mathit{in}\)-Angreifer. Also ist
    genau eines der beiden \(\mathit{in}\), und wir setzen
    \(\tau(v):=\top\) genau dann, wenn \(\lambda(v)=\mathit{in}\).

    Zu (2): Die Angreifer von \(c_C\) mit \(C\in\mathcal{C}^{+}\) sind genau
    die Argumente \(v\) mit \(v\in C\) sowie \(p\). Nach (1) ist ein solcher
    Angreifer genau dann \(\mathit{in}\), wenn das zugehörige Literal unter
    \(\tau\) wahr ist. Da \(\lambda\) \textit{stable} ist, gilt
    \(\lambda(c_C)=\mathit{out}\) genau dann, wenn mindestens ein Angreifer
    \(\mathit{in}\) ist, also genau dann, wenn \(\tau\models p\lor C\). Der Fall
    \(C\in\mathcal{C}^{-}\) verläuft analog mit den Angreifern \(\overline{v}\)
    für \(\neg v\in C\) und \(\overline{n}\).

    Zu (3): Der einzige Angreifer von \(t^{+}\) ist \(\varphi^{+}\). Also ist
    \(t^{+}\) genau dann \(\mathit{out}\), wenn \(\varphi^{+}\) \(\mathit{in}\)
    ist, und wegen \(\mathit{und}(\lambda)=\emptyset\) genau dann
    \(\mathit{in}\), wenn \(\varphi^{+}\) \(\mathit{out}\) ist.

    Zu (4): Die Angreifer von \(\varphi^{+}\) sind die \(c_C\) mit
    \(C\in\mathcal{C}^{+}\) sowie \(t^{+}\). Ist \(\varphi^{+}\)
    \(\mathit{in}\), so sind alle Angreifer \(\mathit{out}\).

    Für die Umkehrung sei \(\tau\) beliebig und seien
    \(\lambda(\varphi^{\pm})\) gemäß (3) und (4) gewählt; die übrigen Labels
    seien durch (1)--(3) festgelegt. Man verifiziert unmittelbar für jedes
    Argument die \textit{stable}-Bedingung: Argumente \(v,\overline{v}\) sind
    genau dann \(\mathit{out}\), wenn ihr Partner \(\mathit{in}\) ist;
    Klauselargumente genau dann, wenn ein Literalargument \(\mathit{in}\) ist;
    \(t^{\pm}\) genau dann, wenn \(\varphi^{\pm}\) \(\mathit{in}\) ist; und
    \(\varphi^{+}\) ist entweder \(\mathit{in}\) (dann sind nach (4) alle
    \(c_C\) und nach (3) auch \(t^{+}\) \(\mathit{out}\)) oder \(\mathit{out}\)
    (dann ist nach (3) \(t^{+}\) \(\mathit{in}\) und greift \(\varphi^{+}\) an).
    Es tritt kein \(\mathit{und}\) auf.

    Der Zusatz folgt: \(\lambda(\varphi^{+})=\mathit{out}\) ist wegen
    \(\lambda(t^{+})=\mathit{in}\) immer zulässig, und
    \(\lambda(\varphi^{+})=\mathit{in}\) ist nach (4) und (2) genau dann
    zulässig, wenn \(\tau\) jede Klausel \(p\lor C\), \(C\in\mathcal{C}^{+}\),
    erfüllt.
\end{proof}

Lemma~\ref{lem:bip-struktur} zeigt insbesondere: Die Labels von
\(\varphi^{+}\) und \(\varphi^{-}\) sind \emph{unabhängig voneinander} frei
wählbar, sobald die jeweiligen Klauselargumente \(\mathit{out}\) sind, da
\(\varphi^{+},t^{+}\) und \(\varphi^{-},t^{-}\) disjunkte Angreifermengen
besitzen. Genau diese Aussage benutzt der Originalbeweis, ohne sie zu begründen.

\section{Korrektheit der Reduktion}
\label{sec:bip-analyse-korrektheit}

\begin{Proposition}
    \label{prop:bip-korrekt}
    Es gilt
    \[
        \Psi\text{ ist gültig}
        \quad\Longleftrightarrow\quad
        \{\varphi^{+},\varphi^{-}\}\bot_{\mathrm{st}}(X\cup\{p,n\})\mid\emptyset
        \text{ in }F_\Psi.
    \]
\end{Proposition}

\begin{proof}
    Wir schreiben \(\mathcal{X}:=\{\varphi^{+},\varphi^{-}\}\) und
    \(\mathcal{Y}:=X\cup\{p,n\}\); beide Mengen sind disjunkt, und
    \(Z=\emptyset\), sodass die Bedingung \(\lambda_1|_Z=\lambda_2|_Z\) aus
    Definition~\ref{def:BU:AFs} für jedes Paar erfüllt ist. Gesucht ist also
    stets ein \(\lambda_3\in\Lambda_{\mathrm{st}}(F_\Psi)\) mit
    \(\lambda_3|_{\mathcal{X}}=\lambda_1|_{\mathcal{X}}\) und
    \(\lambda_3|_{\mathcal{Y}}=\lambda_2|_{\mathcal{Y}}\).

    \enquote{\(\Rightarrow\)}: Sei \(\Psi\) gültig und seien
    \(\lambda_1,\lambda_2\in\Lambda_{\mathrm{st}}(F_\Psi)\). Nach
    Lemma~\ref{lem:bip-struktur} induziert \(\lambda_2\) eine totale Belegung
    \(\tau_2\) von \(V\); insbesondere legt \(\lambda_2|_{\mathcal{Y}}\) die
    Werte \(\tau_2|_{X\cup\{p,n\}}\) fest. Nach Lemma~\ref{lem:bip-qbf-aequiv}
    existiert eine Belegung \(\beta\) von \(Y\), sodass
    \(\tau_3:=\tau_2|_{X\cup\{p,n\}}\cup\beta\) die Formel
    \(\widehat{\varphi}\) erfüllt. Wir wählen \(\lambda_3\) als das durch
    \(\tau_3\) und
    \(\lambda_3(\varphi^{+}):=\lambda_1(\varphi^{+})\),
    \(\lambda_3(\varphi^{-}):=\lambda_1(\varphi^{-})\)
    bestimmte Labelling. Da \(\tau_3\models\widehat{\varphi}\), sind nach
    Lemma~\ref{lem:bip-struktur}(2) sämtliche Klauselargumente
    \(\mathit{out}\); damit sind beide Wahlen für \(\varphi^{+}\) und
    \(\varphi^{-}\) mit (4) verträglich, und \(\lambda_3\) ist
    \textit{stable}. Nach Konstruktion gilt
    \(\lambda_3|_{\mathcal{X}}=\lambda_1|_{\mathcal{X}}\) und
    \(\lambda_3|_{\mathcal{Y}}=\lambda_2|_{\mathcal{Y}}\).

    \enquote{\(\Leftarrow\)}: Sei \(\Psi\) nicht gültig. Dann existiert eine
    Belegung \(\alpha\) von \(X\), sodass \(\varphi(\alpha,\beta)=\bot\) für
    jede Belegung \(\beta\) von \(Y\) gilt. Wir wählen:
    \begin{itemize}
        \item \(\lambda_1\) als ein beliebiges \textit{stable}-Labelling mit
        \(\lambda_1(\varphi^{+})=\lambda_1(\varphi^{-})=\mathit{in}\). Ein
        solches existiert: Man setze \(\tau_1(p)=\top\) und \(\tau_1(n)=\bot\)
        sowie \(\tau_1\) auf \(X\cup Y\) beliebig. Dann sind alle Klauseln
        \(p\lor C\) und \(\neg n\lor C\) erfüllt, also nach
        Lemma~\ref{lem:bip-struktur} alle Klauselargumente \(\mathit{out}\)
        und beide \(\varphi\)-Argumente \(\mathit{in}\) realisierbar.
        \item \(\lambda_2\) als ein \textit{stable}-Labelling zur Belegung
        \(\tau_2:=\alpha\cup\{p\mapsto\bot,\;n\mapsto\top\}\cup\beta_0\) mit
        beliebigem \(\beta_0\).
    \end{itemize}
    Angenommen, es gäbe ein \(\lambda_3\in\Lambda_{\mathrm{st}}(F_\Psi)\) mit
    \(\lambda_3|_{\mathcal{X}}=\lambda_1|_{\mathcal{X}}\) und
    \(\lambda_3|_{\mathcal{Y}}=\lambda_2|_{\mathcal{Y}}\). Sei \(\tau_3\) die
    zugehörige Belegung. Aus
    \(\lambda_3|_{\mathcal{Y}}=\lambda_2|_{\mathcal{Y}}\) folgt
    \(\tau_3|_X=\alpha\), \(\tau_3(p)=\bot\) und \(\tau_3(n)=\top\). Aus
    \(\lambda_3(\varphi^{+})=\lambda_3(\varphi^{-})=\mathit{in}\) folgt mit
    Lemma~\ref{lem:bip-struktur}(4) und (2), dass \(\tau_3\) sämtliche Klauseln
    \(p\lor C\) und \(\neg n\lor C\) erfüllt, also
    \(\tau_3\models\widehat{\varphi}\). Wegen \(\tau_3(p)=\bot\) und
    \(\tau_3(n)=\top\) gilt aber \(\widehat{\varphi}\equiv\varphi\) unter
    \(\tau_3\), also \(\varphi(\alpha,\tau_3|_Y)=\top\) im Widerspruch zur Wahl
    von \(\alpha\). Also existiert kein solches \(\lambda_3\), und es gilt
    \(\mathcal{X}\not\bot_{\mathrm{st}}\mathcal{Y}\mid\emptyset\).
\end{proof}

Die Konstruktion ist offensichtlich in Polynomialzeit berechenbar: Sie erzeugt
\(2|V|+|\mathcal{C}|+4\) Argumente und
\(\mathcal{O}(|V|+\sum_{C\in\mathcal{C}}|C|)\) Angriffe.

\section{Graphklassenzugehörigkeit}
\label{sec:bip-analyse-bipartit}

\begin{Proposition}
    \label{prop:bip-korrekt-bipartit}
    \(F_\Psi\) ist bipartit.
\end{Proposition}

\begin{proof}
    Wir setzen
    \[
    \begin{aligned}
        L:={}& \{v\mid v\in V\}\cup\{c_C\mid C\in\mathcal{C}^{-}\}
              \cup\{\varphi^{+},t^{-}\},\\
        R:={}& \{\overline{v}\mid v\in V\}\cup\{c_C\mid C\in\mathcal{C}^{+}\}
              \cup\{\varphi^{-},t^{+}\}.
    \end{aligned}
    \]
    Offensichtlich gilt \(L\cap R=\emptyset\) und \(L\cup R=A_\Psi\). Für jede
    Art von Angriff prüft man die Seitenwechsel-Eigenschaft:
    \((v,\overline{v})\) und \((\overline{v},v)\) verlaufen zwischen \(L\) und
    \(R\); \((v,c_C)\) mit \(C\in\mathcal{C}^{+}\) sowie \((p,c_C)\) verlaufen
    von \(L\) nach \(R\); \((\overline{v},c_C)\) mit \(C\in\mathcal{C}^{-}\)
    sowie \((\overline{n},c_C)\) verlaufen von \(R\) nach \(L\);
    \((c_C,\varphi^{+})\) mit \(C\in\mathcal{C}^{+}\) verläuft von \(R\) nach
    \(L\); \((c_C,\varphi^{-})\) mit \(C\in\mathcal{C}^{-}\) von \(L\) nach
    \(R\); \((\varphi^{+},t^{+})\) von \(L\) nach \(R\) und
    \((\varphi^{-},t^{-})\) von \(R\) nach \(L\), jeweils mit den
    Gegenangriffen. Also gilt
    \(R_\Psi\subseteq(L\times R)\cup(R\times L)\).
\end{proof}

Hier zeigt sich der eigentliche Grund für die Monotonieforderung: Nur weil
positive Klauseln ausschließlich von unquerstrichenen und negative Klauseln
ausschließlich von querstrichenen Literalargumenten angegriffen werden, lassen
sich die Klauselargumente überhaupt widerspruchsfrei auf die beiden Seiten
verteilen. Bei gemischten Klauseln entstünde ein Klauselargument mit Angreifern
aus \(L\) \emph{und} \(R\), und die Bipartition bräche zusammen.

\section{Kritikpunkte}
\label{sec:bip-analyse-kritik}

Die Reduktion ist nach den Abschnitten~\ref{sec:bip-analyse-korrektheit}
und~\ref{sec:bip-analyse-bipartit} korrekt. Der Beweis in
Kapitel~\ref{ch:graphclasses} weist jedoch die folgenden Mängel auf. Wir
ordnen sie nach abnehmender Schwere.

\begin{enumerate}
    \item \textbf{Unbelegtes Ausgangsproblem.} Der Beweis reduziert von
    \textsc{Balanced Monotone \(\forall\exists\)-3-SAT-\((1,1,2,2)\)} und
    behauptet ohne Quellenangabe und ohne Definition, dieses Problem sei
    \(\Pi_2^\text{P}\)-vollständig. Weder wird gesagt, worauf sich
    \enquote{balanced} bezieht, noch was das Tupel \((1,1,2,2)\) zählt. Das ist
    der einzige Mangel, der die \emph{Gültigkeit} der Aussage betrifft: Eine
    Reduktion von einem Problem unbekannter Komplexität beweist nichts.
    Vorkommensbeschränkungen sind zudem heikel. Nach dem Satz von Tovey ist
    jede \(k\)-CNF-Formel, in der jede Klausel \(k\) verschiedene Literale
    besitzt und jede Variable in höchstens \(k\) Klauseln vorkommt, erfüllbar;
    zu enge Vorkommensschranken können ein Problem also trivialisieren.
    \item \textbf{Die Zusatzrestriktionen werden nicht benötigt.} Die
    Konstruktion verwendet ausschließlich die Monotonie der Klauseln (siehe
    Abschnitt~\ref{sec:bip-analyse-bipartit}); weder die Ausgeglichenheit von
    \(|\mathcal{C}^{+}|\) und \(|\mathcal{C}^{-}|\) noch irgendeine
    Vorkommensschranke noch die Klauselbreite \(3\) geht in Korrektheit oder
    Bipartitheit ein. Der Beweis macht sich also ohne Not von einer stärkeren
    und schwerer zu belegenden Voraussetzung abhängig.
    Abschnitt~\ref{sec:bip-analyse-reparatur} zeigt, dass die Monotonie
    elementar und selbsttragend erreicht werden kann.
    \item \textbf{Vertauschte Rollen von \(\lambda_1\) und \(\lambda_2\).} In
    der Richtung \enquote{\(\Psi\) gültig} übernimmt \(\lambda_1\) die Labels
    von \(\varphi^{+},\varphi^{-}\) und \(\lambda_2\) die Belegung von
    \(X\cup\{p,n\}\), passend zur Reihenfolge der Mengen in
    \(\{\varphi^{+},\varphi^{-}\}\bot_{\mathrm{st}}X\cup\{p,n\}\mid\emptyset\)
    und zu Definition~\ref{def:BU:AFs}. In der Gegenrichtung sind die Rollen
    stillschweigend vertauscht: Dort trägt \(\lambda_1\) die Belegung und
    \(\lambda_2\) die \(\varphi\)-Labels. Inhaltlich ist das folgenlos, weil
    die bedingte Unabhängigkeit symmetrisch ist, formal ist es jedoch ein
    Bruch mit der eigenen Definition. Entweder man benennt konsistent um (wie
    in Proposition~\ref{prop:bip-korrekt}) oder man beruft sich explizit auf
    das Symmetrieaxiom.
    \item \textbf{Undefinierte Mengen \(A\) und \(B\).} Der Satz
    \enquote{[\ldots] das die Einschränkung von \(\lambda_1\) auf \(A\) und die
    Einschränkung von \(\lambda_2\) auf \(B\) kombiniert} verwendet zwei nie
    eingeführte Symbole. Erschwerend kommt hinzu, dass \(A\) in der gesamten
    Arbeit die Argumentmenge eines AFs bezeichnet; gemeint sind offenbar
    \(\{\varphi^{+},\varphi^{-}\}\) und \(X\cup\{p,n\}\).
    \item \textbf{Namenskollision \(L,R\).} Die Bipartitionsklassen heißen
    \(L\) und \(R\), während \(R\) beziehungsweise \(R_\Psi\) die
    Angriffsrelation bezeichnet. Die Schlussformel
    \(R_\Psi\subseteq(L\times R)\cup(R\times L)\) mischt beide Bedeutungen in
    einer Zeile. Vorzuziehen wären etwa \(A_1,A_2\) oder \(L_1,L_2\).
    \item \textbf{Fehlendes Strukturlemma.} Der Beweis benutzt mehrfach
    Eigenschaften von \(\Lambda_{\mathrm{st}}(F_\Psi)\), die er nicht beweist:
    dass in jedem \textit{stable}-Labelling genau eines von \(v,\overline{v}\)
    das Label \(\mathit{in}\) trägt; dass \(c_C\) genau dann \(\mathit{out}\)
    ist, wenn ein Literalargument \(\mathit{in}\) ist; und dass \(t^{+},t^{-}\)
    die freie und \emph{voneinander unabhängige} Wahl der \(\varphi\)-Labels
    erlauben. Lemma~\ref{lem:bip-struktur} schließt diese Lücke.
    \item \textbf{Unvollständige Fallunterscheidung über \(p\) und \(n\).} Der
    Abschnitt \enquote{Konstruktion} diskutiert nur die Kombinationen
    \((p,n)=(\mathit{out},\mathit{in})\) und
    \((p,n)=(\mathit{in},\mathit{out})\). Die ebenfalls in
    \textit{stable}-Labellings auftretenden Kombinationen
    \((\mathit{in},\mathit{in})\) und \((\mathit{out},\mathit{out})\) bleiben
    unerwähnt. Der Korrektheitsbeweis benötigt diese Fallunterscheidung
    tatsächlich nicht --- er läuft über die Formeläquivalenz aus
    Lemma~\ref{lem:bip-qbf-aequiv} ---, der Text suggeriert jedoch eine
    erschöpfende Analyse, die er nicht liefert.
    \item \textbf{Zugehörigkeit fehlt.} Die Proposition behauptet
    \(\Pi_2^\text{P}\)-\emph{Vollständigkeit}, der Beweis zeigt nur die Härte.
    Ein Verweis auf Theorem~\ref{th:ind-st-is-pi2p-c} (die obere Schranke gilt
    für alle AFs, also insbesondere für bipartite) fehlt, anders als beim
    ungeradezykelfreien Fall.
    \item \textbf{Formulierungs- und Notationsmängel.} Die Instanz wird als
    Behauptung formuliert (\enquote{wählen wir
    \(\{\varphi^{+},\varphi^{-}\}\perp_{st}\ldots\)}) statt als Anfrage; die
    Polynomialität der Konstruktion wird nicht erwähnt; \(\perp\) wird sowohl
    für die Unabhängigkeitsrelation als auch (in \enquote{\(p=\perp\)}) für den
    Wahrheitswert \emph{falsch} verwendet, während an anderer Stelle
    \(\bot_{\mathrm{st}}\) für die Unabhängigkeit steht; die Argumentmenge
    heißt hier \(\mathcal{A}_\Psi\), im ungeradezykelfreien Fall dagegen
    \(A_\Psi\); und der Satz \enquote{dann werden die Klauselargumente
    ausschließlich durch die ursprünglichen Literale kontrolliert werden}
    enthält ein doppeltes Verb.
\end{enumerate}

\begin{Corollary}
    \label{cor:bip-subsumiert-ocf}
    Proposition~\ref{prop:bip-st-ind-hard} impliziert
    Proposition~\ref{prop:ocf-af-ci-st-pi^2}.
\end{Corollary}

\begin{proof}
    Ist \(F\) bipartit, so wechselt jeder gerichtete Zyklus in jedem Schritt die
    Bipartitionsklasse und hat daher gerade Länge. Bipartite AFs sind also
    ungeradezykelfrei, und die Härte auf der Teilklasse überträgt sich auf die
    Oberklasse.
\end{proof}

Das ist ein struktureller Hinweis, kein Fehler: Der eigenständige Beweis für
ungeradezykelfreie AFs für \(\sigma=\mathrm{st}\) und \(\sigma=\mathrm{pr}\) ist
entbehrlich, sobald das bipartite Resultat gesichert ist. Umgekehrt bleibt der
gesonderte Beweis für \(\sigma=\mathrm{co}\)
(Proposition~\ref{prop:ocf-af-ci-co-pi^2-c}) nötig, denn das dort verwendete AF
ist wegen der Argumente \(d_i\), die sowohl von \(x_i\) als auch von
\(\overline{x_i}\) angegriffen werden, nicht bipartit.

Schließlich sei angemerkt, dass die Herleitung von
Proposition~\ref{prop:bip-prf-ind-hard} aus dem \textit{stable}-Fall die
Kohärenz ungeradezykelfreier AFs voraussetzt, also
\(\Lambda_{\mathrm{st}}(F)=\Lambda_{\mathrm{pr}}(F)\). Diese Tatsache stammt von
\textcite{dung:1995:af} und sollte dort zitiert werden.

\section{Reparatur: Monotonie ohne fremde Voraussetzung}
\label{sec:bip-analyse-reparatur}

Der schwerwiegendste Kritikpunkt lässt sich vollständig beheben. Statt sich auf
eine exotische Restriktion zu berufen, genügt eine elementare Normalform, die
sich in wenigen Zeilen beweisen lässt.

\begin{Lemma}[Monotone Normalform]
    \label{lem:monotone-normalform}
    Zu jeder Formel \(\Psi=\forall X\exists Y\,\varphi\) mit \(\varphi\) in KNF
    lässt sich in Polynomialzeit eine Formel
    \(\Psi'=\forall X\exists Y'\,\varphi'\) berechnen, sodass jede Klausel von
    \(\varphi'\) entweder nur positive oder nur negative Literale enthält und
    \(\Psi\) genau dann gültig ist, wenn \(\Psi'\) gültig ist.
\end{Lemma}

\begin{proof}
    Für jede Variable \(v\in X\cup Y\) sei \(v'\) eine frische Variable und
    \(Y':=Y\cup\{v'\mid v\in X\cup Y\}\). Die Formel \(\varphi'\) entstehe aus
    \(\varphi\), indem jedes negative Literal \(\neg v\) durch \(v'\) ersetzt
    wird; zusätzlich nehmen wir für jedes \(v\in X\cup Y\) die beiden Klauseln
    \(
        (v\lor v')
    \) und \(
        (\neg v\lor\neg v')
    \)
    auf. Alle aus \(\varphi\) entstandenen Klauseln sind rein positiv, die
    ersten Zusatzklauseln ebenfalls, die zweiten rein negativ. Also ist
    \(\varphi'\) monoton. Die beiden Zusatzklauseln sind zusammen äquivalent zu
    \(v'\leftrightarrow\neg v\).

    Sei \(\Psi\) gültig und \(\alpha\) eine Belegung von \(X\). Wähle \(\beta\)
    mit \(\varphi(\alpha,\beta)=\top\) und setze \(v'\) auf den zu \(v\)
    komplementären Wert. Die Zusatzklauseln sind dann erfüllt. Ist eine
    ursprüngliche Klausel \(C\) durch ein positives Literal \(v\) erfüllt, so
    auch ihr Bild; ist sie durch \(\neg v\) erfüllt, so ist \(v\) falsch, also
    \(v'\) wahr, und das Bild ist ebenfalls erfüllt. Somit ist \(\Psi'\) gültig.

    Sei umgekehrt \(\Psi'\) gültig und \(\alpha\) eine Belegung von \(X\). Es
    existiert eine Belegung von \(Y'\), die \(\varphi'\) erfüllt; wegen der
    Zusatzklauseln gilt dort \(v'=\neg v\) für alle \(v\in X\cup Y\). Sei
    \(\beta\) die Einschränkung auf \(Y\). Jede ursprüngliche Klausel \(C\) hat
    ein erfülltes Bildliteral; ist dieses ein \(v\), so erfüllt \(v\) auch
    \(C\); ist es ein \(v'\), so ist \(v\) falsch und \(\neg v\in C\) erfüllt.
    Also gilt \(\varphi(\alpha,\beta)=\top\) und \(\Psi\) ist gültig.

    Die Konstruktion verdoppelt die Variablenzahl und fügt \(2|X\cup Y|\)
    Klauseln hinzu, ist also polynomiell.
\end{proof}

Damit ist \textsc{Monotone \(\forall\exists\)-SAT} \(\Pi_2^\text{P}\)-hart,
sobald \(\forall\exists\)-SAT es ist, und
Proposition~\ref{prop:bip-st-ind-hard} steht auf einer belegten Grundlage. Wird
zusätzlich Klauselbreite \(3\) gewünscht, füllt man die zweielementigen
Zusatzklauseln durch Wiederholung eines Literals auf; die Monotonie bleibt
erhalten. Alle übrigen Kritikpunkte sind redaktioneller Natur und durch die
Fassung in den Abschnitten~\ref{sec:bip-analyse-struktur}
bis~\ref{sec:bip-analyse-bipartit} bereits ausgeräumt.

\section{Fazit}
\label{sec:bip-analyse-fazit}

Die Aussage von Proposition~\ref{prop:bip-st-ind-hard} ist richtig, und der
Kern der Konstruktion --- die Aufspaltung des Formelarguments in \(\varphi^{+}\)
und \(\varphi^{-}\) sowie die Steuerung durch die frischen Variablen \(p\) und
\(n\) --- ist elegant und trägt die Bipartitheit. Der Beweis sollte jedoch um
Lemma~\ref{lem:bip-qbf-aequiv}, Lemma~\ref{lem:bip-struktur} und
Lemma~\ref{lem:monotone-normalform} ergänzt, in den Rollen von
\(\lambda_1,\lambda_2\) vereinheitlicht und um den Zugehörigkeitsverweis
ergänzt werden. Danach ist er lückenlos.

\chapter{Symmetrische irreflexive AFs}
\label{ch:sym-analyse}

\section{Fragestellung}
\label{sec:sym-frage}

Abschnitt~\ref{sec:sym-afs} hat symmetrische AFs bereits als Extremfall hoher
Angriffsdichte betrachtet und für vollständig gerichtete AFs gezeigt, dass
\textit{stable}-Unabhängigkeit dort stets verletzt ist
(Proposition~\ref{prop:cdg-afs-not-st-ci}). Offen blieb die
komplexitätstheoretische Einordnung der gesamten Klasse. Dieses Kapitel
schließt diese Lücke.

Der Ausgangsverdacht ist, dass die Klasse einfach sein müsste. Symmetrische
irreflexive AFs sind kohärent \cite{Coste:2005:symmetric-af-irflx-coherent},
besitzen stets \textit{stable}-Labellings und ihre klassischen
Entscheidungsprobleme --- Existenz, Verifikation, kredulöse und skeptische
Akzeptanz --- sind sämtlich in Polynomialzeit lösbar
(Proposition~\ref{prop:sym-klassisch-p}). Für alle bisher betrachteten
Graphklassen (Abschnitt~\ref{sec:acyc-afs} bis
Abschnitt~\ref{sec:odd-cycle-free-afs}) war eine solche strukturelle
Vereinfachung stets mit einer Vereinfachung der Unabhängigkeit einhergegangen
oder zumindest verträglich.

Das Hauptresultat dieses Kapitels widerlegt diesen Verdacht für die
\textit{stable}-Semantik: \(\textit{Ind}_{\mathrm{st}}\) bleibt auf
symmetrischen irreflexiven AFs \(\Pi_2^\text{P}\)-vollständig
(Theorem~\ref{th:sym-st-pi2p}). Für die \textit{preferred}-Semantik dagegen
tritt eine echte Vereinfachung ein: Das Problem fällt von
\(\Pi_3^\text{P}\)-Vollständigkeit auf \(\Pi_2^\text{P}\)-Vollständigkeit
(Korollar~\ref{cor:sym-pr-pi2p}). Ergänzend wird ein Fragment identifiziert, in
dem die Unabhängigkeit tatsächlich in Polynomialzeit entscheidbar wird
(Proposition~\ref{prop:sym-singleton-p}).

Wir betrachten für \(\sigma\in\{co,st,pr\}\) das Entscheidungsproblem
\[
\begin{array}{ll}
\textit{Ind}^{\mathrm{sym}}_{\sigma}\\
\text{Input:}  & \text{symmetrisches, irreflexives AF } F=(A,R),\
                 \text{disjunkte } X,Y,Z\subseteq A \\
\text{Output:} & \text{TRUE} \Longleftrightarrow X \bot_{\sigma} Y \mid Z.
\end{array}
\]

\section{Struktur der Labellingmengen}
\label{sec:sym-struktur}

\begin{Definition}
    \label{def:sym-af}
    Ein AF \(F=(A,R)\) heißt \emph{symmetrisch}, falls
    \((a,b)\in R\Rightarrow(b,a)\in R\) für alle \(a,b\in A\) gilt, und
    \emph{irreflexiv}, falls \((a,a)\notin R\) für alle \(a\in A\) gilt. Für
    \(a\in A\) schreiben wir \(N(a):=\{b\in A\mid (a,b)\in R\}\) für die
    \emph{Nachbarschaft} von \(a\) und \(N[a]:=N(a)\cup\{a\}\). Für
    \(S\subseteq A\) sei \(N(S):=\bigcup_{a\in S}N(a)\).
\end{Definition}

In einem symmetrischen irreflexiven AF ist \(R\) die Kantenrelation eines
schlichten ungerichteten Graphen. Eine Menge \(S\subseteq A\) ist genau dann
konfliktfrei, wenn sie in diesem Graphen unabhängig ist, und wir nennen \(S\)
\emph{naiv}, falls \(S\) \(\subseteq\)-maximal konfliktfrei ist. Für
\(S\subseteq A\) bezeichne \(\lambda_S\) das Labelling mit
\(\mathit{in}(\lambda_S)=S\) und \(\mathit{out}(\lambda_S)=A\setminus S\).

\begin{Lemma}
    \label{lem:sym-cf-adm}
    Sei \(F=(A,R)\) symmetrisch und irreflexiv. Dann ist jede konfliktfreie
    Menge \(S\subseteq A\) admissibel.
\end{Lemma}

\begin{proof}
    Sei \(a\in S\) und sei \(b\) ein Angreifer von \(a\), also \((b,a)\in R\).
    Wegen der Konfliktfreiheit von \(S\) gilt \(b\notin S\). Aus der Symmetrie
    folgt \((a,b)\in R\), also greift \(a\in S\) den Angreifer \(b\) an. Damit
    verteidigt \(S\) jedes seiner Elemente.
\end{proof}

\begin{Proposition}
    \label{prop:sym-naiv-st-pr}
    Sei \(F=(A,R)\) symmetrisch und irreflexiv. Dann gilt
    \[
        \Lambda_{\mathrm{st}}(F)
        =\Lambda_{\mathrm{pr}}(F)
        =\{\lambda_S\mid S\subseteq A\text{ naiv}\}.
    \]
    Insbesondere ist \(\Lambda_{\mathrm{st}}(F)\neq\emptyset\).
\end{Proposition}

\begin{proof}
    Sei \(S\) naiv und \(a\in A\setminus S\). Dann ist \(S\cup\{a\}\) nicht
    konfliktfrei. Da \(S\) konfliktfrei und \(F\) irreflexiv ist, muss der
    Konflikt zwischen \(a\) und einem \(b\in S\) bestehen; wegen der Symmetrie
    gilt \((b,a)\in R\). Also greift \(S\) jedes Argument außerhalb von \(S\)
    an, das heißt \(\lambda_S\) ist ein \textit{stable}-Labelling.

    Umgekehrt ist jede \textit{stable}-Extension konfliktfrei und maximal
    konfliktfrei, denn ein hinzunehmbares Argument wäre unangegriffen. Also
    stimmen die \textit{stable}-Extensionen genau mit den naiven Mengen
    überein.

    Für die \textit{preferred}-Semantik gilt allgemein
    \(\Lambda_{\mathrm{st}}(F)\subseteq\Lambda_{\mathrm{pr}}(F)\). Sei
    umgekehrt \(E\) eine \textit{preferred}-Extension, also
    \(\subseteq\)-maximal admissibel. Wäre \(E\) nicht maximal konfliktfrei, so
    gäbe es eine konfliktfreie Menge \(E\subsetneq T\); nach
    Lemma~\ref{lem:sym-cf-adm} wäre \(T\) admissibel, im Widerspruch zur
    Maximalität von \(E\). Also ist \(E\) naiv und damit \textit{stable}. Da
    eine \textit{stable}-Extension \(E\) die Labellingdarstellung
    \(\mathit{in}=E\), \(\mathit{out}=A\setminus E\), \(\mathit{und}=\emptyset\)
    besitzt und dies zugleich das zu \(E\) gehörige
    \textit{preferred}-Labelling ist, stimmen auch die Labellingmengen überein.

    Da \(A\) endlich ist, besitzt jede konfliktfreie Menge eine maximale
    Obermenge; insbesondere existiert mindestens eine naive Menge.
\end{proof}

\begin{Corollary}
    \label{cor:sym-extend}
    Sei \(F=(A,R)\) symmetrisch und irreflexiv und sei \(T\subseteq A\)
    konfliktfrei. Dann existiert ein \(\lambda\in\Lambda_{\mathrm{st}}(F)\) mit
    \(T\subseteq\mathit{in}(\lambda)\). Ist zusätzlich \(a\in A\setminus T\) und
    besitzt \(a\) einen Nachbarn in \(T\), so gilt
    \(\lambda(a)=\mathit{out}\) für jedes solche \(\lambda\).
\end{Corollary}

\begin{proof}
    Man erweitere \(T\) zu einer naiven Menge \(S\) und wende
    Proposition~\ref{prop:sym-naiv-st-pr} an. Besitzt \(a\) einen Nachbarn in
    \(T\subseteq S\), so verhindert die Konfliktfreiheit \(a\in S\).
\end{proof}

Korollar~\ref{cor:sym-extend} ist das Arbeitspferd aller folgenden Beweise: Es
erlaubt, ein Labelling durch Angabe eines konfliktfreien \enquote{Kerns} zu
spezifizieren, sofern jedes Argument, das ausgeschlossen werden soll, einen
Nachbarn im Kern besitzt.

\begin{Proposition}
    \label{prop:sym-klassisch-p}
    Sei \(F=(A,R)\) symmetrisch und irreflexiv und sei \(\sigma\in\{st,pr\}\).
    Dann gilt:
    \begin{enumerate}
        \item \(\Lambda_\sigma(F)\neq\emptyset\).
        \item Die Verifikation, ob ein gegebenes Labelling in
        \(\Lambda_\sigma(F)\) liegt, ist in Polynomialzeit möglich.
        \item Jedes \(a\in A\) ist kredulös akzeptiert.
        \item Ein Argument \(a\in A\) ist genau dann skeptisch akzeptiert, wenn
        \(N(a)=\emptyset\).
    \end{enumerate}
\end{Proposition}

\begin{proof}
    (1) und (2) folgen aus Proposition~\ref{prop:sym-naiv-st-pr}: Zu prüfen ist
    lediglich, ob \(\mathit{in}(\lambda)\) konfliktfrei ist und ob jedes
    Argument außerhalb einen Nachbarn darin besitzt; beides ist in
    \(\mathcal{O}(|A|+|R|)\) möglich.

    (3) Wegen der Irreflexivität ist \(\{a\}\) konfliktfrei; nach
    Korollar~\ref{cor:sym-extend} existiert ein \(\lambda\in\Lambda_\sigma(F)\)
    mit \(\lambda(a)=\mathit{in}\).

    (4) Ist \(N(a)=\emptyset\), so ist \(a\) in jeder maximal konfliktfreien
    Menge enthalten, denn \(a\) steht mit keinem Argument in Konflikt. Ist
    dagegen \(b\in N(a)\), so liefert Korollar~\ref{cor:sym-extend} ein
    \(\lambda\) mit \(\lambda(b)=\mathit{in}\) und folglich
    \(\lambda(a)=\mathit{out}\).
\end{proof}

Proposition~\ref{prop:sym-klassisch-p} formalisiert den Befund von
\textcite{Coste:2005:symmetric-af-irflx-coherent}, dass in symmetrischen
irreflexiven AFs nur zwei sehr einfache und effizient prüfbare Formen der
Akzeptanz auftreten. Nach diesem Bild müsste auch die bedingte Unabhängigkeit
leicht sein: Sie ist nach Definition~\ref{def:BU:AFs} eine Aussage über
\(\Lambda_\sigma(F)\), und alle punktweisen Fragen über \(\Lambda_\sigma(F)\)
sind trivial.

\section{Die eigentliche Schwierigkeit: exakte Realisierbarkeit}
\label{sec:sym-realisierbarkeit}

Die bedingte Unabhängigkeit fragt nicht nach dem Label eines einzelnen
Arguments, sondern danach, ob ein \emph{vollständig vorgeschriebenes
Teillabelling} realisiert werden kann. Genau hier bricht die Einfachheit
zusammen.

\begin{Definition}
    \label{def:sym-real}
    Das Realisierbarkeitsproblem \(\textit{Real}^{\mathrm{sym}}_{\mathrm{st}}\)
    ist wie folgt definiert:
    \[
    \begin{array}{ll}
    \text{Input:}  & \text{symmetrisches, irreflexives AF } F=(A,R),\
                     W\subseteq A,\ \eta:W\to\{\mathit{in},\mathit{out}\}\\
    \text{Output:} & \text{TRUE} \Longleftrightarrow
                     \exists\lambda\in\Lambda_{\mathrm{st}}(F):\lambda|_W=\eta.
    \end{array}
    \]
\end{Definition}

Der Unterschied zur kredulösen Akzeptanz ist der Zwang, Argumente auch
\emph{ausschließen} zu müssen: Ein Argument \(a\) mit \(\eta(a)=\mathit{out}\)
verlangt einen Nachbarn im \(\mathit{in}\)-Teil, und mehrere solche Forderungen
müssen gleichzeitig durch eine einzige konfliktfreie Menge erfüllt werden.

\begin{Proposition}
    \label{prop:sym-real-npc}
    \(\textit{Real}^{\mathrm{sym}}_{\mathrm{st}}\) ist NP-vollständig.
\end{Proposition}

\begin{proof}
    Zugehörigkeit: Man rät \(\lambda\) und verifiziert in Polynomialzeit
    (Proposition~\ref{prop:sym-klassisch-p}(2)), dass
    \(\lambda\in\Lambda_{\mathrm{st}}(F)\) und \(\lambda|_W=\eta\).

    Härte: Wir reduzieren von 3-SAT. Sei \(\varphi=\bigwedge_{k=1}^{m}C_k\) eine
    KNF-Formel über den Variablen \(y_1,\dots,y_l\). Wir definieren ein
    symmetrisches irreflexives AF \(F_\varphi=(A,R)\) durch
    \[
        A:=\{w_j,\overline{w_j}\mid 1\leq j\leq l\}\cup\{c_k\mid 1\leq k\leq m\}
    \]
    und die (symmetrisch abgeschlossene) Angriffsrelation, die genau die
    folgenden Konflikte enthält:
    \begin{itemize}
        \item \(\{w_j,\overline{w_j}\}\) für jedes \(j\),
        \item \(\{w_j,c_k\}\), falls \(y_j\) positiv in \(C_k\) vorkommt,
        \item \(\{\overline{w_j},c_k\}\), falls \(y_j\) negativ in \(C_k\)
        vorkommt.
    \end{itemize}
    Wir setzen \(W:=\{c_1,\dots,c_m\}\) und \(\eta\equiv\mathit{out}\).

    Sei \(\lambda\in\Lambda_{\mathrm{st}}(F_\varphi)\) mit \(\lambda|_W=\eta\)
    und \(S:=\mathit{in}(\lambda)\). Da \(S\) konfliktfrei ist, gilt nicht
    gleichzeitig \(w_j\in S\) und \(\overline{w_j}\in S\); wir definieren
    \(\beta(y_j):=\top\) genau dann, wenn \(w_j\in S\). Nach
    Proposition~\ref{prop:sym-naiv-st-pr} ist \(S\) naiv, also besitzt jedes
    \(c_k\notin S\) einen Nachbarn in \(S\). Die Nachbarn von \(c_k\) sind
    ausschließlich Literalargumente, also enthält \(S\) ein Literalargument von
    \(C_k\). Ist dies \(w_j\), so gilt \(\beta(y_j)=\top\) und \(y_j\in C_k\);
    ist es \(\overline{w_j}\), so ist \(w_j\notin S\), also
    \(\beta(y_j)=\bot\) und \(\neg y_j\in C_k\). In beiden Fällen ist \(C_k\)
    unter \(\beta\) erfüllt. Somit ist \(\varphi\) erfüllbar.

    Sei umgekehrt \(\beta\) eine erfüllende Belegung und
    \(T:=\{w_j\mid\beta(y_j)=\top\}\cup\{\overline{w_j}\mid\beta(y_j)=\bot\}\).
    \(T\) ist konfliktfrei, denn \(T\) enthält aus jedem Paar genau ein
    Argument und die Paare stehen untereinander nicht in Konflikt. Jedes
    \(c_k\) besitzt ein unter \(\beta\) wahres Literal, dessen Literalargument
    in \(T\) liegt. Nach Korollar~\ref{cor:sym-extend} existiert ein
    \(\lambda\in\Lambda_{\mathrm{st}}(F_\varphi)\) mit
    \(T\subseteq\mathit{in}(\lambda)\), und alle \(c_k\) tragen dort
    \(\mathit{out}\). Also gilt \(\lambda|_W=\eta\).

    Die Konstruktion ist offensichtlich polynomiell.
\end{proof}

Proposition~\ref{prop:sym-real-npc} erklärt, warum die Tractability der
klassischen Aufgaben nicht auf die Unabhängigkeit durchschlägt: Der innere
Existenzquantor aus Definition~\ref{def:BU:AFs} ist bereits auf dieser Klasse
NP-vollständig. Was noch fehlt, ist ein Mechanismus, der den äußeren
Allquantor über \(\lambda_1,\lambda_2\) in einen echten \(\forall\)-Quantor
über Belegungen verwandelt. Dieser Mechanismus ist der Gegenstand des nächsten
Abschnitts.

\section{Hauptresultat}
\label{sec:sym-hauptresultat}

Die zentrale konstruktive Schwierigkeit in symmetrischen AFs ist, dass jeder
Angriff in beide Richtungen wirkt. Zwei Konsequenzen sind dabei entscheidend.

Erstens lässt sich die für Reduktionen übliche Invariante \enquote{aus jedem
Paar \(v,\overline{v}\) ist genau eines \(\mathit{in}\)} nicht mehr durch einen
bloßen Zweierzyklus erzwingen: Sobald \(v\) oder \(\overline{v}\) weitere
Nachbarn besitzen, können in einer naiven Menge beide fehlen, weil sie von
anderer Seite dominiert werden. Wir lösen das, indem wir Wert- und Signalknoten
trennen und zu einem Viererkreis verbinden.

Zweitens lässt sich die in gerichteten Konstruktionen übliche Aggregation
\enquote{\(\varphi\) ist \(\mathit{in}\), also sind alle Klauselargumente
\(\mathit{out}\), also ist jede Klausel erfüllt} nicht nachbauen: Ein
Aggregationsargument \(\varphi\), das alle Klauselargumente angreift,
dominiert diese selbst und macht damit jede Forderung an die Literale
gegenstandslos. Wir verzichten deshalb auf ein Aggregationsargument und
schreiben die Klauselargumente stattdessen \emph{direkt} als
\(\mathit{out}\)-Teil der Unabhängigkeitsanfrage vor. Damit die Menge der so
vorschreibbaren Muster klein und beherrschbar bleibt, verbinden wir die
Klauselargumente zu einer Clique.

\begin{Definition}
    \label{def:sym-konstruktion}
    Sei \(\Psi=\forall X\exists Y\,\varphi\) mit
    \(X=\{x_1,\dots,x_n\}\), \(Y=\{y_1,\dots,y_l\}\) und
    \(\varphi=\bigwedge_{k=1}^{m}C_k\) in KNF, \(m\geq 1\). Wir setzen
    \(X^{+}:=X\cup\{x_0\}\) mit einer frischen universellen Variablen \(x_0\)
    und
    \[
        \widehat{\varphi}:=\bigwedge_{k=1}^{m}\bigl(x_0\lor C_k\bigr).
    \]
    Das AF \(G_\Psi=(A_\Psi,R_\Psi)\) besitzt die Argumentmenge
    \[
    \begin{aligned}
        A_\Psi:={}&
        \{u_i,\overline{u_i},s_i,\overline{s_i}\mid x_i\in X^{+}\}
        \;\cup\;
        \{w_j,\overline{w_j}\mid y_j\in Y\}
        \;\cup\;
        \{c_1,\dots,c_m\},
    \end{aligned}
    \]
    und \(R_\Psi\) ist die symmetrische Hülle genau der folgenden Konflikte:
    \begin{enumerate}
        \item[(K1)] \(\{u_i,\overline{u_i}\}\),
        \(\{u_i,\overline{s_i}\}\),
        \(\{\overline{u_i},s_i\}\),
        \(\{s_i,\overline{s_i}\}\)
        \quad für alle \(x_i\in X^{+}\),
        \item[(K2)] \(\{w_j,\overline{w_j}\}\) \quad für alle \(y_j\in Y\),
        \item[(K3)] \(\{s_i,c_k\}\), falls \(x_i\) positiv in \(C_k\) vorkommt,
        und \(\{\overline{s_i},c_k\}\), falls \(x_i\) negativ in \(C_k\)
        vorkommt \quad(\(1\leq i\leq n\)); zusätzlich \(\{s_0,c_k\}\) für alle
        \(k\),
        \item[(K4)] \(\{w_j,c_k\}\), falls \(y_j\) positiv in \(C_k\) vorkommt,
        und \(\{\overline{w_j},c_k\}\), falls \(y_j\) negativ in \(C_k\)
        vorkommt,
        \item[(K5)] \(\{c_k,c_{k'}\}\) für alle \(k\neq k'\).
    \end{enumerate}
    Als Unabhängigkeitsanfrage setzen wir
    \[
        \mathcal{X}:=\{c_1,\dots,c_m\},\qquad
        \mathcal{Y}:=\{u_i,\overline{u_i}\mid x_i\in X^{+}\},\qquad
        Z:=\emptyset.
    \]
\end{Definition}

Das AF \(G_\Psi\) ist nach Konstruktion symmetrisch und irreflexiv, besitzt
\(4(n+1)+2l+m\) Argumente und ist in Polynomialzeit berechenbar. Die Blöcke
(K1) bilden für jede universelle Variable einen Viererkreis
\(u_i-\overline{u_i}-s_i-\overline{s_i}-u_i\); die Wertknoten
\(u_i,\overline{u_i}\) besitzen \emph{keine} weiteren Nachbarn, während die
Signalknoten \(s_i,\overline{s_i}\) die Verbindung zu den Klauselargumenten
herstellen.

\begin{Lemma}[Wertknoten]
    \label{lem:sym-wertknoten}
    Sei \(\lambda\in\Lambda_{\mathrm{st}}(G_\Psi)\) und
    \(S:=\mathit{in}(\lambda)\). Dann gilt für jedes \(x_i\in X^{+}\):
    \(|S\cap\{u_i,\overline{u_i}\}|=1\).
\end{Lemma}

\begin{proof}
    Wegen \(\{u_i,\overline{u_i}\}\in R_\Psi\) und der Konfliktfreiheit von
    \(S\) liegen nicht beide in \(S\). Seien nun beide nicht in \(S\). Nach
    Proposition~\ref{prop:sym-naiv-st-pr} ist \(S\) naiv, also besitzt
    \(u_i\) einen Nachbarn in \(S\). Die Nachbarn von \(u_i\) sind nach (K1)
    genau \(\overline{u_i}\) und \(\overline{s_i}\); folglich gilt
    \(\overline{s_i}\in S\). Analog besitzt \(\overline{u_i}\) die Nachbarn
    \(u_i\) und \(s_i\), also \(s_i\in S\). Wegen
    \(\{s_i,\overline{s_i}\}\in R_\Psi\) widerspricht das der
    Konfliktfreiheit von \(S\).
\end{proof}

Für \(\lambda\in\Lambda_{\mathrm{st}}(G_\Psi)\) mit
\(S=\mathit{in}(\lambda)\) definieren wir daher die \emph{induzierte Belegung}
\[
    \alpha_\lambda:X^{+}\to\{\top,\bot\},\quad
    \alpha_\lambda(x_i)=\top:\iff u_i\in S,
\]
sowie
\[
    \beta_\lambda:Y\to\{\top,\bot\},\quad
    \beta_\lambda(y_j)=\top:\iff w_j\in S.
\]

\begin{Lemma}[Signaltreue]
    \label{lem:sym-signaltreue}
    Sei \(\lambda\in\Lambda_{\mathrm{st}}(G_\Psi)\), \(S=\mathit{in}(\lambda)\).
    Dann gilt:
    \begin{enumerate}
        \item \(s_i\in S\Rightarrow\alpha_\lambda(x_i)=\top\) und
        \(\overline{s_i}\in S\Rightarrow\alpha_\lambda(x_i)=\bot\).
        \item \(w_j\in S\Rightarrow\beta_\lambda(y_j)=\top\) und
        \(\overline{w_j}\in S\Rightarrow\beta_\lambda(y_j)=\bot\).
    \end{enumerate}
\end{Lemma}

\begin{proof}
    (1) Aus \(s_i\in S\) und \(\{\overline{u_i},s_i\}\in R_\Psi\) folgt
    \(\overline{u_i}\notin S\), mit Lemma~\ref{lem:sym-wertknoten} also
    \(u_i\in S\) und damit \(\alpha_\lambda(x_i)=\top\). Der zweite Fall
    verläuft symmetrisch über \(\{u_i,\overline{s_i}\}\in R_\Psi\).

    (2) \(w_j\in S\) liefert \(\beta_\lambda(y_j)=\top\) definitionsgemäß; aus
    \(\overline{w_j}\in S\) folgt \(w_j\notin S\) wegen
    \(\{w_j,\overline{w_j}\}\in R_\Psi\), also \(\beta_\lambda(y_j)=\bot\).
\end{proof}

\begin{Lemma}[Realisierbare Teillabellings]
    \label{lem:sym-teile}
    Sei \(\lambda\in\Lambda_{\mathrm{st}}(G_\Psi)\) und
    \(S=\mathit{in}(\lambda)\). Dann gilt:
    \begin{enumerate}
        \item \(|S\cap\mathcal{X}|\leq 1\).
        \item Zu jedem \(k\in\{1,\dots,m\}\) und jeder Belegung
        \(\alpha:X^{+}\to\{\top,\bot\}\) existiert
        \(\lambda'\in\Lambda_{\mathrm{st}}(G_\Psi)\) mit
        \(\mathit{in}(\lambda')\cap\mathcal{X}=\{c_k\}\) und
        \(\alpha_{\lambda'}=\alpha\).
        \item Zu jeder Belegung \(\alpha:X^{+}\to\{\top,\bot\}\) existiert
        genau dann ein \(\lambda'\in\Lambda_{\mathrm{st}}(G_\Psi)\) mit
        \(\mathit{in}(\lambda')\cap\mathcal{X}=\emptyset\) und
        \(\alpha_{\lambda'}=\alpha\), wenn eine Belegung \(\beta\) von \(Y\)
        mit \(\widehat{\varphi}(\alpha,\beta)=\top\) existiert.
        \item Zu jeder Belegung \(\alpha:X^{+}\to\{\top,\bot\}\) existiert
        \(\lambda'\in\Lambda_{\mathrm{st}}(G_\Psi)\) mit
        \(\alpha_{\lambda'}=\alpha\).
    \end{enumerate}
\end{Lemma}

\begin{proof}
    Wir schreiben \(U_\alpha:=\{u_i\mid\alpha(x_i)=\top\}
    \cup\{\overline{u_i}\mid\alpha(x_i)=\bot\}\) und
    \(S_\alpha:=\{s_i\mid\alpha(x_i)=\top\}
    \cup\{\overline{s_i}\mid\alpha(x_i)=\bot\}\).
    Nach (K1) ist \(U_\alpha\cup S_\alpha\) konfliktfrei: Innerhalb eines
    Viererkreises werden nur gegenüberliegende Knoten gewählt
    (\(u_i\) mit \(s_i\) beziehungsweise \(\overline{u_i}\) mit
    \(\overline{s_i}\)), und zwischen verschiedenen Blöcken bestehen keine
    Konflikte. Ferner besitzt jedes Argument aus
    \(\{u_i,\overline{u_i}\}\setminus U_\alpha\) seinen Partner in
    \(U_\alpha\) als Nachbarn. Nach Korollar~\ref{cor:sym-extend} gilt
    daher für jede Erweiterung eines \(U_\alpha\) enthaltenden konfliktfreien
    Kerns zu einem \textit{stable}-Labelling \(\lambda'\) bereits
    \(\alpha_{\lambda'}=\alpha\).

    (1) Die Klauselargumente bilden nach (K5) eine Clique; eine konfliktfreie
    Menge enthält daher höchstens eines von ihnen.

    (2) Der Kern \(T:=U_\alpha\cup\{c_k\}\) ist konfliktfrei, da Wertknoten
    nach (K1) nur Wert- und Signalknoten desselben Blocks als Nachbarn haben
    und insbesondere zu keinem Klauselargument benachbart sind. Sei
    \(\lambda'\) eine Erweiterung gemäß Korollar~\ref{cor:sym-extend}. Dann gilt
    \(\alpha_{\lambda'}=\alpha\) nach der Vorbemerkung, und jedes \(c_{k'}\)
    mit \(k'\neq k\) besitzt den Nachbarn \(c_k\in T\), ist also
    \(\mathit{out}\). Somit ist
    \(\mathit{in}(\lambda')\cap\mathcal{X}=\{c_k\}\).

    (3) \enquote{\(\Leftarrow\)}: Sei \(\beta\) mit
    \(\widehat{\varphi}(\alpha,\beta)=\top\) und sei
    \[
        T:=U_\alpha\cup S_\alpha\cup
        \{w_j\mid\beta(y_j)=\top\}\cup\{\overline{w_j}\mid\beta(y_j)=\bot\}.
    \]
    \(T\) ist konfliktfrei: Konflikte bestehen nach (K3), (K4) nur zwischen
    Signal- beziehungsweise Literalargumenten und Klauselargumenten, und
    \(T\) enthält keine Klauselargumente; innerhalb der Blöcke wurde die
    Konfliktfreiheit bereits festgestellt. Da \(\widehat{\varphi}\) unter
    \((\alpha,\beta)\) erfüllt ist, besitzt jede Klausel \(x_0\lor C_k\) ein
    wahres Literal, dessen zugehöriges Signal- oder Literalargument nach (K3),
    (K4) in \(T\) liegt und zu \(c_k\) benachbart ist. Nach
    Korollar~\ref{cor:sym-extend} existiert ein \textit{stable}-Labelling
    \(\lambda'\) mit \(T\subseteq\mathit{in}(\lambda')\); dort sind alle
    \(c_k\) \(\mathit{out}\) und es gilt \(\alpha_{\lambda'}=\alpha\).

    \enquote{\(\Rightarrow\)}: Sei \(\lambda'\) mit
    \(\mathit{in}(\lambda')\cap\mathcal{X}=\emptyset\) und
    \(\alpha_{\lambda'}=\alpha\), und sei \(S'=\mathit{in}(\lambda')\). Da
    \(S'\) naiv ist, besitzt jedes \(c_k\) einen Nachbarn in \(S'\). Die
    Nachbarn von \(c_k\) sind nach (K3)--(K5) Signalargumente,
    Literalargumente und andere Klauselargumente; letztere liegen nicht in
    \(S'\). Also enthält \(S'\) ein Signal- oder Literalargument zu einem
    Literal von \(x_0\lor C_k\). Nach Lemma~\ref{lem:sym-signaltreue} ist
    dieses Literal unter \((\alpha,\beta_{\lambda'})\) wahr. Da \(k\) beliebig
    war, gilt \(\widehat{\varphi}(\alpha,\beta_{\lambda'})=\top\).

    (4) Man erweitere \(U_\alpha\) nach Korollar~\ref{cor:sym-extend}.
\end{proof}

\begin{Theorem}
    \label{th:sym-st-pi2p}
    \(\textit{Ind}^{\mathrm{sym}}_{\mathrm{st}}\) ist
    \(\Pi_2^\text{P}\)-vollständig. Die Härte gilt bereits für
    \(Z=\emptyset\).
\end{Theorem}

\begin{proof}
    \emph{Zugehörigkeit.} Symmetrische irreflexive AFs bilden eine Teilklasse
    aller AFs; die obere Schranke folgt daher aus
    Theorem~\ref{th:ind-st-is-pi2p-c}.

    \emph{Härte.} Sei \(\Psi=\forall X\exists Y\,\varphi\) eine Instanz des
    \(\Pi_2^\text{P}\)-vollständigen Problems \(\mathrm{QBF}^2_\forall\) und sei
    \(G_\Psi\) mit \(\mathcal{X},\mathcal{Y},Z=\emptyset\) wie in
    Definition~\ref{def:sym-konstruktion}. Die Mengen \(\mathcal{X}\),
    \(\mathcal{Y}\) und \(Z\) sind paarweise disjunkt, und \(G_\Psi\) ist
    symmetrisch und irreflexiv. Wegen \(Z=\emptyset\) ist die Voraussetzung
    \(\lambda_1|_Z=\lambda_2|_Z\) aus Definition~\ref{def:BU:AFs} für jedes
    Paar erfüllt. Wir zeigen
    \[
        \Psi\text{ ist gültig}
        \quad\Longleftrightarrow\quad
        \mathcal{X}\bot_{\mathrm{st}}\mathcal{Y}\mid\emptyset
        \text{ in }G_\Psi.
    \]

    Vorab halten wir fest, dass wegen
    \(\widehat{\varphi}=\bigwedge_k(x_0\lor C_k)\) gilt:
    \(\Psi\) ist genau dann gültig, wenn
    \(\forall X^{+}\exists Y\,\widehat{\varphi}\) gültig ist. Für
    \(\alpha(x_0)=\top\) ist \(\widehat{\varphi}\) unter jeder Belegung von
    \(Y\) erfüllt, für \(\alpha(x_0)=\bot\) gilt
    \(\widehat{\varphi}\equiv\varphi\).

    \enquote{\(\Rightarrow\)}: Sei \(\Psi\) gültig und seien
    \(\lambda_1,\lambda_2\in\Lambda_{\mathrm{st}}(G_\Psi)\). Sei
    \(\alpha:=\alpha_{\lambda_2}\); nach Lemma~\ref{lem:sym-wertknoten} ist
    \(\lambda_2|_{\mathcal{Y}}\) durch \(\alpha\) eindeutig bestimmt und
    umgekehrt. Nach Lemma~\ref{lem:sym-teile}(1) ist
    \(\mathit{in}(\lambda_1)\cap\mathcal{X}\) entweder leer oder von der Form
    \(\{c_k\}\).

    Im Fall \(\mathit{in}(\lambda_1)\cap\mathcal{X}=\{c_k\}\) liefert
    Lemma~\ref{lem:sym-teile}(2) ein \(\lambda_3\) mit
    \(\mathit{in}(\lambda_3)\cap\mathcal{X}=\{c_k\}\) und
    \(\alpha_{\lambda_3}=\alpha\), also
    \(\lambda_3|_{\mathcal{X}}=\lambda_1|_{\mathcal{X}}\) und
    \(\lambda_3|_{\mathcal{Y}}=\lambda_2|_{\mathcal{Y}}\).

    Im Fall \(\mathit{in}(\lambda_1)\cap\mathcal{X}=\emptyset\) existiert wegen
    der Gültigkeit von \(\forall X^{+}\exists Y\,\widehat{\varphi}\) eine
    Belegung \(\beta\) mit \(\widehat{\varphi}(\alpha,\beta)=\top\);
    Lemma~\ref{lem:sym-teile}(3) liefert das gesuchte \(\lambda_3\).

    Damit gilt \(\mathcal{X}\bot_{\mathrm{st}}\mathcal{Y}\mid\emptyset\).

    \enquote{\(\Leftarrow\)}: Sei \(\Psi\) nicht gültig. Dann existiert eine
    Belegung \(\alpha\) von \(X\), sodass \(\varphi(\alpha,\beta)=\bot\) für
    alle \(\beta\) gilt. Wir setzen
    \(\alpha^{+}:=\alpha\cup\{x_0\mapsto\bot\}\); dann gilt
    \(\widehat{\varphi}(\alpha^{+},\beta)=\bot\) für alle \(\beta\).

    Wähle \(\lambda_2\) nach Lemma~\ref{lem:sym-teile}(4) mit
    \(\alpha_{\lambda_2}=\alpha^{+}\). Wähle \(\lambda_1\) nach
    Lemma~\ref{lem:sym-teile}(3), angewandt auf eine beliebige Belegung
    \(\alpha'\) mit \(\alpha'(x_0)=\top\); wegen
    \(\widehat{\varphi}(\alpha',\beta)=\top\) für jedes \(\beta\) existiert ein
    solches \(\lambda_1\) mit
    \(\mathit{in}(\lambda_1)\cap\mathcal{X}=\emptyset\).

    Angenommen, es existierte \(\lambda_3\in\Lambda_{\mathrm{st}}(G_\Psi)\) mit
    \(\lambda_3|_{\mathcal{X}}=\lambda_1|_{\mathcal{X}}\) und
    \(\lambda_3|_{\mathcal{Y}}=\lambda_2|_{\mathcal{Y}}\). Dann gilt
    \(\mathit{in}(\lambda_3)\cap\mathcal{X}=\emptyset\) und
    \(\alpha_{\lambda_3}=\alpha^{+}\). Nach Lemma~\ref{lem:sym-teile}(3)
    existierte dann ein \(\beta\) mit
    \(\widehat{\varphi}(\alpha^{+},\beta)=\top\), im Widerspruch zur Wahl von
    \(\alpha\). Also gilt
    \(\mathcal{X}\not\bot_{\mathrm{st}}\mathcal{Y}\mid\emptyset\).

    Die Konstruktion ist in Polynomialzeit berechenbar, womit die Härte gezeigt
    ist.
\end{proof}

\begin{Corollary}
    \label{cor:sym-nicht-ocf}
    Theorem~\ref{th:sym-st-pi2p} folgt nicht aus
    Proposition~\ref{prop:ocf-af-ci-st-pi^2}.
\end{Corollary}

\begin{proof}
    Für \(m\geq 3\) enthält \(G_\Psi\) wegen (K5) den Dreierkreis
    \(c_1\to c_2\to c_3\to c_1\) und ist damit nicht ungeradezykelfrei. Ebenso
    ist \(G_\Psi\) wegen dieses Dreiecks nicht bipartit. Die Klasse der
    symmetrischen irreflexiven AFs ist somit weder in der Klasse der
    ungeradezykelfreien noch in der der bipartiten AFs enthalten.
\end{proof}

Bemerkenswert ist, dass symmetrische irreflexive AFs dennoch kohärent sind: Die
Kohärenz wird hier nicht über die Abwesenheit ungerader Zyklen erreicht, sondern
über die Symmetrie selbst (Lemma~\ref{lem:sym-cf-adm}). Die Zyklenfreiheit ist
also hinreichend, aber nicht notwendig für Kohärenz, und für die Komplexität
der Unabhängigkeit ist offenkundig nicht die Kohärenz entscheidend.

\section{Konsequenz für die \textit{preferred}-Semantik}
\label{sec:sym-preferred}

\begin{Lemma}[Transferlemma]
    \label{lem:sym-transfer}
    Sei \(F=(A,R)\) ein AF und seien \(\sigma,\tau\) Semantiken mit
    \(\Lambda_\sigma(F)=\Lambda_\tau(F)\). Dann gilt für alle paarweise
    disjunkten \(X,Y,Z\subseteq A\):
    \(X\bot_\sigma Y\mid Z\) genau dann, wenn \(X\bot_\tau Y\mid Z\).
\end{Lemma}

\begin{proof}
    Die definierende Bedingung aus Definition~\ref{def:BU:AFs} quantifiziert
    ausschließlich über Elemente von \(\Lambda_\sigma(F)\). Ersetzt man diese
    Menge durch die nach Voraussetzung identische Menge \(\Lambda_\tau(F)\), so
    geht die Bedingung Wort für Wort in die entsprechende Bedingung für
    \(\tau\) über.
\end{proof}

\begin{Corollary}
    \label{cor:sym-pr-pi2p}
    \(\textit{Ind}^{\mathrm{sym}}_{\mathrm{pr}}\) ist
    \(\Pi_2^\text{P}\)-vollständig.
\end{Corollary}

\begin{proof}
    Nach Proposition~\ref{prop:sym-naiv-st-pr} gilt
    \(\Lambda_{\mathrm{pr}}(F)=\Lambda_{\mathrm{st}}(F)\) für jedes
    symmetrische irreflexive AF \(F\). Mit Lemma~\ref{lem:sym-transfer} sind
    \(\textit{Ind}^{\mathrm{sym}}_{\mathrm{pr}}\) und
    \(\textit{Ind}^{\mathrm{sym}}_{\mathrm{st}}\) dieselbe Sprache; die
    Behauptung folgt aus Theorem~\ref{th:sym-st-pi2p}.
\end{proof}

Für \(\sigma=\mathrm{pr}\) ist die Klasse also \emph{echt} einfacher: Das
allgemeine Problem ist nach Theorem~\ref{th:ind-pr-is-pi3p-c}
\(\Pi_3^\text{P}\)-vollständig, auf symmetrischen irreflexiven AFs sinkt es um
eine Stufe der polynomiellen Hierarchie. Der Grund ist strukturell: Die
Maximalitätsbedingung der \textit{preferred}-Semantik, die im allgemeinen Fall
den zusätzlichen Quantorenwechsel erzwingt, kollabiert hier zu einer in
Polynomialzeit prüfbaren Bedingung (Proposition~\ref{prop:sym-klassisch-p}(2)).
Unter der üblichen Annahme \(\Pi_2^\text{P}\neq\Pi_3^\text{P}\) ist dieser
Gewinn nicht nur formal.

Für \(\sigma=\mathrm{st}\) tritt dagegen kein Gewinn ein. Die Erklärung liefert
Proposition~\ref{prop:sym-real-npc}: Die Symmetrie vereinfacht alle
\emph{punktweisen} Fragen über \(\Lambda_{\mathrm{st}}(F)\), lässt aber die
\emph{simultane} Realisierbarkeit vorgeschriebener Teillabellings NP-hart. Da
die bedingte Unabhängigkeit genau von dieser simultanen Realisierbarkeit
handelt, überträgt sich die Vereinfachung nicht.

\section{Ein in Polynomialzeit entscheidbares Fragment}
\label{sec:sym-fragment}

Die Härte in Theorem~\ref{th:sym-st-pi2p} benötigt eine unbeschränkt große
Menge \(\mathcal{Y}\). Beschränkt man beide Mengen auf einzelne Argumente und
setzt \(Z=\emptyset\), so wird das Problem tatsächlich einfach. Damit ist die
Klasse in einem präzisen Sinn doch einfacher als der allgemeine Fall, in dem
schon die zugrundeliegenden Realisierbarkeitsfragen NP-vollständig sind.

\begin{Proposition}
    \label{prop:sym-singleton-p}
    Sei \(F=(A,R)\) symmetrisch und irreflexiv und seien \(x,y\in A\) mit
    \(x\neq y\). Dann ist \(\{x\}\bot_{\mathrm{st}}\{y\}\mid\emptyset\) in Zeit
    \(\mathcal{O}(|A|^2+|A|\cdot|R|)\) entscheidbar.
\end{Proposition}

\begin{proof}
    Wegen \(Z=\emptyset\) ist die Bedingung \(\lambda_1|_Z=\lambda_2|_Z\)
    stets erfüllt. Setze
    \[
        \mathcal{R}:=\bigl\{(\lambda(x),\lambda(y))\;\big|\;
        \lambda\in\Lambda_{\mathrm{st}}(F)\bigr\}
        \subseteq\{\mathit{in},\mathit{out}\}^2 .
    \]
    Nach Definition~\ref{def:BU:AFs} gilt
    \(\{x\}\bot_{\mathrm{st}}\{y\}\mid\emptyset\) genau dann, wenn aus
    \((\alpha_1,\beta_1),(\alpha_2,\beta_2)\in\mathcal{R}\) stets
    \((\alpha_1,\beta_2)\in\mathcal{R}\) folgt. Da \(\mathcal{R}\) höchstens
    vier Elemente besitzt, genügt es, die vier Zugehörigkeiten zu bestimmen.
    Wir benutzen durchgehend Proposition~\ref{prop:sym-naiv-st-pr} und
    Korollar~\ref{cor:sym-extend}.

    \emph{\((\mathit{in},\mathit{in})\).} Genau dann realisierbar, wenn
    \((x,y)\notin R\): Ist \(\{x,y\}\) konfliktfrei, erweitere man es; andernfalls
    verbietet die Konfliktfreiheit beide gemeinsam.

    \emph{\((\mathit{in},\mathit{out})\).} Ist \((x,y)\in R\), so erweitere man
    \(\{x\}\); dann ist \(y\) \(\mathit{out}\). Ist \((x,y)\notin R\), so ist
    das Paar genau dann realisierbar, wenn \(N(y)\setminus N[x]\neq\emptyset\).
    Existiert nämlich ein solches \(w\), so ist \(\{x,w\}\) konfliktfrei und
    jede Erweiterung enthält \(x\), aber wegen \(w\in N(y)\) nicht \(y\).
    Existiert umgekehrt \(\lambda\) mit \(\lambda(x)=\mathit{in}\) und
    \(\lambda(y)=\mathit{out}\), so besitzt \(y\) einen Nachbarn
    \(w\in\mathit{in}(\lambda)\); wegen \(x\notin N(y)\) ist \(w\neq x\), und
    wegen der Konfliktfreiheit gilt \(w\notin N(x)\).

    \emph{\((\mathit{out},\mathit{in})\).} Symmetrisch mit vertauschten Rollen.

    \emph{\((\mathit{out},\mathit{out})\).} Genau dann realisierbar, wenn
    \(w_1\in N(x)\setminus\{y\}\) und \(w_2\in N(y)\setminus\{x\}\) mit
    \(w_1=w_2\) oder \((w_1,w_2)\notin R\) existieren. Sind solche \(w_1,w_2\)
    gegeben, so ist \(\{w_1,w_2\}\) konfliktfrei, und jede Erweiterung schließt
    \(x\) und \(y\) aus. Existiert umgekehrt \(\lambda\) mit
    \(\lambda(x)=\lambda(y)=\mathit{out}\), so besitzen \(x\) und \(y\) je einen
    Nachbarn \(w_1,w_2\in\mathit{in}(\lambda)\); diese liegen wegen
    \(\lambda(x)=\lambda(y)=\mathit{out}\) nicht in \(\{x,y\}\) und stehen
    untereinander nicht in Konflikt.

    Alle vier Tests sind durch Absuchen von \(N(x)\), \(N(y)\) beziehungsweise
    aller Paare \((w_1,w_2)\in N(x)\times N(y)\) in der angegebenen Zeit
    durchführbar. Anschließend prüft man die Rechteckbedingung auf der höchstens
    vierelementigen Menge \(\mathcal{R}\) in konstanter Zeit.
\end{proof}

Im allgemeinen Fall sind bereits die vier Realisierbarkeitstests
NP-vollständig, da \((\mathit{in},\cdot)\) die kredulöse Akzeptanz unter
\textit{stable}-Semantik enthält. Proposition~\ref{prop:sym-singleton-p} ist
also eine echte Vereinfachung; sie zeigt zugleich, dass die Härte aus
Theorem~\ref{th:sym-st-pi2p} nicht aus einzelnen Argumentpaaren stammt, sondern
aus der \emph{gemeinsamen} Vorschrift vieler Labels.

Als Konsistenzprüfung sei angemerkt, dass
Proposition~\ref{prop:cdg-afs-not-st-ci} ein Spezialfall dieser Analyse ist: Ein
vollständig gerichtetes AF ist symmetrisch und irreflexiv, seine naiven Mengen
sind genau die einelementigen Teilmengen von \(A\), und für \(x\neq y\) ist
\((\mathit{in},\mathit{in})\) wegen \((x,y)\in R\) nicht realisierbar, während
\((\mathit{in},\mathit{out})\) und \((\mathit{out},\mathit{in})\) realisierbar
sind. Die Rechteckbedingung ist damit verletzt.

\section{Die \textit{complete}-Semantik bleibt offen}
\label{sec:sym-complete}

Für \(\sigma=\mathrm{co}\) lässt sich die Konstruktion aus
Definition~\ref{def:sym-konstruktion} nicht übernehmen. Wir halten zunächst die
Struktur der \textit{complete}-Labellings fest.

\begin{Proposition}
    \label{prop:sym-co-struktur}
    Sei \(F=(A,R)\) symmetrisch und irreflexiv und sei \(\lambda\) ein
    Labelling mit \(S:=\mathit{in}(\lambda)\). Dann ist \(\lambda\) genau dann
    \textit{complete}, wenn \(S\) konfliktfrei ist,
    \(\mathit{out}(\lambda)=N(S)\) gilt und kein \(a\in A\setminus S\) mit
    \(N(a)\subseteq N(S)\) existiert.
\end{Proposition}

\begin{proof}
    Nach Definition~\ref{def:complete-labelling} ist \(a\) genau dann
    \(\mathit{out}\), wenn \(a\) einen \(\mathit{in}\)-Angreifer besitzt, also
    genau dann, wenn \(a\in N(S)\); das ist gleichbedeutend mit
    \(\mathit{out}(\lambda)=N(S)\), und \(S\cap N(S)=\emptyset\) ist genau die
    Konfliktfreiheit von \(S\). Ferner ist \(a\) genau dann \(\mathit{in}\),
    wenn alle Angreifer \(\mathit{out}\) sind, also genau dann, wenn
    \(N(a)\subseteq N(S)\). Für \(a\in S\) ist diese Bedingung wegen der
    Symmetrie automatisch erfüllt; die verbleibende Forderung ist, dass kein
    \(a\) außerhalb von \(S\) sie erfüllt.
\end{proof}

Insbesondere ist das \textit{grounded}-Labelling durch
\(\mathit{in}(\lambda_{\mathrm{gr}})=\{a\in A\mid N(a)=\emptyset\}\) gegeben:
Isolierte Argumente sind unangegriffen, und ein nicht isoliertes Argument \(a\)
besitzt einen Nachbarn, der von einem isolierten Argument nicht angegriffen
werden kann.

Damit scheitert die Reduktion aus Abschnitt~\ref{sec:sym-hauptresultat} an
Lemma~\ref{lem:sym-wertknoten}: Wählt man \(S=\emptyset\), so ist
\(N(S)=\emptyset\), und nach Proposition~\ref{prop:sym-co-struktur} ist das
Labelling, das alle nicht isolierten Argumente mit \(\mathit{und}\) versieht,
\textit{complete}. In diesem Labelling tragen sämtliche Wertknoten
\(u_i,\overline{u_i}\) das Label \(\mathit{und}\); die Zuordnung
\enquote{\(\lambda_2\) kodiert eine totale Belegung} bricht zusammen, und der
äußere Allquantor läuft über wesentlich mehr Labellings als über die
Belegungen von \(X^{+}\). Da zusätzliche \(\lambda_2\) zusätzliche
Kombinationsforderungen erzeugen, kann die Richtung
\enquote{\(\Psi\) gültig \(\Rightarrow\) unabhängig} verletzt werden.

In gerichteten Konstruktionen wird dieses Problem durch Detektorargumente
gelöst, die partielle Labellings erkennen und die eigentliche Formelauswertung
neutralisieren --- so etwa die Argumente \(d_i\), \(s\) und \(g\) im Beweis von
Proposition~\ref{prop:ocf-af-ci-co-pi^2-c}. Ein solcher Detektor benötigt einen
gerichteten Informationsfluss \(\{x_i,\overline{x_i}\}\to d_i\to s\to g\), der
sich in symmetrischen AFs nicht nachbilden lässt: Jeder Detektor griffe seine
Quelle zurück an. Wir halten daher fest:

\begin{Proposition}
    \label{prop:sym-co-obere-schranke}
    \(\textit{Ind}^{\mathrm{sym}}_{\mathrm{co}}\in\Pi_2^\text{P}\).
\end{Proposition}

\begin{proof}
    Folgt aus Theorem~\ref{th:ind-co-is-pi2p-c}, da symmetrische irreflexive
    AFs eine Teilklasse aller AFs bilden.
\end{proof}

Ob \(\textit{Ind}^{\mathrm{sym}}_{\mathrm{co}}\) auch \(\Pi_2^\text{P}\)-hart
ist, bleibt offen. Es ist die einzige der drei betrachteten Semantiken, für die
in dieser Klasse keine vollständige Einordnung erreicht wurde.

\section{Zusammenfassung}
\label{sec:sym-zusammenfassung}

\begin{table}[htbp]
    \myfloatalign
    \centering
    \renewcommand{\arraystretch}{1.15}
    \setlength{\tabcolsep}{1.0em}
    \begin{tabular}{lccc}
        \toprule
        & \(co\) & \(st\) & \(pr\) \\
        \midrule
        allgemeine AFs
            & \(\Pi_2^\text{P}\)-c
            & \(\Pi_2^\text{P}\)-c
            & \(\Pi_3^\text{P}\)-c \\
        symmetrisch, irreflexiv
            & in \(\Pi_2^\text{P}\)
            & \(\Pi_2^\text{P}\)-c
            & \(\Pi_2^\text{P}\)-c \\
        \quad davon \(|X|=|Y|=1\), \(Z=\emptyset\)
            & ---
            & in P
            & in P \\
        \bottomrule
    \end{tabular}
    \caption[Komplexität der Unabhängigkeit auf symmetrischen irreflexiven
    AFs]{Einordnung von \(\sigma\)-Unabhängigkeit auf symmetrischen
    irreflexiven AFs im Vergleich zum allgemeinen Fall. Die Einträge der
    zweiten Zeile stammen aus Proposition~\ref{prop:sym-co-obere-schranke},
    Theorem~\ref{th:sym-st-pi2p} und Korollar~\ref{cor:sym-pr-pi2p}, der
    Eintrag der dritten Zeile aus Proposition~\ref{prop:sym-singleton-p} in
    Verbindung mit Korollar~\ref{cor:sym-pr-pi2p}.}
    \label{tab:sym-komplexitaet}
\end{table}

Die Antwort auf die Ausgangsfrage ist differenziert. Symmetrische irreflexive
AFs sind für die \textit{preferred}-Semantik echt einfacher, weil dort die
Maximalitätsbedingung kollabiert und ein Quantorenwechsel entfällt. Für die
\textit{stable}-Semantik sind sie \emph{nicht} einfacher: Trotz kohärenter
Struktur, garantierter Existenz von Labellings und durchgängig polynomieller
klassischer Entscheidungsprobleme bleibt die bedingte Unabhängigkeit
\(\Pi_2^\text{P}\)-vollständig. Eine Vereinfachung tritt erst ein, wenn man
zusätzlich die Größe der Anfragemengen beschränkt.

Methodisch ist das der interessanteste Befund dieses Kapitels: Die
Tractability der klassischen Aufgaben ist kein verlässlicher Indikator für die
Komplexität der bedingten Unabhängigkeit. Maßgeblich ist stattdessen, ob
vorgeschriebene Teillabellings \emph{simultan} realisierbar sind
(Proposition~\ref{prop:sym-real-npc}) und ob der äußere Allquantor über
hinreichend viele, strukturell unterscheidbare Labellings läuft
(Lemma~\ref{lem:sym-wertknoten}). Beide Bedingungen sind in symmetrischen
irreflexiven AFs erfüllt, obwohl jede einzelne Akzeptanzfrage trivial ist.
